\documentclass[12pt,a4paper]{report}
\usepackage[utf8]{vietnam}
\usepackage[a4paper,margin=2.5cm]{geometry}
\usepackage{amsmath,amssymb}
\usepackage{graphicx}
\usepackage{booktabs,longtable,array}
\usepackage{tikz}
\usetikzlibrary{arrows.meta,positioning,calc}
\usepackage{hyperref}
\hypersetup{colorlinks=true,linkcolor=blue,citecolor=blue,urlcolor=blue}
\setlength{\parindent}{1.2em}
\setlength{\parskip}{0.3em}


\hypersetup{
    pdftitle={Nghiên cứu cơ chế xác thực và quản lý chứng thư số X.509 trong PKI},
    pdfauthor={Đoàn Văn Tuyền (23020151), Vũ Minh Quân (23020142), Nguyễn Chi Tú (23020148)},
    pdfsubject={Mật mã và An toàn thông tin}
}

\begin{document}
\pagenumbering{alph}
\begin{titlepage}
\thispagestyle{empty}
\setlength{\fboxrule}{0.8pt}
\setlength{\fboxsep}{0pt}
\noindent\makebox[\linewidth][c]{\fbox{%
\begin{minipage}[t][\dimexpr\textheight-0.3cm\relax][t]{\dimexpr\textwidth-0.6cm-2\fboxrule\relax}
\centering
\vspace*{0.65cm}
{\bfseries\fontsize{14}{17}\selectfont ĐẠI HỌC QUỐC GIA HÀ NỘI\par}
\vspace{0.15cm}
{\bfseries\fontsize{17}{21}\selectfont TRƯỜNG ĐẠI HỌC CÔNG NGHỆ\par}
\vspace{0.15cm}
{\bfseries\fontsize{12}{15}\selectfont KHOA CÔNG NGHỆ THÔNG TIN\par}

\vspace{1.25cm}
\includegraphics[width=3.4cm]{report/figures/uet.png}\par

\vspace{1.35cm}
{\bfseries\fontsize{16}{20}\selectfont BÁO CÁO BÀI TẬP MÔN HỌC\par}

\vspace{0.95cm}
{\bfseries\fontsize{16}{21}\selectfont
HẠ TẦNG KHÓA CÔNG KHAI (PKI):\\
KIẾN TRÚC, CƠ CHẾ XÁC THỰC CHỨNG THƯ X.509\\
VÀ ỨNG DỤNG\par}
\vspace{0.25cm}
{\normalsize Phân tích chuỗi tin cậy, thu hồi chứng thư và đánh giá thực nghiệm trên CA phân cấp\par}

\vspace{1.9cm}
\begin{minipage}{15.4cm}
\renewcommand{\arraystretch}{1.5}
\begin{tabular}{@{}p{4.1cm}p{6.3cm}p{4.0cm}@{}}
Mã lớp học phần: & 2627I\_INT3230\_2 & \\
Học phần: & Mật mã và An toàn thông tin & \\
Sinh viên thực hiện: & Đoàn Văn Tuyền & MSSV: 23020151 \\
& Vũ Minh Quân & MSSV: 23020142 \\
& Nguyễn Chi Tú & MSSV: 23020148 \\[0.65cm]
Giảng viên môn học: & Lê Thị Hợi & \\
\end{tabular}
\end{minipage}

\vfill
{\bfseries HÀ NỘI -- 2026\par}
\vspace*{0.8cm}
\end{minipage}%
}}
\end{titlepage}
\newgeometry{a4paper,margin=2.5cm}

\pagenumbering{roman}

\chapter*{Tóm tắt}
\addcontentsline{toc}{chapter}{Tóm tắt}
Báo cáo tập trung nghiên cứu cơ chế phân phối, quản lý và xác thực khóa công khai trong Hạ tầng khóa công khai (PKI) chuẩn ITU-T X.509, nhằm giải quyết nguy cơ tấn công thay thế khóa và người đứng giữa (MITM) trong truyền thông mạng. Bằng phương pháp kết hợp phân tích tiêu chuẩn quốc tế (RFC~5280, RFC~6960, RFC~8446, RFC~3647) với thực nghiệm thực tế, đề tài thiết lập một hệ thống PKI phân cấp hai tầng độc lập bằng OpenSSL 3.5.9 và xây dựng bộ kiểm thử tự động gồm 10 ca kiểm thử độc lập (TC01--TC10) cùng các lần chạy đối chứng. Kết quả quan sát thực nghiệm ủng hộ mạnh mẽ giả thuyết $H_1$ và $H_2$: chữ ký số hợp lệ chỉ là điều kiện cần chứ chưa đủ để chứng thư được chấp nhận; quyết định tin cậy của bên dựa vào (relying party) phụ thuộc sâu sắc vào các tham số ngữ cảnh (chuỗi tin cậy về Trust Anchor, thời hạn hiệu lực, đối chiếu SAN, mục đích EKU) và chính sách kiểm tra thu hồi CRL. Phép đo định lượng cục bộ ủng hộ giả thuyết $H_3$, chỉ ra sự gia tăng kích thước tệp CRL ($\approx 47$~bytes/mục) và thời gian thẩm định tăng hơn 2{,}3 lần ở quy mô 10.000 bản ghi thu hồi. Trên cơ sở đó, báo cáo phân tích toàn diện mô hình phòng thủ 6 lớp của PKI hiện đại, làm rõ vai trò của Certificate Transparency (RFC~9162), giao thức tự động hóa ACME (RFC~8555), và vạch ra lộ trình di trú sang các chuẩn mật mã hậu lượng tử của NIST (ML-KEM, ML-DSA, SLH-DSA).

\chapter*{Abstract}
\addcontentsline{toc}{chapter}{Abstract}
This research investigates the mechanisms of public key distribution, management, and certificate validation within the ITU-T X.509 Public Key Infrastructure (PKI) to mitigate key substitution and Man-in-the-Middle (MITM) attacks in open network communications. Combining rigorous analysis of international standards (RFC~5280, RFC~6960, RFC~8446, RFC~3647) with empirical experimentation, we establish an isolated two-tier hierarchical PKI testbed using OpenSSL 3.5.9 and execute an automated test suite comprising ten distinct test cases (TC01--TC10) alongside controlled counter-experiments. The empirical findings strongly corroborate hypotheses $H_1$ and $H_2$: a valid cryptographic signature is merely a necessary but not sufficient condition for certificate acceptance; the relying party's verification decision strictly depends on dynamic contextual constraints (trust path construction to a configured Trust Anchor, validity period, SAN matching, EKU purpose) and client revocation checking policies. Local quantitative benchmarks support hypothesis $H_3$, revealing a linear growth in CRL file size ($\approx 47$~bytes/entry) and a 2.3-fold increase in path validation latency at 10,000 revoked records. Furthermore, the report presents a comprehensive six-layer defense-in-depth model, evaluates Certificate Transparency (RFC~9162) and ACME automation (RFC~8555), and outlines technical migration pathways toward NIST Post-Quantum Cryptography standards (ML-KEM, ML-DSA, and SLH-DSA).

\tableofcontents
\listoftables
\listoffigures
\clearpage
\pagenumbering{arabic}

\chapter{Giới thiệu và xác định vấn đề nghiên cứu}
\label{chap:gioi-thieu}

Chương này trình bày bối cảnh ra đời của mật mã khóa công khai, phân tích lỗ hổng cố hữu trong phân phối khóa trần qua kịch bản tấn công người đứng giữa (MITM), từ đó dẫn dắt đến sự cần thiết của Hạ tầng khóa công khai (PKI). Tiếp đó, chương xác định rõ mục tiêu, 4 câu hỏi nghiên cứu (RQ1--RQ4), phạm vi triển khai và phương pháp luận của đề tài.

\section{Bối cảnh và bài toán}
\label{sec:context-problem}

\subsection{Bối cảnh mật mã khóa công khai}
Trong các hệ thống thông tin hiện đại, việc truyền thông qua các mạng công cộng như Internet luôn đối mặt với nguy cơ bị nghe lén, can thiệp và sửa đổi dữ liệu. Mật mã khóa công khai (Public-Key Cryptography hay mật mã bất đối xứng) ra đời như một bước đột phá quan trọng nhằm giải quyết hạn chế cốt tử của mật mã đối xứng: bài toán phân phối trước một khóa bí mật chung qua kênh truyền an toàn. Trong mật mã khóa công khai, mỗi thực thể tham gia sở hữu một cặp khóa có quan hệ toán học chặt chẽ nhưng không thể suy dẫn khóa này từ khóa kia với năng lực tính toán thực tế:
\begin{itemize}
    \item \textbf{Khóa công khai (Public Key):} Ký hiệu là $PU$ (theo quy ước kinh điển) hoặc $PK$ (theo chuẩn mật mã học hiện đại). Khóa này được công bố rộng rãi cho mọi đối tượng trong hệ thống mà không cần giữ bí mật.
    \item \textbf{Khóa bí mật (Private Key / Secret Key):} Ký hiệu tương ứng là $PR$ hoặc $SK$. Khóa này được chủ sở hữu lưu giữ tuyệt đối bí mật và không bao giờ tiết lộ cho bất kỳ ai.
\end{itemize}

\noindent\textit{Quy ước ký hiệu cặp khóa:} Trong toàn bộ báo cáo này, hai hệ ký hiệu cặp khóa $(PU, PR)$ (Public Key / Private Key) và $(PK, SK)$ (Public Key / Secret Key) được hiểu hoàn toàn tương đương nhau về bản chất mật mã học. Nhằm giữ tính trực quan và gắn liền với kịch bản truyền thông kinh điển giữa các chủ thể (Alice, Bob và Eve), Chương 1 sử dụng hệ ký hiệu $(PU_B, PR_B)$ và $(PU_E, PR_E)$. Từ Chương 2 trở đi, khi đi sâu vào đặc tả kỹ thuật thuật toán, cấu trúc chứng thư X.509 và giao thức TLS 1.3, tài liệu sử dụng thống nhất ký hiệu chuẩn mực $(PK, SK)$ (ví dụ: $SK_{\mathrm{CA}}$, $PK_{\mathrm{server}}$).

Xét kịch bản trao đổi thông tin kinh điển: \textbf{Alice} là bên gửi thông điệp, \textbf{Bob} là bên nhận thông điệp, và \textbf{Eve} là đối thủ tiềm tàng có khả năng giám sát, đánh chặn hoặc sửa đổi các gói tin lưu thông trên mạng. Để gửi một thông điệp bí mật đến Bob, Alice cần sử dụng khóa công khai của Bob để mã hóa. Tuy nhiên, một bài toán then chốt xuất hiện: \textit{Làm thế nào để Alice xác minh rằng khóa công khai mà mình nhận được qua mạng thực sự thuộc về Bob, chứ không phải do đối thủ Eve giả mạo?}

\subsection{Quá trình trao đổi khóa trong điều kiện bình thường}
Trong điều kiện lý tưởng, khi kênh truyền thông không bị can thiệp và khóa công khai được chuyển giao một cách toàn vẹn:
\begin{enumerate}
    \item Bob khởi tạo cặp khóa cá nhân gồm khóa công khai $PU_B$ và khóa bí mật $PR_B$.
    \item Bob cung cấp khóa công khai $PU_B$ cho Alice qua kênh mạng.
    \item Alice sử dụng khóa công khai $PU_B$ của Bob cùng hàm mã hóa $E(\cdot)$ để mã hóa thông điệp $M$, tạo thành bản mã $C$:
    \begin{equation}
        C = E(PU_B, M)
    \end{equation}
    \item Bob nhận bản mã $C$ và sử dụng khóa bí mật duy nhất tương ứng $PR_B$ cùng hàm giải mã $D(\cdot)$ để khôi phục lại thông điệp gốc $M$:
    \begin{equation}
        M = D(PR_B, C)
    \end{equation}
\end{enumerate}

Mối quan hệ toán học giữa hai quá trình được tóm lược bằng biểu thức:
\begin{equation}
    C = E(PU_B, M), \qquad M = D(PR_B, C).
\end{equation}

Quy trình này chỉ đạt được mục tiêu bảo mật dự kiến khi hai điều kiện tiên quyết được thỏa mãn đồng thời: (1) Alice nhận đúng khóa công khai $PU_B$ của Bob, và (2) giao thức cùng thuật toán mã hóa được sử dụng là an toàn. Nếu Alice bị lừa sử dụng một chiếc khóa khác, tính bí mật của toàn bộ thông điệp sẽ bị vô hiệu hóa hoàn toàn.

\subsection{Cuộc tấn công thay thế khóa của Eve}
Khi việc phân phối khóa công khai diễn ra trên kênh truyền mở mà thiếu cơ chế xác thực nguồn gốc, đối thủ Eve có thể thực hiện một cuộc tấn công thay thế khóa theo phương thức người đứng giữa (\textit{Man-in-the-Middle - MITM}). Giả sử Eve sở hữu một cặp khóa của riêng mình gồm khóa công khai $PU_E$ và khóa bí mật $PR_E$. Cuộc tấn công diễn ra theo trình tự 6 bước:
\begin{enumerate}
    \item Bob gửi khóa công khai $PU_B$ của mình cho Alice qua mạng.
    \item Eve chặn thông điệp chứa $PU_B$ trên đường truyền và thay thế bằng khóa công khai của mình là $PU_E$.
    \item Alice nhận được $PU_E$ nhưng nhầm tưởng đó là khóa của Bob; do đó Alice sử dụng $PU_E$ để mã hóa thông điệp $M$:
    \begin{equation}
        C_1 = E(PU_E, M)
    \end{equation}
    \item Alice gửi bản mã $C_1$ đi. Eve chặn bắt $C_1$ và sử dụng khóa bí mật $PR_E$ của mình để giải mã:
    \begin{equation}
        M = D(PR_E, C_1)
    \end{equation}
    Tại thời điểm này, Eve đọc được trọn vẹn nội dung thông điệp gốc $M$ của Alice.
    \item Để duy trì sự ngụy trang, Eve sử dụng khóa công khai thật sự của Bob ($PU_B$) để mã hóa lại thông điệp $M$:
    \begin{equation}
        C_2 = E(PU_B, M)
    \end{equation}
    sau đó Eve chuyển tiếp bản mã $C_2$ tới Bob dưới danh nghĩa Alice.
    \item Bob nhận được $C_2$, sử dụng khóa bí mật $PR_B$ của mình để giải mã:
    \begin{equation}
        M = D(PR_B, C_2)
    \end{equation}
    Bob khôi phục được đúng thông điệp $M$ ban đầu và hoàn toàn không nhận ra cuộc tấn công nếu không có cơ chế xác thực thích hợp.
\end{enumerate}

Ý nghĩa của các biểu thức trong cuộc tấn công:
\begin{itemize}
    \item $C_1 = E(PU_E, M)$ thể hiện bản mã tạo ra do Alice mã hóa bằng khóa công khai của kẻ tấn công Eve thay vì của Bob.
    \item $M = D(PR_E, C_1)$ chứng minh Eve có khả năng giải mã trực tiếp và thu thập nội dung thông tin mật nhờ nắm giữ khóa bí mật $PR_E$.
    \item $C_2 = E(PU_B, M)$ thể hiện việc Eve tái mã hóa thông điệp bằng khóa công khai hợp lệ của Bob để che đậy sự can thiệp của mình.
\end{itemize}

Cuộc tấn công người đứng giữa (MITM) này thành công vì Eve đã lợi dụng việc Alice hoàn toàn không có khả năng xác minh nguồn gốc và tính hợp lệ của khóa công khai nhận được. Mô hình trao đổi bình thường và cuộc tấn công thay thế khóa được minh họa tổng hợp tại Hình~\ref{fig:mitm-key-exchange}.

\begin{figure}[!htbp]
\centering
\resizebox{0.96\textwidth}{!}{%
\begin{tikzpicture}[
    scale=0.94,
    sectiontitle/.style={anchor=west, font=\bfseries\normalsize, text=black},
    actor/.style={draw=black, fill=gray!8, thick, rounded corners=5pt, minimum width=3.2cm, minimum height=1.1cm, align=center, font=\small, text=black},
    eve/.style={draw=black, fill=gray!22, very thick, double, rounded corners=5pt, minimum width=4.0cm, minimum height=1.1cm, align=center, font=\small, text=black},
    note/.style={draw=black!75, fill=white, font=\scriptsize, align=center, inner sep=5pt, rounded corners=3pt, text=black},
    evenote/.style={draw=black, fill=gray!12, dashed, font=\scriptsize, align=center, inner sep=5pt, rounded corners=3pt, text=black},
    warnnote/.style={draw=black, fill=gray!18, thick, font=\scriptsize, align=center, inner sep=5pt, rounded corners=4pt, text=black},
    msg/.style={-{Stealth[scale=1.15]}, thick, black},
    evemsg/.style={-{Stealth[scale=1.15]}, very thick, dashed, black},
    lifeline/.style={thick, dotted, black!50},
    evelifeline/.style={thick, dashed, black!65}
]

    % ==============================================================
    % (a) KỊCH BẢN BÌNH THƯỜNG (LÝ TƯỞNG)
    % ==============================================================
    \node[sectiontitle] at (0, 13.0) {(a) Trao đổi bình thường qua kênh truyền không bảo vệ (Kịch bản lý tưởng)};

    % Diễn viên
    \node[actor] (bobA) at (2.2, 11.6) {\textbf{Bob}\\[2pt]{\scriptsize Cặp khóa: $(PU_B, PR_B)$}};
    \node[actor] (aliceA) at (12.3, 11.6) {\textbf{Alice}\\[2pt]{\scriptsize Nhận đúng khóa $PU_B$}};

    % Đường thời gian (Lifelines)
    \draw[lifeline] (2.2, 11.0) -- (2.2, 7.9);
    \draw[lifeline] (12.3, 11.0) -- (12.3, 7.9);

    % Bước 1: Bob gửi PU_B
    \draw[msg] (2.2, 9.9) -- (12.3, 9.9)
        node[midway, above=3pt, font=\footnotesize] {1. Bob gửi khóa công khai $PU_B$};

    % Bước 2: Alice gửi bản mã C
    \draw[msg] (12.3, 8.6) -- (2.2, 8.6)
        node[midway, above=3pt, font=\footnotesize] {2. Alice gửi bản mã $C = E(PU_B, M)$};

    % Ghi chú giải mã & xác thực
    \node[note, text width=4.0cm] at (2.2, 7.3) {
        \textbf{Bob giải mã bằng $PR_B$:}\\
        $M = D(PR_B, C)$
    };
    \node[note, text width=4.0cm] at (12.3, 7.3) {
        \textbf{Alice an tâm:}\\
        Chỉ người có $PR_B$ mới giải mã được
    };

    % Đường phân cách hai kịch bản
    \draw[dash pattern=on 5pt off 4pt, black!50, thick] (0, 6.1) -- (14.5, 6.1);

    % ==============================================================
    % (b) KỊCH BẢN TẤN CÔNG MAN-IN-THE-MIDDLE (EVE TRÁO ĐỔI KHÓA)
    % ==============================================================
    \node[sectiontitle] at (0, 5.3) {(b) Kịch bản tấn công Man-In-The-Middle (Eve can thiệp và tráo đổi khóa)};

    % Diễn viên
    \node[actor] (bobB) at (1.8, 3.9) {\textbf{Bob}\\[2pt]{\scriptsize Cặp khóa: $(PU_B, PR_B)$}};
    \node[eve]   (eveB) at (7.25, 3.9) {\textbf{Eve (Kẻ tấn công MITM)}\\[2pt]{\scriptsize Cặp khóa giả mạo: $(PU_E, PR_E)$}};
    \node[actor] (aliceB) at (12.7, 3.9) {\textbf{Alice}\\[2pt]{\scriptsize Nhầm $PU_E$ là khóa của Bob}};

    % Đường thời gian (Lifelines)
    \draw[lifeline] (1.8, 3.3) -- (1.8, -4.5);
    \draw[evelifeline] (7.25, 3.3) -- (7.25, 2.85);
    \draw[evelifeline] (7.25, 2.15) -- (7.25, -0.95);
    \draw[lifeline] (12.7, 3.3) -- (12.7, -4.5);

    % Bước 1: Bob gửi PU_B -> Eve chặn
    \node[evenote, text width=2.6cm] (chaneve) at (7.25, 2.5) {\textbf{Eve chặn bắt $PU_B$}};
    \draw[msg] (1.8, 2.5) -- (chaneve.west)
        node[midway, above=3pt, font=\footnotesize] {1. Bob gửi khóa $PU_B$};

    % Bước 2: Eve thay bằng PU_E gửi Alice
    \draw[evemsg] (7.25, 1.1) -- (12.7, 1.1)
        node[midway, above=3pt, font=\footnotesize] {2. Gửi $PU_E$ (mạo danh Bob)};

    % Bước 3: Alice gửi C1 về cho Eve (nhầm là Bob)
    \draw[msg] (12.7, -0.3) -- (7.25, -0.3)
        node[midway, above=3pt, font=\footnotesize] {3. Alice gửi bản mã $C_1 = E(PU_E, M)$};

    % Hộp Eve giải mã, đọc trộm và mã hóa lại
    \node[evenote, text width=6.6cm] (evedecrypt) at (7.25, -1.8) {
        \textbf{Eve giải mã bằng $PR_E$ \& đọc trộm:}\\
        $M = D(PR_E, C_1)$ \textit{(Thu được thông điệp gốc $M$)}\\[3pt]
        \textbf{Eve mã hóa lại bằng $PU_B$ thật của Bob:}\\
        $C_2 = E(PU_B, M)$
    };

    % Bước 4: Eve gửi C2 cho Bob
    \draw[evemsg] (7.25, -3.5) -- (1.8, -3.5)
        node[midway, above=3pt, font=\footnotesize] {4. Chuyển tiếp bản mã $C_2 = E(PU_B, M)$};

    % Hàng kết quả cuối cùng: 3 khối cân xứng
    \node[note, text width=3.2cm] at (1.8, -5.3) {
        \textbf{Bob giải mã ($PR_B$):}\\
        $M = D(PR_B, C_2)$\\
        {\scriptsize (Nhận đúng $M$, không nghi ngờ)}
    };

    \node[warnnote, text width=6.2cm] at (7.25, -5.3) {
        \textbf{Hậu quả an ninh nghiêm trọng:}\\
        Eve đọc/sửa trộm $M$ mà hai bên hoàn toàn không biết\\
        do thiếu cơ chế xác thực danh tính khóa công khai!
    };

    \node[note, text width=3.2cm] at (12.7, -5.3) {
        \textbf{Alice tin tưởng:}\\
        Tin Bob nhận được $C_1$\\
        {\scriptsize (Không biết bị tráo khóa $PU_E$)}
    };

\end{tikzpicture}%
}

\caption{Mô hình trao đổi thông tin thông thường và kịch bản tấn công thay thế khóa của Eve.}
\label{fig:mitm-key-exchange}
\end{figure}

\subsection{Phân tích bản chất bài toán và vai trò của PKI}
\label{subsec:analysis-pki-role}
Phân tích kịch bản tấn công của Eve làm bộc lộ những bài học cơ bản về an toàn mật mã:
\begin{itemize}
    \item \textbf{Phân biệt tính bí mật của dữ liệu và tính xác thực của khóa công khai:} Mật mã hóa chỉ đảm bảo tính bí mật (confidentiality) của thông điệp đối với đối tượng sở hữu khóa bí mật tương ứng. Tuy nhiên, tính bí mật đó trở nên vô nghĩa nếu bản thân chiếc khóa công khai được dùng để mã hóa không bảo đảm tính xác thực (authenticity) và không có căn cứ xác thực nguồn gốc (origin authenticity).
    \item \textbf{Nghịch lý toán học trong bảo mật hệ thống:} Các thuật toán mật mã cơ sở (như các hàm mã hóa phi đối xứng) vẫn hoạt động hoàn toàn chính xác theo đúng các định lý toán học ($C_1$ giải mã chính xác bằng $PR_E$, và $C_2$ giải mã chính xác bằng $PR_B$). Mặc dù thuật toán không hề bị bẻ gãy, hệ thống bảo mật tổng thể vẫn bị phá vỡ vì khóa đưa vào thuật toán không thuộc về đúng đối tượng dự kiến.
    \item \textbf{Các nguy cơ và hậu quả bảo mật:}
    \begin{itemize}
        \item \textit{Lộ lọt dữ liệu bí mật:} Kẻ tấn công có thể đọc trộm toàn bộ thông điệp trao đổi giữa các bên.
        \item \textit{Giả mạo danh tính (Impersonation):} Kẻ tấn công có thể thay phiên đóng vai cả hai bên đối thoại một cách trót lọt.
        \item \textit{Sửa đổi và thao túng thông điệp (Tampering):} Eve không chỉ dừng lại ở việc đọc trộm mà còn có thể can thiệp, sửa đổi nội dung $M$ thành $M'$ trước khi mã hóa lại thành $C_2 = E(PU_B, M')$, gây sai lệch hoàn toàn dữ liệu nếu giao thức không có cơ chế bảo vệ tính toàn vẹn và chữ ký số thích hợp.
    \end{itemize}
\end{itemize}

Do đó, thách thức cốt lõi đặt ra là: \textbf{Phải có một cơ chế đáng tin cậy để liên kết (bind) một khóa công khai cụ thể với danh tính thực tế của chủ sở hữu hợp pháp}.

Để thiết lập mối liên kết đáng tin cậy giữa danh tính và khóa công khai mà không đòi hỏi các bên phải gặp mặt trực tiếp hoặc thiết lập kênh bí mật trước, \textbf{Hạ tầng khóa công khai (Public Key Infrastructure - PKI)} đã được chuẩn hóa và triển khai như một giải pháp nền tảng. PKI là một hệ thống hạ tầng toàn diện bao gồm phần cứng, phần mềm, chính sách và tiêu chuẩn hỗ trợ quản lý, phân phối, lưu trữ và xác thực mối liên hệ giữa khóa công khai với danh tính của các thực thể thông qua các cấu phần và cơ chế then chốt:
\begin{itemize}
    \item \textbf{Tổ chức chứng thực số (Certificate Authority - CA):} Là bên thứ ba đáng tin cậy (Trusted Third Party). CA có trách nhiệm thẩm tra tính chính xác về danh tính của đối tượng yêu cầu trước khi cấp phát chứng thư số.
    \item \textbf{Chứng thư số (Digital Certificate):} Là cấu trúc dữ liệu điện tử (thường theo chuẩn quốc tế ITU-T X.509) đóng gói khóa công khai cùng với các thông tin định danh của chủ thể (Subject) sở hữu khóa đó.
    \item \textbf{Chữ ký số (Digital Signature):} CA sử dụng khóa bí mật của chính mình để tạo chữ ký số trên toàn bộ nội dung của chứng thư, bảo đảm tính toàn vẹn và chống chối bỏ cho mối liên kết giữa danh tính và khóa công khai.
\end{itemize}

Trong môi trường có sự hỗ trợ của PKI, khi Alice muốn liên lạc với Bob, thay vì nhận một khóa công khai trần dễ bị đánh tráo, Alice sẽ nhận được một chứng thư số của Bob do một CA đáng tin cậy ban hành. Alice có thể thực hiện quy trình kiểm tra toàn diện:
\begin{enumerate}
    \item \textbf{Kiểm tra chữ ký của CA:} Alice sử dụng khóa công khai của CA (vốn đã được cài đặt sẵn và tin cậy trước trong hệ thống) để xác thực chữ ký số trên chứng thư của Bob.
    \item \textbf{Xác minh danh tính chủ thể:} Alice kiểm tra xem thông tin định danh (như tên chủ thể, tên miền) có khớp chính xác với Bob hay không.
    \item \textbf{Kiểm tra thời hạn hiệu lực:} Alice đối chiếu thời điểm hiện tại với khoảng thời gian hợp lệ được quy định trên chứng thư (các trường \texttt{notBefore} và \texttt{notAfter}).
    \item \textbf{Kiểm tra trạng thái thu hồi:} Alice kiểm tra xem chứng thư có bị thu hồi trước thời hạn (ví dụ do lộ khóa bí mật) hay không thông qua danh sách thu hồi chứng thư (CRL) hoặc giao thức trạng thái chứng thư trực tuyến (OCSP) theo các chính sách áp dụng.
\end{enumerate}

\noindent\textbf{Giới hạn của PKI và sự cố thực tế kinh điển (DigiNotar 2011):}
Cần khẳng định rằng, PKI không phải là một giải pháp loại bỏ hoàn toàn mọi nguy cơ an ninh. PKI về thực chất là một hạ tầng quản trị niềm tin (\textit{trust management}); an ninh của hệ thống bị ràng buộc chặt chẽ bởi các chính sách vận hành và tính toàn vẹn của các CA gốc (Trust Anchors). Nếu một CA bị xâm nhập hoặc cấp sai chứng thư, toàn bộ chuỗi tin cậy sẽ bị lợi dụng để hợp thức hóa các hành vi tấn công mà trình duyệt máy trạm khó có thể tự phát hiện bằng các kiểm tra toán học thông thường.

Minh chứng điển hình cho giới hạn này là \textbf{sự cố DigiNotar (2011)}~\cite{prins2011diginotar}, khi một CA công cộng của Hà Lan bị tin tặc chiếm quyền và phát hành hàng loạt chứng thư giả mạo (\textit{rogue certificates}) có chữ ký mật mã hoàn toàn hợp lệ cho tên miền \texttt{*.google.com}. Kẻ tấn công đã sử dụng các chứng thư này để triển khai tấn công người đứng giữa (MITM) đối với hàng trăm nghìn người dùng Internet mà không hề làm phát sinh cảnh báo lỗi chứng thư, buộc các hãng công nghệ lớn phải xóa bỏ hoàn toàn DigiNotar khỏi Root Trust Store. Sự cố khẳng định một nguyên lý cốt lõi: chữ ký mật mã đúng là điều kiện cần nhưng chưa đủ để bảo đảm an toàn nếu niềm tin vào CA bị vi phạm (phân tích chi tiết về vectơ tấn công và bài học kiến trúc từ sự cố này sẽ được mổ xẻ tại Mục~\ref{sec:threat-model} ở Chương~\ref{chap:threat-security}).

Chính vì vậy, bài toán nghiên cứu trọng tâm được xác định là: \textbf{Làm thế nào để phân phối và xác thực khóa công khai một cách đáng tin cậy, hạn chế tối đa nguy cơ giả mạo và tấn công thay thế khóa, đồng thời đánh giá được các điều kiện chấp nhận và giới hạn an ninh của hệ thống trong thực tế triển khai?} Đây là động lực thúc đẩy các mục tiêu và câu hỏi nghiên cứu cụ thể trong đề tài này.

\section{Mục tiêu và câu hỏi nghiên cứu}
\label{sec:objectives-rqs}
Mục tiêu của đề tài là phân tích kiến trúc hạ tầng khóa công khai (PKI), tìm hiểu cơ chế xác thực chứng thư số theo tiêu chuẩn X.509 và xây dựng một môi trường thực nghiệm (lab) nhằm kiểm chứng các điều kiện chấp nhận hoặc từ chối chứng thư. Thông qua đó, đề tài làm rõ cách PKI thiết lập mối liên hệ tin cậy giữa danh tính và khóa công khai, đồng thời đánh giá vai trò của cơ chế thu hồi chứng thư trong việc bảo đảm an toàn thông tin.

Để đạt được mục tiêu trên, đề tài tập trung giải quyết các câu hỏi nghiên cứu sau:
\begin{enumerate}
  \item \textbf{RQ1:} PKI liên kết danh tính với khóa công khai bằng cách nào?
  \item \textbf{RQ2:} Những điều kiện nào khiến chứng thư X.509 được chấp nhận hoặc từ chối?
  \item \textbf{RQ3:} CRL và OCSP giúp giảm thiểu rủi ro liên quan đến chứng thư bị thu hồi như thế nào?
  \item \textbf{RQ4:} Những điểm yếu nào vẫn tồn tại trong PKI ngay cả khi các thuật toán mật mã được sử dụng vẫn bảo đảm an toàn?
\end{enumerate}

Trong đó, mối liên hệ giữa các câu hỏi nghiên cứu với cấu trúc các chương trong báo cáo được phân bổ cụ thể:
\begin{itemize}
    \item \textbf{RQ1} được làm sáng tỏ qua phân tích lý thuyết về kiến trúc PKI, mô hình tin cậy và cấu trúc chứng thư X.509 ở Chương~\ref{chap:kien-truc-pki} và Chương~\ref{chap:x509-xac-thuc};
    \item \textbf{RQ2} và \textbf{RQ3} được đối chiếu giữa lý thuyết thuật toán xác thực đường dẫn, cơ chế thu hồi (Chương~\ref{chap:x509-xac-thuc}) và kiểm chứng trực quan bằng thực nghiệm 10 ca kiểm thử chuyên sâu (TC01--TC10) ở Chương~\ref{chap:experiment};
    \item \textbf{RQ4} được mổ xẻ chuyên sâu qua mô hình hóa rủi ro, phân tích giới hạn của chữ ký hợp lệ tại Chương~\ref{chap:threat-security} và các hướng tiếp cận nâng cao (như Certificate Transparency, tự động hóa ACME) ở Chương~\ref{chap:ung-dung}.
\end{itemize}

Thông qua việc trả lời các câu hỏi nghiên cứu và kiểm chứng bằng thực nghiệm, đề tài hướng đến việc làm rõ các yếu tố quyết định độ tin cậy của PKI, nhận diện những hạn chế trong quá trình xác thực và quản lý chứng thư số, từ đó nâng cao hiểu biết về khả năng ứng dụng cũng như các rủi ro an toàn thông tin của hệ thống PKI trong thực tế.

\section{Phạm vi và phương pháp nghiên cứu}
\label{sec:scope-methodology}
\begin{itemize}
    \item \textbf{Phạm vi nghiên cứu:} Đề tài tập trung nghiên cứu mô hình PKI phân cấp chuẩn ITU-T X.509 v3 và cơ chế xác thực máy chủ (Server Authentication) trong giao thức TLS trên môi trường phòng thí nghiệm nội bộ (lab). Trọng tâm nghiên cứu bao gồm cấu trúc các trường chứng thư và phần mở rộng quan trọng (SAN, Basic Constraints, Key Usage), thuật toán xác minh chuỗi tin cậy, hai cơ chế kiểm tra thu hồi (CRL và OCSP), cùng các mô hình đe dọa và hạn chế an ninh cố hữu của PKI.
    \item \textbf{Phương pháp nghiên cứu:} Đề tài kết hợp phương pháp nghiên cứu tài liệu chuẩn hóa quốc tế (RFC 5280, RFC 6960, RFC 8446, RFC 3647) với phương pháp thực nghiệm khoa học. Nhóm nghiên cứu thiết kế và triển khai một hạ tầng PKI hoàn chỉnh bằng công cụ OpenSSL, xây dựng chuỗi chứng thư nhiều cấp và thực thi 10 ca kiểm thử độc lập (TC01--TC10) nhằm ghi nhận, đo lường và đối chiếu trực quan hành vi thẩm định chứng thư thực tế với các ràng buộc lý thuyết.
\end{itemize}

\section{Cấu trúc của báo cáo}
\label{sec:report-structure}
Báo cáo nghiên cứu được bố cục thành 7 chương chuyên môn và phần kết luận, cụ thể như sau:
\begin{itemize}
    \item \textbf{Chương 1 -- Giới thiệu và xác định vấn đề nghiên cứu:} Trình bày bối cảnh ra đời của mật mã khóa công khai, phân tích lỗ hổng cố hữu trong phân phối khóa trần qua kịch bản tấn công người đứng giữa (MITM), sự cần thiết của PKI, mục tiêu, 4 câu hỏi nghiên cứu (RQ1--RQ4), phạm vi triển khai và phương pháp luận của đề tài.
    \item \textbf{Chương 2 -- Cơ sở mật mã khóa công khai và PKI:} Trình bày cơ sở lý thuyết mật mã bất đối xứng, hàm băm mật mã, chữ ký số, quản lý vòng đời khóa và vai trò của các lớp thuật toán trong giao thức TLS 1.3.
    \item \textbf{Chương 3 -- Kiến trúc PKI và vòng đời chứng thư:} Phân tích kiến trúc tổng thể của PKI, các thành phần và mô hình tin cậy, quy trình quản lý vòng đời chứng thư và chính sách vận hành (CP/CPS) theo khung RFC 3647.
    \item \textbf{Chương 4 -- X.509, chuỗi tin cậy và cơ chế xác thực:} Đi sâu vào cấu trúc chứng thư số chuẩn ITU-T X.509 v3, thuật toán xây dựng và thẩm định đường dẫn chứng thư, cùng cơ chế kiểm tra trạng thái thu hồi (CRL và OCSP).
    \item \textbf{Chương 5 -- Phân tích bảo mật, mô hình tấn công và so sánh:} Mô hình hóa đe dọa, so sánh đối chiếu các mô hình tin cậy (PKI, Web of Trust, TOFU), và phân tích giới hạn của chữ ký hợp lệ thông qua nghiên cứu điển hình về sự cố DigiNotar (2011).
    \item \textbf{Chương 6 -- Thiết kế và đánh giá thực nghiệm:} Thiết kế môi trường lab OpenSSL, kết quả thực nghiệm 10 ca kiểm thử (TC01--TC10) và các thảo luận, đánh giá chuyên sâu.
    \item \textbf{Chương 7 -- Ứng dụng và hướng phát triển PKI:} Khảo sát các ứng dụng thực tế của PKI, cơ chế minh bạch chứng thư (Certificate Transparency - CT), tự động hóa quản lý chứng thư qua giao thức ACME và định hướng chuyển đổi sang mật mã hậu lượng tử (PQC).
    \item \textbf{Kết luận và hướng phát triển:} Tổng kết các phát hiện chính, đối chiếu và giải đáp 4 câu hỏi nghiên cứu, chỉ rõ các giới hạn của mô hình thực nghiệm và đề xuất hướng nghiên cứu mở rộng trong tương lai.
\end{itemize}

\chapter{Cơ sở mật mã khóa công khai và PKI}
\label{chap:co-so-mat-ma}
\section{Yêu cầu bảo mật và quản lý khóa}
\label{sec:security-reqs-key-mgmt}

\subsection{Yêu cầu bảo mật trong truyền thông và vai trò của PKI}
Trong an toàn thông tin hiện đại, việc bảo vệ luồng truyền thông và dữ liệu giao dịch xoay quanh mô hình các mục tiêu an ninh cơ bản (mô hình CIA mở rộng):
\begin{itemize}
    \item \textbf{Tính bí mật (Confidentiality):} Đảm bảo thông tin chỉ được truy cập và đọc hiểu bởi các đối tượng có thẩm quyền (sở hữu khóa giải mã hợp lệ), ngăn chặn hành vi nghe lén hoặc giải mã trái phép từ bên thứ ba trên kênh truyền thông công cộng.
    \item \textbf{Tính toàn vẹn (Integrity):} Đảm bảo dữ liệu không bị sửa đổi, chèn thêm, xóa bớt hoặc tái sắp xếp trái phép trong quá trình lưu trữ và truyền tải. Mọi biến đổi dù là nhỏ nhất trên thông điệp đều phải được phát hiện kịp thời.
    \item \textbf{Tính xác thực (Authenticity):} Đảm bảo nguồn gốc của dữ liệu và định danh thực tế của thực thể tham gia giao tiếp là chính xác và có thể kiểm chứng được, ngăn ngừa hành vi giả mạo danh tính (\textit{impersonation}) hoặc tấn công người đứng giữa (MITM).
    \item \textbf{Tính sẵn sàng (Availability):} Đảm bảo hệ thống, tài nguyên và dịch vụ xác thực luôn sẵn sàng phục vụ các bên tin cậy khi có yêu cầu hợp lệ, không bị gián đoạn bởi sự cố kỹ thuật hoặc các cuộc tấn công từ chối dịch vụ (DoS/DDoS).
    \item \textbf{Tính chống chối bỏ (Non-repudiation):} Đảm bảo thực thể đã thực hiện một hành động (như ký duyệt văn bản hay phát hành chứng thư) không thể phủ nhận hành vi của mình đối với bên thứ ba độc lập.
\end{itemize}

Hạ tầng khóa công khai (PKI) về bản chất là một cơ chế quản trị niềm tin và định danh số. Do đó, PKI không trực tiếp giải quyết toàn bộ các mục tiêu an ninh nói trên mà có sự phân định rõ ràng giữa thuộc tính hỗ trợ trực tiếp và thuộc tính cần phối hợp với các cơ chế bảo mật khác:
\begin{itemize}
    \item \textbf{Các thuộc tính PKI hỗ trợ trực tiếp:}
    \begin{itemize}
        \item \textit{Tính xác thực (Authenticity):} Đây là mục tiêu cốt lõi của PKI. Thông qua chứng thư số chuẩn ITU-T X.509 và chữ ký số của Tổ chức chứng thực (CA), PKI liên kết chặt chẽ danh tính thực tế của chủ thể (Subject) với khóa công khai tương ứng, giảm thiểu và ngăn chặn nguy cơ phân phối khóa công khai trần không rõ nguồn gốc trong điều kiện CA được bảo vệ an toàn.
        \item \textit{Tính toàn vẹn (Integrity) và Chống chối bỏ (Non-repudiation):} Bằng việc đóng gói dữ liệu và ký số bằng khóa bí mật của CA hoặc của chủ thể, PKI bảo đảm tính toàn vẹn của chứng thư số và thông điệp giao dịch; đồng thời ngăn chặn bên ký chối bỏ trách nhiệm về dữ liệu đã ban hành.
    \end{itemize}
    \item \textbf{Các thuộc tính cần kết hợp cơ chế bảo vệ khác:}
    \begin{itemize}
        \item \textit{Tính bí mật (Confidentiality):} Bản thân PKI và chứng thư X.509 không trực tiếp thực hiện việc mã hóa dữ liệu ứng dụng. Khóa công khai trong chứng thư chủ yếu phục vụ xác thực danh tính và thỏa thuận khóa; việc bảo vệ bí mật cho luồng dữ liệu đòi hỏi các giao thức truyền thông an toàn (như TLS 1.3~\cite{rfc8446}) sử dụng cơ chế trao đổi khóa phiên (như ECDHE) kết hợp thuật toán mã hóa đối xứng khối lượng lớn (như AES-GCM hoặc ChaCha20-Poly1305).
        \item \textit{Tính sẵn sàng (Availability):} Dù PKI định nghĩa các dịch vụ kiểm tra trạng thái chứng thư (CRL, OCSP), tính sẵn sàng của hạ tầng phụ thuộc hoàn toàn vào kiến trúc mạng và giải pháp phân phối (như cân bằng tải, cơ chế lưu đệm OCSP Stapling, mạng phân phối nội dung CDN) nhằm tránh nghẽn cổ chai và ngăn ngừa tấn công làm tê liệt điểm kiểm tra thu hồi.
    \end{itemize}
\end{itemize}

Bảng~\ref{tab:pki-security-goals} tóm tắt phạm vi hỗ trợ của PKI đối với các mục tiêu an ninh thông tin.

\begin{table}[!htbp]
\centering
\small
\renewcommand{\arraystretch}{1.3}
\begin{tabular}{>{\raggedright\arraybackslash}p{3.5cm}>{\raggedright\arraybackslash}p{2.6cm}>{\raggedright\arraybackslash}p{8.5cm}}
\toprule
\textbf{Mục tiêu an ninh} & \textbf{Vai trò của PKI} & \textbf{Cơ chế thực hiện / Phối hợp} \\
\midrule
\textbf{Xác thực (Authenticity)} & Hỗ trợ trực tiếp & Chứng thư số X.509, chữ ký số của CA trên Subject PublicKey. \\
\midrule
\textbf{Toàn vẹn (Integrity)} & Hỗ trợ trực tiếp & Hàm băm mật mã và chữ ký số trên cấu trúc chứng thư. \\
\midrule
\textbf{Chống chối bỏ (Non-repudiation)} & Hỗ trợ trực tiếp & Chữ ký số gắn liền với khóa bí mật duy nhất của thực thể. \\
\midrule
\textbf{Bí mật (Confidentiality)} & Kết hợp bên ngoài & Giao thức TLS bắt tay thỏa thuận khóa phiên (ECDHE) và mã hóa đối xứng (AES-GCM). \\
\midrule
\textbf{Sẵn sàng (Availability)} & Kết hợp bên ngoài & Hạ tầng mạng phân tán, bộ đệm chứng thực, OCSP Stapling, CDN. \\
\bottomrule
\end{tabular}
\caption[Phân định phạm vi hỗ trợ của PKI đối với mục tiêu an ninh thông tin]{Phân định phạm vi hỗ trợ của PKI đối với các mục tiêu an ninh thông tin.}
\label{tab:pki-security-goals}
\end{table}



\subsection{Nguyên lý quản lý khóa và vòng đời khóa mật mã}
\label{subsec:key-lifecycle}
Theo \textit{Nguyên lý Kerckhoffs}, độ an toàn của một hệ thống mật mã không phụ thuộc vào việc giữ bí mật thuật toán mà nằm hoàn toàn ở sự an toàn của khóa mật mã. Trong hệ thống khóa bất đối xứng, hai thành phần của cặp khóa có yêu cầu bảo vệ trái ngược nhau:
\begin{itemize}
    \item \textbf{Khóa công khai (Public Key):} Đòi hỏi tính \textit{toàn vẹn} và tính \textit{xác thực} nguồn gốc tuyệt đối (đây chính là nhiệm vụ của chứng thư số PKI), nhưng không cần giữ bí mật.
    \item \textbf{Khóa bí mật (Private Key):} Đòi hỏi tính \textit{bí mật} và \textit{bảo vệ quyền truy cập} tuyệt đối. Nếu khóa bí mật bị lộ, toàn bộ các bảo đảm an ninh về xác thực và chống chối bỏ đều sụp đổ.
\end{itemize}

Quản lý khóa (Key Management) là tập hợp toàn bộ các chính sách và kỹ thuật nhằm kiểm soát cặp khóa xuyên suốt vòng đời (Key Lifecycle):
\begin{enumerate}
    \item \textbf{Khởi tạo khóa (Generation):} Phải sử dụng bộ sinh số ngẫu nhiên giả an toàn mật mã (CSPRNG) và tuân thủ độ dài khóa khuyến nghị (ví dụ RSA $\ge$ 2048-bit, ECC $\ge$ 256-bit) để chống thám mã vét cạn.
    \item \textbf{Lưu trữ và bảo vệ khóa bí mật (Storage):} Khóa bí mật của các CA gốc và Intermediate CA bắt buộc phải được tạo sinh và lưu trữ trong các mô-đun phần cứng an toàn chuyên dụng (Hardware Security Module -- HSM) với chính sách kiểm soát truy cập nghiêm ngặt (quản trị phân quyền ngưỡng $m$-trên-$n$). Khóa tuyệt đối không được xuất ra ngoài bộ nhớ dưới dạng bản rõ.
    \item \textbf{Phân phối và chứng thực (Distribution):} Khóa công khai được đóng gói thành yêu cầu cấp chứng thư (CSR) và gửi tới CA để thẩm định, cấp phát chứng thư X.509.
    \item \textbf{Thời hạn sử dụng an toàn của khóa (Cryptoperiod):} Mọi khóa mật mã đều cần được ấn định thời hạn sử dụng an toàn (Cryptoperiod) nhằm giới hạn khối lượng bản mã tích lũy cho đối thủ thám mã và hạn chế rủi ro thỏa hiệp khóa theo thời gian. Cần phân biệt rõ Cryptoperiod (vòng đời mật mã của khóa) với thời hạn hiệu lực của chứng thư số (Certificate Validity Period xác định bởi các trường \texttt{notBefore} và \texttt{notAfter}), vốn là khoảng thời gian chứng thư được CA bảo chứng về mối liên kết giữa định danh và khóa công khai (như phân tích tại Mục~\ref{sec:certificate-lifecycle}).
    \item \textbf{Thu hồi và tiêu hủy (Revocation \& Destruction):} Khi khóa bí mật nghi ngờ bị lộ, chủ sở hữu phải lập tức yêu cầu CA đưa chứng thư vào danh sách thu hồi (CRL/OCSP). Khi hết vòng đời, khóa bí mật phải bị xóa hủy an toàn (zeroize) để ngăn ngừa phục hồi dữ liệu từ phần cứng.
\end{enumerate}


\section{Khóa công khai, hàm băm và chữ ký số}
\label{sec:crypto_primitives}

Để phân tích chính xác cơ chế hoạt động của hạ tầng PKI, cần phân định rành mạch giữa các nguyên hàm mật mã cơ sở gồm cặp khóa bất đối xứng, hàm băm một chiều, chữ ký số và chứng thư số:

\begin{itemize}
    \item \textbf{Cặp khóa công khai và khóa bí mật ($PK, SK$):} Trong mật mã bất đối xứng, mỗi thực thể sở hữu một cặp khóa có quan hệ toán học chặt chẽ. Khóa công khai ($PK$) được phổ biến rộng rãi cho toàn bộ mạng lưới, trong khi khóa bí mật ($SK$) phải được lưu trữ và bảo vệ tuyệt đối bởi chủ sở hữu. Tuy nhiên, bản thân cặp khóa chỉ tồn tại dưới dạng các giá trị nhị phân trừu tượng; nó hoàn toàn không chứa thông tin định danh và không tự chứng minh được danh tính thực tế của chủ thể nắm giữ khóa.
    
    \item \textbf{Hàm băm mật mã ($H$):} Là hàm toán học một chiều ánh xạ một thông điệp $m$ có kích thước tùy ý thành một chuỗi bit có độ dài cố định gọi là giá trị tóm lược (digest): $h = H(m)$. Hàm băm mật mã chuẩn bắt buộc phải thỏa mãn ba tính chất: kháng tiền ảnh (pre-image resistance), kháng tiền ảnh thứ hai (second pre-image resistance) và kháng va chạm (collision resistance). Mặc dù hàm băm đảm bảo phát hiện mọi sửa đổi dữ liệu (tính toàn vẹn), nó không thể xác thực nguồn gốc thông điệp vì không chứa yếu tố bí mật.
    
    \item \textbf{Chữ ký số (Digital Signature):} Là sự kết hợp giữa hàm băm và mật mã khóa công khai nhằm bảo đảm tính toàn vẹn, xác thực nguồn gốc và chống chối bỏ. Về mặt mô hình khái niệm, chữ ký số được định nghĩa thông qua cặp thuật toán:
    \begin{equation}
        \sigma = \mathrm{Sign}(SK, H(m))
    \end{equation}
    \begin{equation}
        \mathrm{Verify}(PK, H(m), \sigma) \in \{\mathrm{True}, \mathrm{False}\}
    \end{equation}
    
    Cần lưu ý rằng các biểu thức trên chỉ đại diện cho mô hình khái niệm trừu tượng. Trong thực tế triển khai, thao tác tính toán phụ thuộc chặt chẽ vào cấu trúc toán học của từng lược đồ chữ ký và các cơ chế đệm (padding) chuẩn hóa:
    \begin{itemize}
        \item \textit{RSA (PKCS\#1 v1.5 / RSA-PSS):} Chuẩn cổ điển PKCS\#1 v1.5 áp dụng cơ chế đệm tất định (deterministic padding); trong khi chuẩn hiện đại RSA-PSS sử dụng sơ đồ đệm xác suất (probabilistic signature scheme) với muối ngẫu nhiên (salt) nhằm ngăn chặn triệt để các dạng tấn công đại số trước khi thực hiện lũy thừa modulo: $\sigma = (H_{pad}(m))^d \pmod N$.
        \item \textit{ECDSA (Elliptic Curve Digital Signature Algorithm):} Thực hiện trên nhóm điểm của đường cong elip, yêu cầu sinh một giá trị ngẫu nhiên bí mật $k$ trên mỗi phiên ký để tạo cặp giá trị chữ ký $(r, s)$.
        \item \textit{Ed25519 (EdDSA):} Sử dụng đường cong xoắn Edwards với cơ chế sinh tham số ngẫu nhiên mang tính tất định dẫn xuất từ khóa bí mật và thông điệp. Đáng chú ý, ở chế độ PureEdDSA (RFC 8032), thuật toán ký trực tiếp trên toàn bộ thông điệp $m$ và thực hiện băm nội tại bằng SHA-512 chứ không đòi hỏi một bước băm $H(m)$ riêng biệt từ bên ngoài như các lược đồ truyền thống.
    \end{itemize}

    \item \textbf{Chứng thư số (Digital Certificate):} Chữ ký số chứng minh được một thông điệp được tạo ra bởi người sở hữu khóa $SK$, nhưng không thể trả lời câu hỏi thực thể đó là ai ngoài đời thực. Chứng thư số chuẩn ITU-T X.509 đóng vai trò đóng gói khóa công khai $PK$ cùng định danh của chủ thể (Subject) và được một bên thứ ba tin cậy (Certificate Authority - CA) xác thực bằng chữ ký số sử dụng khóa bí mật $SK_{CA}$. Bằng cách này, chứng thư số hoàn thiện mắt xích liên kết tin cậy (binding) giữa danh tính và khóa công khai.
\end{itemize}

Bản chất kỹ thuật và giới hạn của các thành phần mật mã được tổng hợp và đối chiếu tại Bảng~\ref{tab:crypto_comparison}.

\clearpage

\begin{table}[!htbp]
\centering
\caption{So sánh bản chất và vai trò của các thành phần mật mã trong PKI}
\label{tab:crypto_comparison}
\small
\renewcommand{\arraystretch}{1.3}
\begin{tabularx}{\textwidth}{@{}l >{\raggedright\arraybackslash}p{3.2cm} >{\raggedright\arraybackslash}p{2.6cm} >{\raggedright\arraybackslash}X >{\raggedright\arraybackslash}p{3.2cm}@{}}
\toprule
\textbf{Khái niệm} & \textbf{Bản chất kỹ thuật} & \textbf{Khóa bí mật ($SK$)} & \textbf{Giá trị an ninh cung cấp} & \textbf{Hạn chế nội tại} \\
\midrule
\textbf{Cặp khóa ($PK, SK$)} & 
Hai khóa có liên hệ toán học; $PK$ công khai, $SK$ bảo mật. & 
Chỉ duy nhất chủ sở hữu hợp pháp nắm giữ và bảo vệ. & 
Nền tảng cho ký số, xác thực sở hữu và thỏa thuận khóa phiên. & 
Không tự liên kết được với danh tính thực thể ngoài đời thực. \\
\midrule
\textbf{Hàm băm ($H$)} & 
Ánh xạ thông điệp tùy ý thành digest có độ dài cố định. & 
Không sử dụng khóa. & 
Bảo đảm tính toàn vẹn, tóm lược dữ liệu, kháng va chạm. & 
Không xác thực nguồn gốc do bất kỳ ai cũng tính toán được. \\
\midrule
\textbf{Chữ ký số} & 
Cặp thuật toán Sign/Verify thao tác trên giá trị băm thông điệp. & 
Ký bằng $SK$ của chủ thể; kiểm tra bằng $PK$ tương ứng. & 
Xác thực nguồn gốc dữ liệu, bảo đảm tính toàn vẹn và chống chối bỏ. & 
Chỉ xác minh tính hợp lệ của dữ liệu với $PK$, không định danh người ký. \\
\midrule
\textbf{Chứng thư số} & 
Cấu trúc dữ liệu liên kết $PK$ với danh tính, được CA ký. & 
Ký bởi $SK_{CA}$; cấu trúc chứng thư không bao giờ chứa $SK$. & 
Thiết lập liên kết tin cậy giữa khóa công khai và danh tính chủ thể. & 
Phụ thuộc độ tin cậy của CA, thời hạn hiệu lực và trạng thái thu hồi. \\
\bottomrule
\end{tabularx}
\end{table}


\clearpage

\section{Vai trò thuật toán trong TLS}
\label{sec:crypto_roles_in_tls}

Một trong những sai lầm phổ biến khi tiếp cận mật mã học là quan niệm cho rằng khóa công khai trong chứng thư số X.509 được sử dụng để mã hóa trực tiếp toàn bộ dữ liệu ứng dụng trao đổi trên kênh truyền mạng. Trong các hệ thống bảo mật hiện đại như giao thức Transport Layer Security phiên bản 1.3 (TLS 1.3, chuẩn hóa tại RFC 8446~\cite{rfc8446}), các thuật toán mật mã được phân công trách nhiệm độc lập theo từng lớp chức năng riêng biệt nhằm tối ưu hóa hiệu năng tính toán và bảo đảm an toàn tối đa:

\begin{itemize}
    \item \textbf{Xác thực danh tính (Authentication Layer):}
    \begin{itemize}
        \item \textit{Bản chất:} Chứng thư số của máy chủ chứa khóa công khai dài hạn ($PK_{server}$) và được chứng thực thông qua chữ ký số của một CA tin cậy ($SK_{CA}$) nhằm giải quyết bài toán liên kết định danh.
        \item \textit{Cơ chế trong TLS 1.3:} Trong giai đoạn bắt tay (Handshake), máy chủ gửi chứng thư X.509 cho client và phát hành thông điệp xác thực \texttt{CertificateVerify}. Thông điệp này được ký bằng khóa bí mật dài hạn ($SK_{server}$) lên toàn bộ bản ghi lịch sử (\textit{transcript}) của các gói tin bắt tay trước đó.
        \item \textit{Mục tiêu:} Sử dụng các thuật toán chữ ký số hiện đại (như RSA-PSS, ECDSA với đường cong secp256r1/secp384r1, hoặc Ed25519) nhằm chứng minh máy chủ thực sự sở hữu khóa bí mật tương ứng với chứng thư (\textit{Proof of Possession}), ngăn chặn việc đối thủ đánh cắp và tái sử dụng chứng thư; tuy nhiên bước này không chứng minh được chứng thư đó có bị CA cấp phát sai hoặc bị giả mạo hay không nếu CA đã bị xâm nhập. Khóa công khai này tuyệt đối không dùng để mã hóa lưu lượng dữ liệu.
    \end{itemize}

    \item \textbf{Thỏa thuận khóa phiên (Key Exchange / Agreement Layer):}
    \begin{itemize}
        \item \textit{Bản chất:} Nhằm thiết lập một khóa bí mật chung giữa client và server qua kênh truyền không tin cậy mà không bao giờ truyền trực tiếp khóa đó trên mạng.
        \item \textit{Cơ chế trong TLS 1.3:} Khác biệt căn bản với các phiên bản TLS cũ (vốn cho phép dùng khóa RSA của chứng thư để mã hóa và gửi khóa tiền phiên \textit{pre-master secret}), TLS 1.3 đã loại bỏ hoàn toàn cơ chế trao đổi khóa RSA tĩnh. Trong các phiên bắt tay xác thực dựa trên chứng thư số (Certificate-based Handshake), giao thức bắt buộc sử dụng cơ chế trao đổi khóa tạm thời dựa trên Diffie-Hellman: ECDHE (\textit{Ephemeral Elliptic Curve Diffie-Hellman}, như X25519) hoặc DHE (ngoại trừ chế độ chỉ sử dụng khóa chia sẻ trước PSK-only).
        \item \textit{Đặc tính Bí mật chuyển tiếp hoàn hảo (Perfect Forward Secrecy - PFS):} Cặp khóa bất đối xứng tạm thời (\textit{ephemeral key pair}) được sinh ngẫu nhiên độc lập cho từng phiên kết nối và bị hủy ngay sau khi bắt tay. Do đó, ngay cả khi khóa bí mật dài hạn của máy chủ bị xâm phạm trong tương lai, kẻ tấn công cũng không thể giải mã hồi cứu các phiên truyền thông đã lưu vết trong quá khứ.
    \end{itemize}

    \item \textbf{Mã hóa đối xứng dữ liệu ứng dụng (Record Layer Encryption):}
    \begin{itemize}
        \item \textit{Bản chất:} Dữ liệu trao đổi thực tế (HTTP requests/responses, JSON, hình ảnh) có dung lượng rất lớn. Việc dùng mật mã khóa công khai để mã hóa từng gói tin sẽ làm tê liệt hạ tầng mạng do chi phí tính toán modulo và đường cong elip cực kỳ đắt đỏ.
        \item \textit{Cơ chế trong TLS 1.3:} Giá trị bí mật chung sau ECDHE được đưa qua hàm suy dẫn khóa HKDF (\textit{HMAC-based Extract-and-Expand Key Derivation Function}) để sinh ra các khóa phiên đối xứng. Tầng Record sau đó sử dụng độc quyền các thuật toán mã hóa đối xứng có xác thực: AEAD (\textit{Authenticated Encryption with Associated Data}), tiêu biểu như AES-128-GCM, AES-256-GCM hoặc ChaCha20-Poly1305.
        \item \textit{Mục tiêu:} Bảo đảm đồng thời tính bí mật (Confidentiality) và tính toàn vẹn (Integrity) cho từng bản ghi dữ liệu ứng dụng với tốc độ xử lý phần cứng cao (tận dụng tập lệnh AES-NI) và độ trễ thấp nhất.
    \end{itemize}
\end{itemize}

Sự phân định vai trò và phối hợp giữa các thành phần thuật toán trong TLS 1.3 được tổng hợp tại Bảng~\ref{tab:tls_algorithms}.

\begin{table}[htbp]
\centering
\caption{Phân định vai trò các thuật toán mật mã trong giao thức TLS 1.3}
\label{tab:tls_algorithms}
\small
\renewcommand{\arraystretch}{1.3}
\begin{tabularx}{\textwidth}{@{}>{\raggedright\arraybackslash}p{2.8cm} >{\raggedright\arraybackslash}p{3.2cm} >{\raggedright\arraybackslash}p{3.4cm} >{\raggedright\arraybackslash}X@{}}
\toprule
\textbf{Lớp chức năng} & \textbf{Thuật toán tiêu biểu} & \textbf{Khóa mật mã sử dụng} & \textbf{Mục tiêu an ninh đáp ứng} \\
\midrule
\textbf{Xác thực (Authentication)} & 
RSA-PSS, ECDSA (secp256r1), Ed25519. & 
Khóa bất đối xứng dài hạn ($PK, SK$) trong chứng thư X.509. & 
Xác thực nguồn gốc máy chủ, ngăn chặn tấn công mạo danh và người đứng giữa (MITM). \\
\midrule
\textbf{Thỏa thuận khóa (Key Exchange)} & 
ECDHE (X25519, secp256r1), DHE. & 
Cặp khóa bất đối xứng tạm thời (ephemeral), sinh mới mỗi phiên. & 
Thiết lập bí mật chung an toàn; bảo đảm Bí mật chuyển tiếp hoàn hảo (PFS). \\
\midrule
\textbf{Mã hóa dữ liệu (Record Layer)} & 
AEAD (AES-128-GCM, AES-256-GCM, ChaCha20-Poly1305). & 
Khóa đối xứng phiên (Session Key) suy dẫn qua HKDF. & 
Bảo đảm tính bí mật và toàn vẹn dữ liệu ứng dụng với tốc độ và hiệu năng tối ưu. \\
\bottomrule
\end{tabularx}
\end{table}


Trong phạm vi của đề tài nghiên cứu này, các nội dung tiếp theo (Chương 3, Chương 4 và kết quả thực nghiệm tại Chương 6) sẽ tập trung làm rõ lớp chức năng đầu tiên: cấu trúc chứng thư X.509, quy trình thẩm định chuỗi tin cậy và cơ chế thu hồi của hạ tầng PKI.

\chapter{Kiến trúc PKI và vòng đời chứng thư}
\label{chap:kien-truc-pki}
\section{Thành phần và mô hình tin cậy}
\label{sec:pki-components-trust-model}

Hạ tầng khóa công khai (Public Key Infrastructure -- PKI) là một hệ thống tích hợp toàn diện bao gồm phần cứng, phần mềm, quy trình, nhân sự và các chính sách quản lý nhằm tạo lập, lưu trữ, phân phối và thu hồi chứng thư số theo chuẩn ITU-T X.509 (chuẩn hóa chi tiết tại RFC~5280~\cite{rfc5280} và khung chính sách RFC~3647~\cite{rfc3647}). Về bản chất, PKI giải quyết bài toán niềm tin phân tán trên môi trường mạng mở bằng cách chuyển dịch việc tin cậy lẫn nhau giữa các đối tượng ngang hàng (vốn đòi hỏi $O(n^2)$ mối quan hệ bảo mật trực tiếp) sang việc cùng đặt niềm tin vào một bên thứ ba có thẩm quyền và uy tín (Trusted Third Party).

\subsection{Các thành phần cốt lõi trong PKI}
\label{subsec:pki-core-components}

Một hệ thống PKI hoàn chỉnh bao gồm các cấu phần chức năng phối hợp chặt chẽ với nhau trong việc quản lý và duy trì chuỗi tin cậy:

\begin{itemize}
    \item \textbf{Gốc tin cậy (Trust Anchor):} Là điểm khởi đầu tuyệt đối của niềm tin mà bên tin cậy (Relying Party) chấp nhận trước một cách vô điều kiện mà không cần bất kỳ chữ ký của bên thứ ba nào xác nhận thêm. Trong thực tế triển khai, Trust Anchor thường là khóa công khai của một Root CA được lưu trữ trực tiếp, an toàn trong kho chứng thư tin cậy (Trust Store) của hệ điều hành (như Windows, Linux, macOS) hoặc trình duyệt web độc lập (như Mozilla Firefox).
    
    \item \textbf{Tổ chức chứng thực gốc (Root CA -- Root Certificate Authority):} Là thực thể đứng ở đỉnh cao nhất của mô hình PKI phân cấp. Root CA sở hữu cặp khóa ký (signing key pair) tối cao và tự phát hành một chứng thư tự ký (Self-signed Certificate), trong đó trường định danh người phát hành (Issuer) trùng khớp hoàn toàn với trường định danh chủ thể (Subject). Khóa bí mật của Root CA ($SK_{\mathrm{Root}}$) là tài sản tối mật, thường được duy trì ở trạng thái ngoại tuyến (offline) trong các mô-đun phần cứng an toàn chuyên dụng (Hardware Security Module -- HSM) nhằm triệt tiêu nguy cơ bị xâm nhập qua mạng. Root CA chỉ được đưa vào vận hành theo quy trình kiểm toán nghiêm ngặt khi cần ký cấp phát hoặc thu hồi chứng thư cho các Intermediate CA.
    
    \item \textbf{Tổ chức chứng thực trung gian (Intermediate CA / Sub-CA):} Là các CA được Root CA (hoặc một Intermediate CA cấp cao hơn) ủy quyền và cấp chứng thư để thực hiện việc vận hành hàng ngày. Intermediate CA chịu trách nhiệm trực tiếp ký phát hành chứng thư cho các thực thể cuối (End-entity) và phát hành danh sách thu hồi (CRL). Việc thiết lập mô hình phân tầng qua Intermediate CA tạo ra một lớp bảo vệ và cô lập rủi ro (isolation layer): nếu một Intermediate CA bị thỏa hiệp hoặc lộ khóa, hệ thống chỉ cần thu hồi chứng thư của riêng Sub-CA đó mà không làm sụp đổ toàn bộ cây tin cậy của Root CA.
    
    \item \textbf{Cơ quan đăng ký (Registration Authority -- RA):} Đóng vai trò là giao diện trung gian kết nối giữa người dùng/chủ thể và CA. RA có nhiệm vụ tiếp nhận yêu cầu cấp chứng thư (CSR), thực hiện quy trình thẩm tra tính xác thực về mặt định danh, tư cách pháp nhân hoặc quyền sở hữu tên miền của đối tượng yêu cầu trước khi chuyển yêu cầu đã được phê duyệt sang cho CA ký phát hành. Bản thân RA không trực tiếp ký chứng thư và không được phép nắm giữ khóa bí mật của CA.
    
    \item \textbf{Thực thể cuối (End-entity) và Chủ thể (Subject):} Cần phân biệt rõ: \textit{Subject} là trường định danh chủ thể được đóng gói trong cấu trúc chứng thư số X.509; còn \textit{End-entity} là vai trò của thực thể ở tầng lá trong kiến trúc phân cấp PKI (như máy chủ Web TLS, thiết bị hoặc người dùng cuối). Thực thể cuối sở hữu cặp khóa mật mã và không có quyền cấp phát chứng thư cho các thực thể khác. Theo chuẩn RFC~5280~\cite{rfc5280}, phần mở rộng \texttt{basicConstraints} của chứng thư thực thể cuối có thể vắng mặt (omitted) hoặc nếu hiện diện thì cờ \texttt{cA} bắt buộc phải mang giá trị \texttt{FALSE} (\texttt{cA:FALSE}).
    
    \item \textbf{Bên tin cậy (Relying Party):} Là bất kỳ thực thể hoặc ứng dụng nào (như trình duyệt web máy trạm, ứng dụng xác thực) nhận chứng thư số từ Subject và cần thẩm định xem chứng thư đó có đáng tin cậy hay không. Relying party đối chiếu với danh mục Trust Anchor cấu hình sẵn trong Trust Store nội bộ để tiến hành xây dựng và xác minh đường dẫn chứng thư (Certification Path Validation Algorithm).
    
    \item \textbf{Kho lưu trữ và công bố chứng thư (Repository / Directory Service):} Là hệ thống máy chủ dịch vụ công cộng (thường chạy qua giao thức HTTP hoặc LDAP) chịu trách nhiệm lưu trữ và xuất bản các chứng thư CA, danh sách thu hồi chứng thư (CRL) hoặc cung cấp dịch vụ phản hồi trạng thái thu hồi trực tuyến (OCSP Responder) để relying party có thể truy vấn và tải về kiểm tra bất kỳ lúc nào.
\end{itemize}

Bảng~\ref{tab:pki-components-summary} tổng hợp vai trò, loại khóa quản lý và trạng thái hoạt động của các thành phần trong hệ thống PKI.

\begin{table}[!htbp]
\centering
\small
\renewcommand{\arraystretch}{1.25}
\begin{tabularx}{\textwidth}{>{\raggedright\arraybackslash}p{2.6cm}>{\raggedright\arraybackslash}p{3.1cm}>{\raggedright\arraybackslash}p{2.7cm}>{\raggedright\arraybackslash}X}
\toprule
\textbf{Thành phần} & \textbf{Khóa quản lý} & \textbf{Trạng thái vận hành} & \textbf{Chức năng chính} \\
\midrule
\textbf{Trust Anchor} & Khóa công khai Root CA & Tích hợp sẵn trong OS/Browser & Điểm mốc tin cậy tuyệt đối để neo chuỗi xác minh. \\
\midrule
\textbf{Root CA} & $SK_{\mathrm{Root}}$, $PK_{\mathrm{Root}}$ & Ngoại tuyến (Offline, HSM) & Ký phát hành chứng thư tự ký và ủy quyền cho Sub-CA. \\
\midrule
\textbf{Intermediate CA} & $SK_{\mathrm{Sub}}$, $PK_{\mathrm{Sub}}$ & Trực tuyến (Online) & Ký phát hành chứng thư cho End-entity, ký danh sách CRL. \\
\midrule
\textbf{RA} & Khóa xác thực danh tính & Trực tuyến & Thẩm tra định danh, duyệt yêu cầu trước khi gửi sang CA. \\
\midrule
\textbf{Chủ thể (Subject)} & $SK_{\mathrm{Entity}}$, $PK_{\mathrm{Entity}}$ & Trực tuyến (Web/Client) & Nắm giữ khóa bí mật, dùng chứng thư để chứng minh danh tính. \\
\midrule
\textbf{Bên tin cậy (RP)} & Sở hữu Trust Store & Trực tuyến (Client) & Thẩm định đường dẫn chứng thư, kiểm tra trạng thái thu hồi. \\
\midrule
\textbf{Kho lưu trữ (Repo)} & Không giữ khóa ký & Trực tuyến (HTTP, LDAP) & Lưu trữ, xuất bản chứng thư CA, CRL và dịch vụ OCSP. \\
\bottomrule
\end{tabularx}
\caption{Tổng hợp vai trò và đặc điểm vận hành của các thành phần trong PKI.}
\label{tab:pki-components-summary}
\end{table}



\subsection{Mô hình tin cậy phân cấp (Hierarchical Trust Model)}
\label{subsec:hierarchical-trust-model}

Trong hệ thống PKI chuẩn ITU-T X.509, mô hình tin cậy phổ biến và chuẩn mực nhất là \textbf{Mô hình cây phân cấp (Hierarchical Tree)}:

\begin{itemize}
    \item \textbf{Cơ chế ủy quyền niềm tin (Trust Delegation):} Niềm tin được kế thừa theo hướng một chiều từ gốc xuống ngọn ($\mathrm{Root} \rightarrow \mathrm{Intermediate} \rightarrow \mathrm{Leaf}$). Relying party chỉ cần đặt niềm tin tiên nghiệm vào Root CA nằm trong Trust Store của mình. Mọi chứng thư do Intermediate CA hợp lệ phát hành đều được tin cậy nếu toàn bộ đường dẫn chứng thư thỏa mãn các điều kiện kiểm tra xác thực ở Mục~\ref{sec:path-validation} (Chương~\ref{chap:x509-xac-thuc}).
    
    \item \textbf{Phân nhánh CA theo mục đích sử dụng (Purpose-based CA Segmentation):} Mô hình cây cho phép phân nhánh quyền hạn linh hoạt theo từng mục đích kỹ thuật chuyên biệt (ví dụ: một Intermediate CA chuyên cấp chứng thư TLS Web Server, một nhánh cấp chứng thư ký mã nguồn Code Signing, và một nhánh độc lập cấp chứng thư xác thực định danh cá nhân). Cần lưu ý phân biệt khái niệm này với nguyên tắc \textit{Segregation of Duties} trong quản trị an ninh thông tin, vốn đề cập đến việc phân tách nhiệm vụ giữa các nhân sự vận hành hệ thống nhằm ngăn ngừa xung đột lợi ích và lạm quyền.
    
    \item \textbf{Cơ chế cô lập rủi ro (Risk Isolation):} Bảo vệ Root CA ở trạng thái ngoại tuyến dài hạn; các hoạt động cấp phát tần suất cao và dễ phát sinh rủi ro trên mạng đều do Intermediate CA đảm nhiệm. Khi có sự cố lộ khóa hoặc thỏa hiệp an ninh tại một nhánh, quản trị viên chỉ cần thu hồi chứng thư của Intermediate CA tương ứng mà không làm ảnh hưởng đến các nhánh dịch vụ khác và không buộc người dùng toàn cầu phải cập nhật lại Root Trust Store.
\end{itemize}

Trong mô hình lab thực nghiệm của nhóm (Chương~\ref{chap:experiment}), các thành phần cốt lõi được triển khai gồm: Root CA, Intermediate CA (Sub-CA), End-entity (Web Server) và Relying Party (OpenSSL verification client). Các cấu phần như RA hoặc kho dịch vụ công bố trực tuyến (OCSP/LDAP) được tinh giản để tập trung nghiên cứu bản chất giải thuật xác minh chuỗi tin cậy. Đồng thời, ngoài mô hình cây phân cấp đơn gốc này, các biến thể kiến trúc khác như Chứng nhận chéo (Cross-certification), Cầu nối tin cậy (Bridge CA) và Mạng lưới tin cậy ngang hàng (Web of Trust) sẽ được phân tích đối chiếu chuyên sâu tại Mục~\ref{sec:trust-models-comparison} (Chương~\ref{chap:threat-security}).

Hình~\ref{fig:pki-architecture-components} minh họa tổng thể mối quan hệ tương tác, luồng dữ liệu cấp phát và quy trình xác thực giữa các thành phần trong mô hình PKI phân cấp.

\clearpage

\begin{figure}[!htbp]
\centering
\resizebox{0.86\textwidth}{!}{%
\begin{tikzpicture}[
    % Định dạng hộp với text width cố định
    rootbox/.style={draw=red!80!black, fill=red!8, thick, rounded corners=6pt, text width=6.6cm, align=center, inner sep=6pt, font=\small},
    cabox/.style={draw=orange!85!black, fill=orange!10, thick, rounded corners=6pt, text width=5.8cm, align=center, inner sep=6pt, font=\small},
    rabox/.style={draw=purple!80!black, fill=purple!8, thick, rounded corners=6pt, text width=3.4cm, align=center, inner sep=6pt, font=\small},
    repobox/.style={draw=teal!80!black, fill=teal!8, thick, rounded corners=6pt, text width=3.6cm, align=center, inner sep=6pt, font=\small},
    leafbox/.style={draw=blue!80!black, fill=blue!8, thick, rounded corners=6pt, text width=6.0cm, align=center, inner sep=6pt, font=\small},
    rpbox/.style={draw=green!60!black, fill=green!8, thick, rounded corners=6pt, text width=6.8cm, align=center, inner sep=6pt, font=\small},
    % Kiểu mũi tên
    arrow/.style={-{Stealth[scale=1.1]}, thick},
    dashedarrow/.style={-{Stealth[scale=1.05]}, thick, dashed},
    trustarrow/.style={-{Stealth[scale=1.1]}, very thick, blue!85!black, dashed}
]
    % --- TẦNG 1: Root CA (Đỉnh trên cùng) ---
    \node[rootbox] (root) at (9.0, 15.0) {
        \textbf{Tổ chức chứng thực gốc (Root CA)}\\[2pt]
        {\scriptsize Khóa $SK_{\mathrm{Root}}$ giữ Ngoại tuyến\\(Offline tuyệt đối trong thiết bị HSM)}
    };

    % --- TẦNG 2: Intermediate CA (Sub-CA) ---
    \node[cabox] (subca) at (9.0, 11.2) {
        \textbf{Intermediate CA (Sub-CA)}\\[2pt]
        {\scriptsize Khóa $(SK_{\mathrm{Sub}}, PK_{\mathrm{Sub}})$ Trực tuyến (Online),\\ký phát hành chứng thư máy chủ TLS}
    };

    % --- TẦNG 3: RA (Trái) và Kho lưu trữ / OCSP (Phải) ---
    \node[rabox] (ra) at (1.8, 7.0) {
        \textbf{Cơ quan đăng ký (RA)}\\[2pt]
        {\scriptsize Thẩm tra định danh,\\tiếp nhận \& duyệt CSR}
    };
    
    \node[repobox] (repo) at (16.2, 7.0) {
        \textbf{Kho lưu trữ \& OCSP}\\[2pt]
        {\scriptsize Phân phối danh sách CRL,\\phản hồi OCSP trực tuyến}
    };

    % --- TẦNG 4: Chủ thể - Web Server ---
    \node[leafbox] (leaf) at (9.0, 3.2) {
        \textbf{Chủ thể -- Web Server (Subject)}\\[2pt]
        {\scriptsize Sở hữu cặp khóa $(SK_{\mathrm{Server}}, PK_{\mathrm{Server}})$\\Thực thể cuối: End-Entity ($cA = \text{FALSE}$)}
    };

    % --- TẦNG 5: Bên tin cậy - Trình duyệt / Client ---
    \node[rpbox] (rp) at (9.0, -0.6) {
        \textbf{Bên tin cậy -- Client / Trình duyệt (Relying Party)}\\[2pt]
        {\scriptsize Sở hữu kho Root Trust Store nội bộ,\\thẩm định chuỗi chứng thư và kiểm tra trạng thái thu hồi}
    };

    % ==============================================================
    % CÁC LUỒNG TƯƠNG TÁC (Tọa độ tuyệt đối độc lập, triệt tiêu đè chữ)
    % ==============================================================

    % 1. Root CA ký cấp ủy quyền cho Sub-CA
    \draw[arrow, red!80!black] (root.south) -- (subca.north)
        node[midway, right=8pt, font=\small\bfseries] {1. Ký ủy quyền Sub-CA};

    % 2. Chủ thể nộp CSR cho RA (Tâm X=3.9 cân bằng chính xác giữa trục đứng X=1.8 và hộp Chủ thể X=6.0)
    \draw[dashedarrow, purple!80!black] (leaf.west) -| (ra.south);
    \node[purple!80!black, font=\small\bfseries, above=4pt] at (3.9, 3.2) {2. Gửi yêu cầu CSR};

    % 3. RA duyệt CSR chuyển sang Sub-CA (Tâm X=3.95 cân bằng chính xác giữa trục đứng X=1.8 và Sub-CA X=6.1)
    \draw[dashedarrow, purple!80!black] (ra.north) |- (subca.west);
    \node[purple!80!black, font=\small\bfseries, above=4pt] at (3.95, 11.2) {3. Phê duyệt CSR};

    % 4. Sub-CA ký phát hành chứng thư cho Chủ thể (Hành lang trung tâm rộng rãi)
    \draw[arrow, orange!85!black] (subca.south) -- (leaf.north)
        node[midway, right=8pt, font=\small\bfseries, align=left] {4. Ký cấp chứng thư\\TLS cho Chủ thể};

    % Xuất bản CRL: Từ Sub-CA sang đỉnh Kho lưu trữ (Căn giữa khoảng thoáng X=12.0 đến X=16.2)
    \draw[arrow, teal!80!black] (subca.east) -| (repo.north);
    \node[teal!80!black, font=\small\bfseries, above=4pt] at (14.0, 11.2) {Xuất bản CRL};

    % 5. Chủ thể trình chứng thư số trong bắt tay TLS cho Client
    \draw[arrow, blue!80!black] (leaf.south) -- (rp.north)
        node[midway, right=8pt, font=\small\bfseries] {5. Trình chứng thư (TLS Handshake)};

    % 6. Client truy vấn kiểm tra thu hồi (Cách đường đứng 10pt, hoàn toàn không chạm đường)
    \draw[dashedarrow, teal!80!black] (rp.east) -| (repo.south);
    \node[teal!80!black, font=\small\bfseries, align=right, left=10pt] at (16.2, 2.5) {6. Truy vấn\\OCSP / CRL};

    % 7. Chuỗi tin cậy: Relying Party neo niềm tin vào Root CA qua Trust Store (bọc ngoài lề trái x=-1.2)
    \draw[trustarrow, rounded corners=12pt] (rp.west) -- (-1.2, -0.6) -- (-1.2, 15.0) -- (root.west);
    \node[blue!85!black, font=\scriptsize, align=left, right=8pt] at (-1.2, 12.8) {
        \textbf{Gốc tin cậy (Trust Anchor):}\\
        Client lưu sẵn khóa Root CA\\
        trong Root Trust Store nội bộ
    };

\end{tikzpicture}%
}

\caption{Kiến trúc các thành phần và luồng dữ liệu cấp phát, xác thực trong mô hình PKI phân cấp.}
\label{fig:pki-architecture-components}
\end{figure}

\section{Vòng đời chứng thư}
\label{sec:certificate-lifecycle}
Trong khi Mục~\ref{subsec:key-lifecycle} đã trình bày về vòng đời của khóa mật mã (tập trung vào sự an toàn vật lý và toán học của khóa), Mục~\ref{sec:certificate-lifecycle} tập trung vào vòng đời của chứng thư số (quản lý trạng thái pháp lý, cấu trúc dữ liệu và chuỗi tin cậy). Hai khái niệm này có mối liên hệ mật thiết nhưng tách biệt nhau: một chứng thư có thể hết hạn hoặc bị thu hồi trước khi cặp khóa tương ứng bị tiêu hủy, và một cặp khóa có thể tiếp tục được sử dụng qua nhiều thế hệ chứng thư nếu áp dụng hình thức gia hạn giữ khóa (renewal). 

Đồng thời, bản thân các tổ chức CA cũng có vòng đời riêng: quy trình thay thế và đổi mới khóa CA (CA Key Rollover/Rekey theo mục 5.6 của dàn bài RFC~3647~\cite{rfc3647}) đòi hỏi phải tạo cặp khóa mới, phát hành chứng thư chéo hoặc phân phối lại Root Trust Store mà không làm gián đoạn các chứng thư thuê bao đang lưu hành; và quy trình ngừng hoạt động CA (CA Termination theo mục 5.8 của dàn bài RFC~3647) quy định việc thu hồi toàn bộ chứng thư đã cấp, hủy khóa CA an toàn và chuyển giao nhật ký kiểm toán cho cơ quan lưu trữ pháp lý (chi tiết xem tại Mục~\ref{sec:operating-policy}).

\subsection{Các giai đoạn và nhánh rẽ trong vòng đời chứng thư}
\label{subsec:lifecycle-stages}
Vòng đời chứng thư X.509 gồm 10 giai đoạn có các nhánh rẽ và vòng lặp thay thế nhau, được quy hoạch thành 4 pha chính như minh họa trong Hình~\ref{fig:certificate-lifecycle}.

\clearpage

\begin{figure}[!htbp]
\centering
\resizebox{0.95\textwidth}{!}{%
\begin{tikzpicture}[
    scale=0.92,
    pbox/.style={draw=blue!80!black, fill=blue!6, thick, rounded corners=6pt, text width=9.6cm, align=left, inner sep=7pt, font=\small},
    renewbox/.style={draw=teal!85!black, fill=teal!8, thick, rounded corners=6pt, text width=5.0cm, align=left, inner sep=7pt, font=\small},
    normbox/.style={draw=gray!75!black, fill=gray!8, thick, rounded corners=6pt, text width=5.2cm, align=left, inner sep=7pt, font=\small},
    revbox/.style={draw=red!80!black, fill=red!8, thick, rounded corners=6pt, text width=5.2cm, align=left, inner sep=7pt, font=\small},
    destbox/.style={draw=purple!85!black, fill=purple!8, thick, rounded corners=6pt, text width=12.0cm, align=left, inner sep=7pt, font=\small},
    arr/.style={-{Stealth[scale=1.15]}, very thick, blue!80!black},
    looparr/.style={-{Stealth[scale=1.15]}, thick, teal!85!black, dashed},
    revarr/.style={-{Stealth[scale=1.15]}, very thick, red!85!black},
    normarr/.style={-{Stealth[scale=1.15]}, very thick, gray!80!black}
]

    % ==============================================================
    % PHA 1: KHỞI TẠO & ĐĂNG KÝ
    % ==============================================================
    \node[pbox] (p1) at (5.2, 15.6) {
        \centering\textbf{\color{blue!90!black}Pha 1: Khởi tạo \& Đăng ký (Bước 1--3)}\par\vspace{3pt}
        {\footnotesize
        $\bullet$ \textbf{1. Sinh cặp khóa:} Tạo $(PK, SK)$ cục bộ qua CSPRNG ($SK$ tuyệt đối bí mật).\par
        $\bullet$ \textbf{2. Tạo CSR:} Đóng gói PKCS\#10, tự ký bởi $SK$ để chứng minh sở hữu (PoP).\par
        $\bullet$ \textbf{3. Thẩm định danh tính:} RA thẩm tra hồ sơ theo CA/B Forum (DV / OV / EV).}
    };

    % ==============================================================
    % PHA 2: CẤP PHÁT & TRIỂN KHAI
    % ==============================================================
    \node[pbox] (p2) at (5.2, 11.4) {
        \centering\textbf{\color{blue!90!black}Pha 2: Cấp phát \& Triển khai (Bước 4--5)}\par\vspace{3pt}
        {\footnotesize
        $\bullet$ \textbf{4. Cấp phát chứng thư:} CA ký \texttt{TBSCertificate} ($SK_{\text{CA}}$); nộp nhật ký CT Log.\par
        $\bullet$ \textbf{5. Triển khai cấu hình:} Cài đặt chứng thư kèm Intermediate CA Bundle lên Server.}
    };

    % ==============================================================
    % PHA 3: VẬN HÀNH & SỬ DỤNG
    % ==============================================================
    \node[pbox] (p3) at (5.2, 7.2) {
        \centering\textbf{\color{blue!90!black}Pha 3: Vận hành \& Sử dụng (Bước 6--7)}\par\vspace{3pt}
        {\footnotesize
        $\bullet$ \textbf{6. Xác minh đường dẫn:} Relying Party thẩm định chuỗi tin cậy, SAN, hạn dùng.\par
        $\bullet$ \textbf{7. Giám sát trạng thái:} Định kỳ đối chiếu hiệu lực \& kiểm tra thu hồi (CRL/OCSP).}
    };

    % ==============================================================
    % KHỐI GIA HẠN / RE-KEYING (Bên phải Pha 3)
    % ==============================================================
    \node[renewbox] (renew) at (15.0, 7.2) {
        \centering\textbf{\color{teal!90!black}Gia hạn / Re-keying (7b)}\par
        {\scriptsize\textit{\color{teal!80!black}(Thực hiện trước khi hết hạn)}}\par\vspace{3pt}
        {\scriptsize
        $\bullet$ \textbf{Renewal:} Giữ khóa cũ, gia hạn ngày.\par
        $\bullet$ \textbf{Re-keying:} Sinh cặp khóa mới hoàn toàn.\par
        $\bullet$ \textbf{Tự động hóa:} Triển khai RFC 8555 (ACME).}
    };

    % ==============================================================
    % PHA 4: KẾT THÚC VÒNG ĐỜI (HẾT HẠN HOẶC THU HỒI)
    % ==============================================================
    \node[normbox] (expire) at (1.8, 1.8) {
        \centering\textbf{\color{gray!90!black}8. Hết hạn (Expiration)}\par
        {\scriptsize\textit{\color{gray!80!black}(Nhánh kết thúc bình thường)}}\par\vspace{3pt}
        {\scriptsize
        $\bullet$ Vượt khoảng \texttt{[notBefore, notAfter]}.\par
        $\bullet$ Chữ ký số CA vẫn đúng toán học.\par
        $\bullet$ Bị từ chối theo chính sách xác minh.\par
        \textit{\color{blue!70!black}(Kiểm chứng thực nghiệm TC03)}}
    };

    \node[revbox] (revoke) at (8.6, 1.8) {
        \centering\textbf{\color{red!90!black}9. Thu hồi (Revocation)}\par
        {\scriptsize\textit{\color{red!75!black}(Nhánh kết thúc sự cố / khẩn cấp)}}\par\vspace{3pt}
        {\scriptsize
        $\bullet$ Khóa lộ (\texttt{keyCompromise}), đổi pháp nhân.\par
        $\bullet$ Đưa vào danh sách thu hồi CRL / OCSP.\par
        $\bullet$ Tồn tại độ trễ đồng bộ (Propagation Delay).\par
        \textit{\color{blue!70!black}(Kiểm chứng thực nghiệm TC07)}}
    };

    % ==============================================================
    % PHA 4B: TIÊU HỦY & LƯU TRỮ KHÓA
    % ==============================================================
    \node[destbox] (destruct) at (5.2, -2.0) {
        \centering\textbf{\color{purple!90!black}10. Tiêu hủy \& Lưu trữ khóa (Destruction \& Key Archival)}\par\vspace{3pt}
        {\scriptsize
        $\bullet$ \textbf{Khóa ký số (Digital Signature):} Tiêu hủy an toàn vĩnh viễn (Crypto-erase trên SSD/HSM) để ngăn mạo danh hồi tố.\par
        $\bullet$ \textbf{Khóa mã hóa bảo mật (Confidentiality):} Lưu trữ an toàn dài hạn (Key Archival) phục vụ giải mã tài liệu kiểm toán.}
    };

    % ==============================================================
    % CÁC ĐƯỜNG LIÊN KẾT VÀ MŨI TÊN
    % ==============================================================
    % Pha 1 -> Pha 2 -> Pha 3
    \draw[arr] (p1.south) -- (p2.north);
    \draw[arr] (p2.south) -- (p3.north);

    % Nhánh Gia hạn
    \draw[arr, teal!85!black] (p3.east) -- (renew.west);
    
    % Vòng lặp từ Gia hạn về Pha 1 & Pha 2
    \draw[looparr] (renew.east) -- ++(0.6, 0) |- (p1.east)
        node[pos=0.75, above=2pt, font=\scriptsize\bfseries, text=teal!85!black] {Re-keying (về Pha 1)};
    \draw[looparr] (renew.north) |- (p2.east)
        node[pos=0.7, above=2pt, font=\scriptsize\bfseries, text=teal!85!black] {Renewal (về Pha 2)};

    % Nhánh rẽ sang Pha 4 (Hết hạn vs Thu hồi)
    \draw[normarr] ($(p3.south)+(-2.2, 0)$) -- (expire.north)
        node[pos=0.45, left=4pt, font=\scriptsize\bfseries, text=gray!80!black] {Nhánh bình thường};
    
    \draw[revarr] ($(p3.south)+(2.2, 0)$) -- (revoke.north)
        node[pos=0.45, right=4pt, font=\scriptsize\bfseries, text=red!85!black] {Nhánh sự cố bảo mật};

    % Xuống Tiêu hủy
    \draw[normarr] (expire.south) -- ($(destruct.north)+(-3.4, 0)$);
    \draw[revarr] (revoke.south) -- ($(destruct.north)+(3.4, 0)$);

\end{tikzpicture}%
}

\caption[Các giai đoạn và nhánh rẽ trong vòng đời chứng thư X.509]{Mô hình các giai đoạn, nhánh rẽ và vòng lặp trong vòng đời chứng thư X.509.}
\label{fig:certificate-lifecycle}
\end{figure}

\vspace{0.6cm}
\subsubsection*{Pha 1: Khởi tạo và đăng ký (Bước 1--3)}
\begin{enumerate}
    \item \textbf{Sinh cặp khóa (Key Generation):} Chủ thể (Subject) tự sinh cặp khóa mật mã bất đối xứng ($PK, SK$) trên thiết bị cục bộ bằng bộ sinh số ngẫu nhiên giả an toàn mật mã (CSPRNG). Khóa bí mật ($SK$) phải được bảo vệ nghiêm ngặt, tuyệt đối không gửi qua mạng cho CA hay bất kỳ bên thứ ba nào.
    \item \textbf{Tạo yêu cầu ký chứng thư (CSR Generation):} Chủ thể đóng gói khóa công khai ($PK$) cùng thông tin định danh đề nghị (tên miền/FQDN, tổ chức) thành tệp CSR theo chuẩn RFC~2986 (PKCS\#10)~\cite{rfc2986}. Chủ thể dùng chính khóa bí mật của mình ký lên CSR để cung cấp bằng chứng sở hữu khóa (Proof of Possession -- PoP). Cần nhấn mạnh: chữ ký trên CSR chỉ chứng minh thực thể gửi yêu cầu có nắm giữ khóa bí mật tương ứng với khóa công khai, hoàn toàn không chứng minh được danh tính thực tế của người gửi.
    \item \textbf{Xác minh danh tính (Identity Verification):} Cơ quan đăng ký (RA) tiếp nhận CSR và thẩm tra thông tin định danh theo các quy định trong Tiêu chuẩn Baseline Requirements của CA/Browser Forum~\cite{cabforum_br} và chính sách CP/CPS của CA. Tùy theo loại chứng thư, quy trình thẩm định có thể là:
    \begin{itemize}
        \item \textit{Domain Validation (DV):} Xác minh quyền kiểm soát đối với tên miền thông qua bản ghi DNS hoặc phản hồi HTTP.
        \item \textit{Organization Validation (OV):} Thẩm tra sự tồn tại pháp lý của tổ chức đăng ký qua cổng thông tin đăng ký doanh nghiệp chính thức.
        \item \textit{Extended Validation (EV):} Quy trình thẩm định nghiêm ngặt bổ sung nhiều thủ tục đối chiếu giấy phép kinh doanh, quyền hạn đại diện pháp luật và xác minh độc lập bằng điện thoại theo hướng dẫn riêng EV Guidelines của CA/Browser Forum~\cite{cabforum_ev} (được chuẩn hóa tách biệt khỏi bộ tiêu chuẩn Baseline Requirements cơ bản).
    \end{itemize}
\end{enumerate}

\subsubsection*{Pha 2: Cấp phát và triển khai (Bước 4--5)}
\begin{enumerate}
    \setcounter{enumi}{3}
    \item \textbf{Cấp phát chứng thư và cơ chế minh bạch chứng thư (CT):} 
    CA không ký lên toàn bộ tập tin chứng thư hoàn chỉnh mà chỉ ký lên cấu trúc dữ liệu tóm lược \texttt{TBSCertificate} (To-Be-Signed Certificate) theo chuẩn RFC~5280~\cite{rfc5280}.
    \begin{itemize}
        \item CA trích xuất khóa công khai và thông tin danh tính đã duyệt, ấn định số sê-ri duy nhất và khoảng thời gian hiệu lực \texttt{[notBefore, notAfter]}.
        \item \textit{Bẫy cấu hình sao chép Extension (OpenSSL copy\_extensions):} CA không tự động sao chép toàn bộ các extension có trong CSR của người dùng. Trong môi trường OpenSSL, hành vi này phụ thuộc vào chỉ thị \texttt{copy\_extensions} trong tệp cấu hình CA (\texttt{openssl.cnf}). Đây là một bẫy kỹ thuật kinh điển khi thiết lập lab: nếu CA cấu hình thiếu cẩn trọng hoặc không bật \texttt{copy\_extensions = copy}, chứng thư phát hành sẽ bị thiếu mất phần mở rộng bắt buộc như Subject Alternative Name (SAN), khiến các trình duyệt hiện đại từ chối chứng thư ngay lập tức (kinh nghiệm này được đúc kết trực tiếp khi thiết kế lab và chuẩn bị các ca kiểm thử TC01, TC04, TC08 tại Chương~\ref{chap:experiment}; trong đó TC01 được cấu hình đầy đủ SAN để làm mốc đối chứng chuẩn).
        \item \textit{Quy trình nhúng SCT qua Certificate Transparency (CT):} Đối với Web PKI công cộng, theo chính sách bắt buộc của các trình duyệt web hiện đại (Google Chrome, Apple Safari) nhằm ngăn ngừa CA cấp lén lút chứng thư giả mạo: trước khi cấp chứng thư chính thức, CA phải tạo một chứng thư tiền phát hành (Precertificate) gửi lên các máy chủ nhật ký CT Log công khai (RFC~9162~\cite{rfc9162}). CT Log phản hồi bằng bằng chứng ghi nhận thời gian SCT (Signed Certificate Timestamp). CA nhúng các SCT này vào phần mở rộng của chứng thư chính thức rồi mới ký phát hành ($SK_{\text{CA}}$). Cần lưu ý rằng việc bắt buộc nộp CT Log là quy định trong chính sách Root Store của trình duyệt, không nằm trong chuẩn gốc RFC 5280 (chi tiết xem tại Mục~\ref{sec:certificate-transparency}).
    \end{itemize}
    \item \textbf{Cài đặt và triển khai (Installation / Deployment):} Chủ thể cấu hình chứng thư số cùng với tệp chuỗi chứng thư trung gian (Intermediate CA certificates / Certificate Bundle) lên máy chủ dịch vụ. Nếu thiếu chứng thư trung gian trên máy chủ, client không thể hoàn thành việc xây dựng chuỗi tin cậy dẫn tới Trust Anchor (liên hệ tới ca kiểm thử TC06 tại Chương~\ref{chap:experiment}).
\end{enumerate}

\subsubsection*{Pha 3: Vận hành (Bước 6--7)}
\begin{enumerate}
    \setcounter{enumi}{5}
    \item \textbf{Xác minh và sử dụng (Validation \& Usage):} Khi máy chủ gửi chứng thư cho Relying Party (ví dụ trong pha bắt tay TLS~1.3~\cite{rfc8446}), bên tin cậy sẽ kích hoạt thuật toán xác minh đường dẫn đầy đủ. Quy trình bao gồm: xây dựng chuỗi dẫn tới Trust Anchor đã cấu hình, kiểm tra chữ ký từng mắt xích, thẩm định tính hợp lệ của Basic Constraints và độ sâu chuỗi (\texttt{pathLenConstraint}), đối chiếu mục đích sử dụng khóa (Key Usage/Extended Key Usage), khớp tên máy chủ (SAN), và kiểm tra thời hạn cùng trạng thái thu hồi. Chi tiết toàn bộ thuật toán được phân tích chuyên sâu tại Mục~\ref{sec:path-validation} (Chương~\ref{chap:x509-xac-thuc}).
    \item \textbf{Gia hạn (Renewal / Re-keying):} Trước khi chứng thư mãn hạn, chủ thể thực hiện gia hạn để duy trì kết nối liên tục:
    \begin{itemize}
        \item Nếu thực hiện \textit{Renewal}, chủ thể tái sử dụng cặp khóa cũ và chỉ xin cấp chứng thư mới với thời hạn sử dụng mới.
        \item Nếu thực hiện \textit{Re-keying}, chủ thể quay trở lại Bước~1 để tạo cặp khóa hoàn toàn mới.
    \end{itemize}
    \textit{Lộ trình rút ngắn thời hạn chứng thư của CA/Browser Forum:} Theo phiên bản mới nhất của Tiêu chuẩn Baseline Requirements (v2.3.1)~\cite{cabforum_br}, CA/Browser Forum đã thiết lập lộ trình liên tục rút ngắn thời hạn hiệu lực tối đa đối với chứng thư số TLS máy chủ công cộng:
    \begin{itemize}
        \item Tối đa \textbf{200 ngày}: áp dụng từ ngày 15/03/2026.
        \item Tối đa \textbf{100 ngày}: áp dụng từ ngày 15/03/2027.
        \item Tối đa \textbf{47 ngày}: áp dụng từ ngày 15/03/2029.
    \end{itemize}
    Xu hướng rút ngắn mạnh mẽ này nhằm hạn chế thời gian tổn thương khi khóa bị lộ, nhưng đồng thời kéo theo sự bắt buộc phải ứng dụng các giao thức tự động hóa quản lý vòng đời chứng thư như ACME (RFC~8555)~\cite{rfc8555} (phân tích chi tiết tại Mục~\ref{sec:acme-automation}).
\end{enumerate}

\subsubsection*{Pha 4: Kết thúc vòng đời (Bước 8--10: Nhánh rẽ thay thế và Tiêu hủy)}
\begin{enumerate}
    \setcounter{enumi}{7}
    \item \textbf{Hết hạn (Expiration):} Chứng thư không còn vượt qua bước kiểm tra thời hạn hiệu lực nếu thời điểm xác minh nằm ngoài khoảng \texttt{[notBefore, notAfter]} quy định trong cấu trúc chứng thư.
    Cần lưu ý rằng: chữ ký số của CA trên một chứng thư đã hết hạn vẫn hoàn toàn đúng về mặt toán học. Việc từ chối chứng thư phụ thuộc chặt chẽ vào chính sách xác minh và tính chính xác của đồng hồ hệ thống phía relying party; một client bị sai lệch đồng hồ hoặc cấu hình lỏng lẻo bỏ qua cờ kiểm tra thời gian vẫn có thể chấp nhận chứng thư này. Luận điểm này minh chứng thực nghiệm cho giả thuyết H1 và được kiểm chứng qua ca kiểm thử TC03 (Mục~\ref{sec:exp-results}).
    \item \textbf{Thu hồi chứng thư (Revocation):} Là nhánh rẽ kết thúc khẩn cấp thay thế cho việc hết hạn tự nhiên. Khi phát sinh sự cố, CA sẽ đưa số sê-ri của chứng thư vào danh sách thu hồi kèm theo mã lý do quy chuẩn theo RFC~5280~\cite{rfc5280}, tiêu biểu như:
    \begin{itemize}
        \item \texttt{keyCompromise} (khóa bí mật bị lộ/thỏa hiệp),
        \item \texttt{cACompromise} (khóa của CA bị xâm nhập),
        \item \texttt{affiliationChanged} (chủ thể thay đổi quan hệ công tác/pháp nhân),
        \item \texttt{superseded} (chứng thư đã được thay thế bằng chứng thư mới),
        \item \texttt{cessationOfOperation} (dịch vụ ngừng hoạt động hoàn toàn).
    \end{itemize}
    \textit{Độ trễ thu hồi (Propagation Delay):} Việc thu hồi không có hiệu lực tức thời đối với mọi client. Relying party chỉ biết chứng thư bị thu hồi khi danh sách CRL mới được định kỳ phát hành và tải về thành công (phụ thuộc vào các trường \texttt{thisUpdate} và \texttt{nextUpdate}), hoặc thông qua truy vấn OCSP trực tiếp (phân tích sâu tại Mục~\ref{sec:revocation-mechanisms} và kiểm chứng qua ca kiểm thử TC07).
    \item \textbf{Tiêu hủy và lưu trữ khóa (Destruction \& Key Archival):}
    \begin{itemize}
        \item \textit{Khóa ký số:} Khi chứng thư hết hạn hoặc bị thu hồi và không còn giá trị sử dụng, khóa bí mật dùng cho chữ ký số phải bị tiêu hủy vĩnh viễn. Trên các ổ cứng thể rắn (SSD) hay bộ nhớ flash hiện đại, các lệnh xóa file thông thường của hệ điều hành không đảm bảo dữ liệu vật lý bị triệt tiêu do cơ chế Wear Leveling; vì vậy cần áp dụng các cơ chế xóa an toàn bằng mật mã (Crypto-erase) hoặc phá hủy trên các phân vùng bảo mật của phần cứng chuyên dụng (HSM/TPM).
        \item \textit{Khóa mã hóa dữ liệu:} Ngược lại với khóa ký số, nếu một cặp khóa từng được sử dụng để mã hóa lưu trữ tài liệu dữ liệu lâu dài (Confidentiality key), khóa bí mật giải mã có thể phải được đưa vào quy trình lưu trữ an toàn dài hạn (Key Archival) dưới sự bảo vệ của chính sách sao lưu nghiêm ngặt, nhằm phục vụ công tác giải mã kiểm toán trong tương lai ngay cả khi chứng thư liên kết đã hết hiệu lực.
    \end{itemize}
\end{enumerate}

\subsection{Ánh xạ giai đoạn vòng đời chứng thư với công cụ OpenSSL và kịch bản thực nghiệm}
Nhằm cung cấp cái nhìn xuyên suốt từ quy trình chuẩn hóa đến hoạt động kỹ thuật cụ thể trong phòng lab, Bảng~\ref{tab:lifecycle-openssl-tc} tổng hợp mối liên kết giữa các giai đoạn trong vòng đời chứng thư X.509, các câu lệnh OpenSSL tương ứng và các ca kiểm thử thực nghiệm (TC01--TC10) được triển khai tại Chương~\ref{chap:experiment}.

\begin{table}[!htbp]
\centering
\small
\renewcommand{\arraystretch}{1.25}
\begin{tabularx}{\textwidth}{>{\raggedright\arraybackslash}p{2.6cm}>{\raggedright\arraybackslash}p{5.5cm}>{\raggedright\arraybackslash}X}
\toprule
\textbf{Giai đoạn vòng đời} & \textbf{Lệnh OpenSSL thực thi trong Lab} & \textbf{Ca kiểm thử liên quan} \\
\midrule
\textbf{1. Sinh cặp khóa} & \texttt{openssl genrsa -aes256 ...} hoặc \texttt{openssl ecparam ...} & Nền tảng mọi kịch bản \\
\midrule
\textbf{2. Tạo CSR (PoP)} & \texttt{openssl req -new -key ... -out ...} & Tạo yêu cầu ký chứng thư kèm PoP \\
\midrule
\textbf{4. Cấp phát chứng thư} & \texttt{openssl ca -in ... -out ... -extfile ...} & Cấp phát X.509 với extension; TC05 (chữ ký bị sửa đổi bit), TC08 (cấp sai cờ EKU) \\
\midrule
\textbf{5. Triển khai cấu hình} & Cấu hình tệp chứng thư và chuỗi Intermediate CA bundle & TC06 (thiếu Intermediate CA bundle) \\
\midrule
\textbf{6. Xác minh đường dẫn} & \texttt{openssl verify -CAfile ... -untrusted ... -purpose ... -verify\_hostname ...} & TC01 (thẩm định thành công toàn diện), TC02 (self-signed không thuộc trust store), TC04 (hostname không khớp SAN), TC08 (sai mục đích EKU), TC10 (CA ngoài trust store) \\
\midrule
\textbf{8. Kiểm tra thời hạn} & \texttt{openssl verify -attime ...} (hoặc so khớp thời gian thực) & TC03 (thời điểm kiểm tra nằm ngoài hạn hiệu lực) \\
\midrule
\textbf{9. Thu hồi chứng thư} & \texttt{openssl ca -revoke ...}, \texttt{openssl ca -gencrl ...} và xác minh với \texttt{-crl\_check} & TC07 (chứng thư bị thu hồi đưa vào CRL) \\
\bottomrule
\end{tabularx}
\caption{Ánh xạ các giai đoạn vòng đời chứng thư X.509 với lệnh OpenSSL và ca kiểm thử.}
\label{tab:lifecycle-openssl-tc}
\end{table}



\clearpage

\subsection{Phân biệt Yêu cầu ký chứng thư (CSR) và Chứng thư số đã phát hành}
\label{subsec:csr-vs-cert}
Để tránh nhầm lẫn giữa định dạng yêu cầu đăng ký và đối tượng chứng thực, Bảng~\ref{tab:csr-vs-x509} đối chiếu bản chất kỹ thuật giữa tệp CSR (RFC~2986~\cite{rfc2986}) và chứng thư số X.509 v3 (RFC~5280~\cite{rfc5280}):

\begin{table}[!htbp]
\centering
\small
\renewcommand{\arraystretch}{1.25}
\begin{tabularx}{\textwidth}{>{\raggedright\arraybackslash}p{3.2cm}>{\raggedright\arraybackslash}X>{\raggedright\arraybackslash}X}
\toprule
\textbf{Tiêu chí phân định} & \textbf{Yêu cầu ký chứng thư (CSR -- PKCS\#10)} & \textbf{Chứng thư số X.509 v3} \\
\midrule
\textbf{Bản chất kỹ thuật} & Cấu trúc dữ liệu chứa thông tin đề nghị và khóa công khai do chủ thể tự tạo ra. & Cấu trúc dữ liệu có chữ ký số của CA, liên kết định danh với khóa công khai theo quy trình CP/CPS. \\
\midrule
\textbf{Thực thể thực hiện ký} & Ký bởi chính khóa bí mật của chủ thể ($SK_{\text{subject}}$). & Ký bởi khóa bí mật của tổ chức cấp phát ($SK_{\text{CA}}$). \\
\midrule
\textbf{Dữ liệu được ký} & Thông tin danh tính yêu cầu và trường khóa công khai (\texttt{SubjectPublicKeyInfo}). & Cấu trúc dữ liệu chuẩn hóa \texttt{TBSCertificate} (To-Be-Signed Certificate) của CA. \\
\midrule
\textbf{Ý nghĩa của chữ ký} & Chỉ chứng minh thực thể gửi yêu cầu có nắm giữ khóa bí mật tương ứng (Proof of Possession -- PoP); không chứng minh danh tính. & Xác nhận CA đã thẩm định thông tin theo chính sách vận hành và liên kết danh tính với khóa công khai. \\
\midrule
\textbf{Quyết định về Extensions} & Chỉ mang tính chất đề xuất từ phía client. & Do CA toàn quyền ấn định hoặc chọn lọc (qua cờ cấu hình như \texttt{copy\_extensions} trong OpenSSL). \\
\midrule
\textbf{Phạm vi sử dụng} & Chỉ dùng để nộp cho CA/RA trong pha đăng ký cấp phát. & Dùng trong giao thức truyền thông (TLS) để relying party xác thực chuỗi tin cậy. \\
\bottomrule
\end{tabularx}
\caption[Phân biệt bản chất kỹ thuật giữa CSR và chứng thư số X.509]{Phân biệt bản chất kỹ thuật giữa Yêu cầu ký chứng thư (CSR -- RFC~2986) và Chứng thư số X.509 v3 (RFC~5280).}
\label{tab:csr-vs-x509}
\end{table}



\clearpage

\section{Chính sách vận hành}
\label{sec:operating-policy}

Một hạ tầng khóa công khai (PKI) không chỉ bao gồm các thuật toán mật mã và giao thức kỹ thuật, mà còn phụ thuộc cốt lõi vào hệ thống chính sách và quy trình vận hành bảo đảm mức độ tin cậy. Về mặt cấu trúc tài liệu, RFC~3647~\cite{rfc3647} cung cấp một khung dàn bài chuẩn hóa gồm 9 phần lớn, liệt kê danh mục các chủ đề kỹ thuật, an ninh và quản trị mà một CA cần đặc tả. Cần lưu ý rằng các ký hiệu số mục được dẫn chiếu trong Mục~\ref{sec:operating-policy} này (như mục 3, 5.2, 5.4, 6.2.2, 8, 9) là số mục trong dàn bài chuẩn của tài liệu CP/CPS theo RFC~3647, không phải số mục điều khoản của chính bản thân văn bản RFC~3647. Bản thân RFC~3647 chỉ đóng vai trò phân loại và định hướng cấu trúc nội dung; các tiêu chí kỹ thuật cụ thể hay mức độ ràng buộc thực tế sẽ do CA tự xác định hoặc phải tuân thủ theo tiêu chuẩn của các bên quản lý kho khóa gốc (Root Store Policies) và các tổ chức quản lý ngành như CA/Browser Forum~\cite{cabforum_br}.

\subsection{Phân định vai trò giữa CP và CPS}
\label{subsec:cp-vs-cps}

Khung dàn bài RFC~3647 phân chia hai tầng tài liệu chính sách:
\begin{itemize}
    \item \textbf{Chính sách Chứng thư (Certificate Policy -- CP):} Xác định các quy tắc về khả năng áp dụng của một nhóm chứng thư cho một cộng đồng người dùng cụ thể hoặc lớp ứng dụng có chung yêu cầu an ninh. CP mang tính định hướng cấp cao, làm rõ mục đích sử dụng và mức độ bảo đảm mà CA cam kết hướng tới bên tin cậy (Relying Party). Bản thân CP không tự là hợp đồng pháp lý hoàn chỉnh; hiệu lực ràng buộc thực tế thường được thiết lập qua thỏa thuận thuê bao (Subscriber Agreement) hoặc thỏa thuận bên tin cậy (Relying Party Agreement) có dẫn chiếu đến CP đó.
    
    \item \textbf{Tuyên bố Thực hành Chứng thực (Certification Practice Statement -- CPS):} Đặc tả chi tiết các biện pháp kỹ thuật, quy trình thủ tục và cơ chế kiểm soát cụ thể mà một CA áp dụng để hiện thực hóa các yêu cầu được nêu trong một hoặc nhiều CP tương ứng.
\end{itemize}

\textbf{Mối liên kết kỹ thuật và thực tiễn triển khai:} Về mặt kỹ thuật theo chuẩn RFC~5280~\cite{rfc5280}, CP được định danh trên chứng thư số thông qua một mã định danh đối tượng (OID) đặt trong phần mở rộng \texttt{certificatePolicies}, có thể đi kèm một đường dẫn liên kết tĩnh (CPS URI) trỏ trực tiếp đến tài liệu thực hành của CA. Trên thực tế, Mục 2.2 của Tiêu chuẩn Baseline Requirements (CA/Browser Forum)~\cite{cabforum_br} yêu cầu rằng tài liệu CP và/hoặc CPS của CA công cộng phải được cấu trúc bám sát dàn bài chuẩn của RFC~3647 và chứa đầy đủ nội dung theo quy định. Bản thân Mục 2.2 không bắt buộc phải gộp CP và CPS thành một văn bản; việc xây dựng tài liệu gộp chung mang tên Certificate Policy and Certification Practice Statement (Combined CP/CPS) là giải pháp thực tế phổ biến được nhiều CA thương mại lựa chọn nhằm tinh giản quản trị và bảo đảm tính đồng bộ.

\textbf{Lưu ý tránh nhầm lẫn trong thực nghiệm:} Trong tệp cấu hình của công cụ OpenSSL (\texttt{openssl.cnf}), chỉ thị \texttt{policy = policy\_strict} (hoặc \texttt{policy\_loose}) chỉ đơn thuần là quy tắc kỹ thuật cấp thấp nhằm đối chiếu sự trùng khớp của các trường tên phân biệt (Distinguished Name -- DN) giữa yêu cầu cấp phát (CSR) và chứng thư CA. Chỉ thị này hoàn toàn không phải là một Certificate Policy (CP) theo định nghĩa của chuẩn RFC~3647.

\subsection{Quản lý và bảo vệ khóa bí mật của CA (Mục 6 của dàn bài RFC 3647)}
\label{subsec:ca-key-protection}

Mục 6 (Technical Security Controls) trong dàn bài RFC~3647 liệt kê các chủ đề kỹ thuật mà CPS cần làm rõ, bao gồm việc tạo sinh cặp khóa, bảo vệ khóa bí mật của CA và quản lý vòng đời khóa:
\begin{itemize}
    \item \textbf{Bảo vệ khóa bằng phần cứng:} Trong thực hành của các CA công cộng theo Mục 6.2.7 của chuẩn Baseline Requirements~\cite{cabforum_br}, khóa bí mật của Root CA và Intermediate CA bắt buộc phải được tạo sinh và lưu trữ trong các mô-đun phần cứng an toàn chuyên dụng (Hardware Security Module -- HSM) đạt tiêu chuẩn tối thiểu \textbf{FIPS 140-2 Level 3} hoặc \textbf{FIPS 140-3 Level 3}~\cite{fips140_3}, hoặc đạt chứng chỉ \textbf{Common Criteria EAL 4+} trở lên. Tiêu chuẩn FIPS 140 Level 2 chỉ được chấp nhận đối với khóa của thuê bao (subscriber) ký mã nguồn (code signing), hoàn toàn không đủ điều kiện cho khóa của các tổ chức CA phát hành chứng thư công cộng.
    
    \item \textbf{Kiểm soát nhiều người (Multiperson Control -- Mục con 6.2.2 của dàn bài RFC 3647):} Nhằm ngăn chặn hành vi lạm quyền của một cá nhân đơn lẻ, các CA thường áp dụng cơ chế phân mảnh bí mật với ngưỡng $m$-trên-$n$ (như thuật toán Shamir's Secret Sharing~\cite{shamir1979}, ví dụ yêu cầu kích hoạt đồng thời của 3 trên 5 thẻ khóa vật lý) đối với các thao tác nhạy cảm như kích hoạt khóa CA, ký phát chứng thư trung gian hoặc khôi phục sao lưu.
    
    \item \textbf{Nghi thức tạo khóa (Key Generation Ceremony):} Khóa bí mật của CA cấp cao thường được khởi tạo theo một nghi thức chính thức diễn ra trong môi trường an toàn vật lý được cách ly mạng, có nhân sự giám sát độc lập và được lập biên bản ghi nhận đầy đủ nhằm phục vụ công tác kiểm toán.
\end{itemize}

\subsection{Phân quyền và các biện pháp kiểm soát thủ tục (Mục 5.2 của dàn bài RFC 3647)}
\label{subsec:procedural-controls}

Mục 5.2 trong dàn bài RFC~3647 có tiêu đề là Procedural Controls (các biện pháp kiểm soát thủ tục), liệt kê các yêu cầu mà CPS cần đặc tả về vai trò tin cậy (Trusted Roles), số lượng nhân sự cho mỗi tác vụ và nguyên tắc tách biệt nhiệm vụ:
\begin{itemize}
    \item \textbf{Các vai trò tin cậy phổ biến:} Khung RFC~3647 không ấn định sẵn tên gọi chức danh, nhưng thực tiễn công nghiệp thường phân tách thành các vai trò độc lập điển hình như:
    \begin{itemize}
        \item \textit{Quản trị viên (Administrator):} Cài đặt, bảo trì hạ tầng mạng, máy chủ và phần mềm nền tảng.
        \item \textit{Nhân viên vận hành CA (CA Operator):} Tiếp nhận hồ sơ xác minh, thực hiện các thao tác cấp phát và thu hồi chứng thư số.
        \item \textit{Cán bộ kiểm toán (Audit Officer):} Độc lập hoàn toàn với nhóm kỹ thuật/vận hành, chịu trách nhiệm lưu trữ và giám sát nhật ký an ninh.
        \item \textit{Cán bộ an ninh (Security Officer):} Quản lý kiểm soát an ninh vật lý và cấp phát quyền truy cập thiết bị HSM.
    \end{itemize}
    
    \item \textbf{Tách biệt nhiệm vụ và kiểm soát kép:} CPS cần làm rõ nguyên tắc một nhân sự không được phép kiêm nhiệm các vai trò mâu thuẫn lợi ích (đặc biệt là giữa người vận hành hệ thống và người giám sát nhật ký). Những thao tác trọng yếu thường yêu cầu cơ chế kiểm soát kép (Dual Control), đòi hỏi sự hiện diện và đồng thuận của tối thiểu hai cá nhân có thẩm quyền.
\end{itemize}

\subsection{Ghi nhận nhật ký và đánh giá tuân thủ (Mục 5.4 và Mục 8 của dàn bài RFC 3647)}
\label{subsec:audit-log-compliance}

\begin{itemize}
    \item \textbf{Nhật ký kiểm toán (Mục 5.4 của dàn bài RFC 3647):} CPS cần xác định danh mục các sự kiện phải ghi log (khởi động/dừng dịch vụ, truy cập quản trị, tạo khóa, phát hành chứng thư, tạo CRL) kèm nhãn thời gian tin cậy. Dữ liệu nhật ký được khuyến nghị sao lưu định kỳ vào các phân vùng lưu trữ khó sửa đổi hoặc hệ thống thu thập log tập trung được phân quyền độc lập.
    
    \item \textbf{Đánh giá tuân thủ (Mục 8 của dàn bài RFC 3647):} RFC~3647 liệt kê các chủ đề về tần suất đánh giá, phạm vi kiểm toán và tư cách của tổ chức đánh giá. Theo chính sách quản lý Root Store và quy định của CA/Browser Forum~\cite{cabforum_br}, các CA công cộng thực tế phải trải qua các đợt kiểm toán định kỳ hằng năm bởi các tổ chức kiểm toán độc lập đủ điều kiện theo các chuẩn mực như WebTrust for CA~\cite{webtrust} hoặc ETSI EN 319 411~\cite{etsi_en_319_411}.
    
    \item \textbf{Bài học từ sự cố DigiNotar (Fox-IT, 2011)~\cite{prins2011diginotar}:} Báo cáo điều tra pháp y độc lập của Fox-IT~\cite{prins2011diginotar} về vụ xâm nhập nhà cấp phát chứng thư số DigiNotar (Hà Lan) làm sáng tỏ những lỗ hổng nghiêm trọng trong quản trị hạ tầng CA: Mặc dù DigiNotar có trang bị thiết bị HSM (nCipher netHSM) và khóa CA chỉ được kích hoạt thông qua thẻ thông minh kèm mã PIN, toàn bộ 8 máy chủ quản lý CA đều bị tin tặc chiếm quyền điều khiển, cho phép kẻ tấn công thực hiện ký phát hàng loạt chứng thư giả mạo ngay trong lúc khóa đang ở trạng thái kích hoạt (activated). Thêm vào đó, tất cả máy chủ CA đều gia nhập cùng một miền Windows (Active Directory) với chính sách mật khẩu yếu dễ bị vét cạn, hệ thống không cài đặt phần mềm chống mã độc, nhật ký an ninh lưu trữ cục bộ trên chính máy bị chiếm quyền và đã bị can thiệp xóa dấu vết. Cấu trúc mạng nội bộ dù chia 24 phân đoạn và có vùng Secure-net nhưng các chính sách tường lửa lỏng lẻo với nhiều luật ngoại lệ. Sự cố DigiNotar là minh chứng lịch sử đắt giá cho thấy: phần cứng HSM, văn bản chính sách CP/CPS và kiểm toán định kỳ chỉ là điều kiện cần; an ninh thực sự của hệ thống PKI phụ thuộc cốt lõi vào sự nghiêm ngặt trong triển khai phân tách mạng, quản trị danh tính và giám sát kỹ thuật liên tục.
\end{itemize}

\subsection{Đối chiếu chính sách CA thực tế và môi trường thực nghiệm của nhóm}
\label{subsec:real-vs-lab-ca}

Chính sách vận hành hoàn chỉnh của một CA trong môi trường sản xuất không thể được suy diễn hay đại diện chỉ thông qua các kịch bản thực nghiệm trong phòng lab. Cần nhấn mạnh rằng, theo điều khoản phạm vi áp dụng (Scope) trong Tiêu chuẩn Baseline Requirements của CA/Browser Forum~\cite{cabforum_br}, bộ tiêu chuẩn BR không điều chỉnh các hệ thống PKI nội bộ doanh nghiệp mà Root CA không được các nhà phát triển phần mềm/hệ điều hành (như Microsoft, Apple, Google, Mozilla) phân phối mặc định trong Root Trust Store; các yêu cầu kỹ thuật và thủ tục nghiêm ngặt của BR chỉ mang tính bắt buộc khi được các chương trình Root Store chấp nhận và thực thi. Luận điểm này là cơ sở giải thích vì sao môi trường lab thực nghiệm của nhóm nằm ngoài phạm vi điều chỉnh bắt buộc của BR và được thiết kế tinh giản phục vụ học thuật.

Bảng~\ref{tab:rfc3647-vs-lab} tổng hợp sự khác biệt cơ bản giữa yêu cầu của một CA công cộng thực tế theo khung RFC~3647 và môi trường thử nghiệm đã được triển khai:

\begin{table}[!htbp]
\centering
\small
\renewcommand{\arraystretch}{1.25}
\begin{tabularx}{\textwidth}{>{\raggedright\arraybackslash}p{3.2cm}>{\raggedright\arraybackslash}X>{\raggedright\arraybackslash}X}
\toprule
\textbf{Chủ đề theo dàn bài RFC~3647} & \textbf{CPS của CA công cộng thường đặc tả} & \textbf{Triển khai trong bài lab thực nghiệm của nhóm} \\
\midrule
\textbf{Xác minh danh tính (Mục~3)} & Thường đặc tả quy trình thẩm định hồ sơ pháp lý theo phân hạng chứng thư (DV, OV, EV), đối chiếu cơ sở dữ liệu doanh nghiệp và xác minh quyền kiểm soát tên miền. & Bỏ qua hoàn toàn quy trình xác minh thực tế; thông tin chủ thể (Subject DN) được tự gán trực tiếp qua tệp cấu hình hoặc câu lệnh OpenSSL. \\
\midrule
\textbf{Kiểm soát thủ tục \& vai trò tin cậy (Mục~5.2)} & Thường phân tách các vai trò độc lập (Quản trị viên, Vận hành, Kiểm toán, An ninh); áp dụng cơ chế kiểm soát kép (Dual Control) cho tác vụ nhạy cảm. & Toàn bộ thao tác do nhóm nghiên cứu thực hiện trên một máy trạm, một tài khoản hệ điều hành, không phân tách vai trò. \\
\midrule
\textbf{Nhật ký kiểm toán (Mục~5.4)} & Ghi vết tự động các sự kiện an ninh với nhãn thời gian tin cậy; chuyển log về máy chủ tập trung an toàn hoặc phân vùng lưu trữ bất biến. & Chỉ dựa vào tệp văn bản thô \texttt{index.txt} và \texttt{serial} của OpenSSL; không có cơ chế lưu vết truy cập quản trị hay chống chỉnh sửa log. \\
\midrule
\textbf{Bảo vệ khóa bí mật của CA (Mục~6.2, 6.2.2)} & Lưu trữ bắt buộc trong thiết bị HSM đạt chuẩn tối thiểu FIPS 140-2/140-3 Level 3; kích hoạt bằng cơ chế chia sẻ bí mật ngưỡng nhiều người ($m$-trên-$n$). & Khóa bí mật của Root CA và Intermediate CA lưu trực tiếp dưới dạng tệp tin phần mềm (\texttt{.key}) không đặt mật khẩu (\texttt{-nodes}) trên ổ đĩa máy tính, không sử dụng thiết bị phần cứng bảo mật HSM. \\
\midrule
\textbf{Đánh giá tuân thủ (Mục~8)} & Định kỳ thực hiện kiểm toán độc lập hằng năm theo các chuẩn mực quốc tế (WebTrust, ETSI) để duy trì sự tin cậy trên Root Store. & Không có quy trình đánh giá tuân thủ; là hệ thống tự tạo và tự tin cậy cục bộ phục vụ học thuật. \\
\midrule
\textbf{Pháp lý và trách nhiệm (Mục~9)} & Quy định rõ về quỹ bảo hiểm rủi ro, cam kết tài chính với bên tin cậy và cơ chế tài phán giải quyết khiếu nại. & Hoàn toàn không phát sinh cam kết pháp lý hay ràng buộc tài chính. \\
\bottomrule
\end{tabularx}
\caption[So sánh đối chiếu giữa khung RFC 3647, CA thực tế và bài lab]{So sánh đối chiếu giữa khung dàn bài RFC~3647, triển khai CA thực tế và bài lab của nhóm.}
\label{tab:rfc3647-vs-lab}
\end{table}



Sự đối chiếu trên làm rõ rằng các ca kiểm thử từ TC01 đến TC10 được thực hiện trong báo cáo chỉ nhằm mục đích quan sát phản ứng của quy trình xác minh chuỗi chứng thư đối với từng điều kiện vi phạm kỹ thuật cụ thể theo chuẩn RFC~5280~\cite{rfc5280}. Các kết quả này hoàn toàn không phản ánh mức độ an toàn vận hành, năng lực phòng vệ hay tính tuân thủ quy trình an ninh mạng mà một CA thực tế thường phải thiết lập.

\chapter{X.509, chuỗi tin cậy và cơ chế xác thực}
\label{chap:x509-xac-thuc}
\section{Cấu trúc chứng thư X.509 v3}
\label{sec:x509-structure}

Chuẩn ITU-T X.509 phiên bản 3 (RFC~5280~\cite{rfc5280}) định nghĩa định dạng chứng thư số bằng cú pháp ASN.1 (Abstract Syntax Notation One). Cấu trúc tổng thể của chứng thư X.509 v3 gồm ba trường chính ở cấp cao nhất:
\begin{verbatim}
Certificate  ::=  SEQUENCE  {
    tbsCertificate       TBSCertificate,
    signatureAlgorithm   AlgorithmIdentifier,
    signatureValue       BIT STRING  }
\end{verbatim}

Trong đó:
\begin{itemize}
    \item \texttt{tbsCertificate} (To-Be-Signed Certificate): Khối dữ liệu chứa toàn bộ thông tin định danh và tham số kỹ thuật cần được ký bảo đảm.
    \item \texttt{signatureAlgorithm}: Định danh thuật toán mật mã mà CA sử dụng để tạo chữ ký số (ví dụ: \texttt{sha256WithRSAEncryption} hoặc \texttt{ecdsa-with-SHA256}).
    \item \texttt{signatureValue}: Giá trị chữ ký số thực tế (chuỗi bit \texttt{BIT STRING}) do CA tính toán trên cấu trúc \texttt{tbsCertificate} mã hóa theo quy tắc chuẩn hóa DER. Việc sửa đổi dù chỉ một bit trên chứng thư đã phát hành sẽ phá hủy tính hợp lệ toán học của chữ ký và bị từ chối ngay lập tức (phép thử tính toàn vẹn này được kiểm chứng tại ca kiểm thử TC05).
\end{itemize}

Khối \texttt{tbsCertificate} là đối tượng cốt lõi phản ánh toàn bộ nội dung quản trị và kỹ thuật của chứng thư. Theo Mục~4.1 và 4.2 của RFC~5280~\cite{rfc5280}, các trường chuẩn và phần mở rộng thiết yếu bao gồm:

\subsection{Các trường dữ liệu chuẩn (Standard Fields)}
\label{subsec:standard-fields}

\begin{itemize}
    \item \textbf{Số sê-ri (Serial Number):} Số nguyên dương duy nhất do CA gán cho mỗi chứng thư ban hành, độ dài không quá 20 octet theo RFC~5280 (Mục 4.1.2.2). Để chống tấn công va chạm (collision) trên hàm băm, theo BR~\cite{cabforum_br}, CA công cộng phải sinh serial có ít nhất 64 bit ngẫu nhiên thay vì sử dụng cơ chế tăng tuần tự đơn giản.

    \item \textbf{Tên đơn vị cấp phát (Issuer):} Định danh CA đã ký và phát hành chứng thư dưới dạng Tên phân biệt (Distinguished Name -- DN) theo chuẩn X.500, gồm chuỗi các thuộc tính như Quốc gia ($C$), Tổ chức ($O$), Tên phổ biến ($CN$). Đối với chứng thư tự ký của Root CA, trường \texttt{issuer} trùng khớp hoàn toàn với trường \texttt{subject}.

    \item \textbf{Thời hạn hiệu lực (Validity):} Xác định khoảng thời gian mà CA cam kết duy trì thông tin trạng thái chứng thư, xác định bởi cặp mốc \texttt{notBefore} và \texttt{notAfter} (dùng \texttt{UTCTime} trước năm 2050 và \texttt{GeneralizedTime} từ năm 2050 trở đi). Thuật toán xác minh yêu cầu kiểm tra thời điểm xác minh nằm trong khoảng này; kết quả phụ thuộc chính sách của client (kiểm chứng tại ca kiểm thử TC03).

    \item \textbf{Tên chủ thể (Subject):} Định danh thực thể sở hữu khóa công khai bên trong chứng thư dưới dạng DN. RFC~5280 cho phép Subject rỗng khi danh tính nằm hoàn toàn trong SAN (khi đó SAN phải critical); chứng thư TLS hiện đại định danh tên máy chủ chủ yếu qua SAN.

    \item \textbf{Thông tin khóa công khai của chủ thể (Subject Public Key Info):} Chứa khóa công khai thực tế của chủ thể cùng mã OID định danh thuật toán và các tham số toán học liên quan (như modulus $N$, số mũ $e$ của RSA, hoặc định danh đường cong elip của ECDSA). Trường này cho phép relying party trích xuất đúng khóa và thuật toán để xác thực chữ ký của chủ thể mà không cần biết trước cấu hình.
\end{itemize}

\subsection{Các phần mở rộng X.509 v3 quan trọng (Standard Extensions)}
\label{subsec:standard-extensions}

Phần mở rộng v3 cho phép bổ sung các thuộc tính kiểm soát chính sách và bảo mật linh hoạt. Cấu trúc ASN.1 chuẩn của mỗi phần mở rộng gồm:
\begin{verbatim}
Extension ::= SEQUENCE {
    extnID      OBJECT IDENTIFIER,
    critical    BOOLEAN DEFAULT FALSE,
    extnValue   OCTET STRING }
\end{verbatim}

Mỗi extension gồm mã định danh (\texttt{extnID}), cờ \texttt{critical} và giá trị dữ liệu (\texttt{extnValue}). Nếu \texttt{critical: TRUE}, validator bắt buộc phải nhận diện và xử lý extension đó; nếu không hiểu, validator bắt buộc phải từ chối chứng thư. Các phần mở rộng cốt lõi bao gồm:
\begin{itemize}
    \item \textbf{Subject Alternative Name (SAN -- OID \texttt{2.5.29.17}):} Cho phép gán thêm nhiều danh tính cho chủ thể dưới các định dạng như tên miền (\texttt{dNSName}) hoặc địa chỉ IP (\texttt{iPAddress}). Các client hiện đại đối chiếu tên máy chủ với SAN, không dựa vào CN (kiểm chứng tại ca kiểm thử TC04).

    \item \textbf{Ràng buộc cơ bản -- Basic Constraints (OID \texttt{2.5.29.19}):} Xác định thực thể có quyền đóng vai trò CA để ký cấp chứng thư hay không qua cờ \texttt{cA} và giới hạn độ sâu tùy chọn \texttt{pathLenConstraint}. Chứng thư CA phải có extension này và phải critical (RFC~5280 mục 4.2.1.9). Thực thể cuối có thể không có extension này hoặc có \texttt{cA:FALSE}; extension tự nó không ngăn chặn, mà validator từ chối dùng chứng thư \texttt{cA:FALSE} làm CA (ràng buộc này ảnh hưởng trực tiếp đến việc xây dựng chuỗi ở ca TC06).

    \item \textbf{Mục đích sử dụng khóa -- Key Usage (OID \texttt{2.5.29.15}):} Chuỗi bit giới hạn các thao tác mật mã cơ sở của khóa công khai (như \texttt{digitalSignature}, \texttt{keyCertSign}, \texttt{cRLSign}), thường đặt \texttt{critical: TRUE} để ngăn ngừa dùng một khóa cho các ngữ cảnh mâu thuẫn. Cần lưu ý rằng \texttt{keyEncipherment} không cần cho TLS 1.3/ECDHE (chỉ cần \texttt{digitalSignature}).

    \item \textbf{Mục đích sử dụng khóa mở rộng (EKU -- OID \texttt{2.5.29.37}):} Xác định giao thức ứng dụng cấp cao mà khóa được áp dụng qua danh sách OID, tiêu biểu như \texttt{serverAuth} (xác thực máy chủ TLS) hoặc \texttt{codeSigning} (ký mã nguồn). Kết nối TLS máy chủ bị từ chối theo chính sách xác minh khi EKU có mặt mà không chứa \texttt{serverAuth} (kiểm chứng tại ca TC08). Chứng thư không có EKU thường không bị giới hạn mục đích.
\end{itemize}

\subsection{Trích xuất thực tế cấu trúc chứng thư trong môi trường Lab}
\label{subsec:lab-cert-extraction}

Dưới đây là kết quả giải mã cấu trúc chi tiết của chứng thư máy chủ TLS (\texttt{server-a.crt}) được cấp bởi Intermediate CA trong môi trường lab bằng lệnh:
\begin{verbatim}
$ openssl x509 -in demo/out/servers/server-a.crt -text -noout
\end{verbatim}

\begin{figure}[!htbp]
    \centering
    \includegraphics[width=0.85\textwidth, height=0.78\textheight, keepaspectratio]{openssl-x509-terminal-real.png}
    \caption{Kết quả thực thi lệnh trích xuất cấu trúc chứng thư X.509 v3 của máy chủ TLS (\texttt{server-a.crt}) trong môi trường thực nghiệm.}
    \label{fig:openssl-x509-terminal-real}
\end{figure}

\noindent\textbf{Phân tích từ kết quả thực nghiệm:}

Từ kết quả trích xuất trên Hình~\ref{fig:openssl-x509-terminal-real}, các trường cấu trúc của chứng thư máy chủ phản ánh đúng các thiết lập chính sách đã nghiên cứu:
\begin{itemize}
    \item \textbf{Phiên bản và Số sê-ri:} Chứng thư mang phiên bản \texttt{3 (0x2)} cho phép sử dụng các extension v3, và số sê-ri \texttt{4096 (0x1000)} được quản lý trong cơ sở dữ liệu \texttt{index.txt} của Intermediate CA.
    \item \textbf{Cặp định danh Issuer và Subject:} Thể hiện rõ mô hình phân tầng: \texttt{Issuer} trỏ về \texttt{CN = Lab Intermediate CA}, còn \texttt{Subject} định danh máy chủ dịch vụ \texttt{CN = server-a.test}.
    \item \textbf{SAN:} Chỉ định tên miền hợp lệ \texttt{DNS:server-a.test}, đáp ứng yêu cầu đối chiếu tên máy chủ của client (đối chứng với ca kiểm thử TC04).
    \item \textbf{Basic Constraints:} Thiết lập \texttt{critical, CA:FALSE}, định vị đây là chứng thư thực thể cuối (End-entity) và không có quyền ký phát chứng thư phái sinh.
    \item \textbf{Key Usage và EKU:} Bật \texttt{Digital Signature}, \texttt{Key Encipherment}; \texttt{EKU} kích hoạt \texttt{serverAuth}, đáp ứng điều kiện xác thực máy chủ TLS (đối chứng ca TC08).
    \item \textbf{Định danh khóa (SKI và AKI):} Trường SKI và AKI chứa mã băm nhận diện khóa công khai của máy chủ và của Intermediate CA, hỗ trợ thuật toán xây dựng và neo chuỗi tin cậy.
\end{itemize}

\section{Xây dựng và xác minh đường dẫn}
\label{sec:path-validation}

Một chứng thư máy chủ không tự mang đủ bằng chứng để được tin cậy. Relying party phải chỉ ra được một đường dẫn chứng thư (certification path) nối chứng thư đó tới một trust anchor mà chính nó đã cấu hình. Việc này gồm hai giai đoạn: xây dựng đường dẫn (path building) và xác minh đường dẫn (path validation, RFC~5280~\cite{rfc5280}, mục~6).

\subsection{Chuỗi chứng thư trong lab}
\label{subsec:lab-trust-chain}

Lab dùng chuỗi ba mắt xích $\mathrm{Root\ CA} \rightarrow \mathrm{Intermediate\ CA} \rightarrow \text{chứng thư máy chủ}$, như minh họa ở Hình~\ref{fig:trust-chain-validation}.
\begin{itemize}
    \item \textbf{Root CA} (\texttt{Lab Root CA}): Chứng thư tự ký ($\mathrm{Issuer} \equiv \mathrm{Subject}$), \texttt{basicConstraints} critical với \texttt{CA:TRUE}. Khóa $SK_{\mathrm{Root}}$ dùng để ký chứng thư của Intermediate CA.
    \item \textbf{Intermediate CA} (\texttt{Lab Intermediate CA}): Được Root ký phát với chữ ký $\sigma_{\mathrm{Sub}} = \mathrm{Sign}_{SK_{\mathrm{Root}}}(\mathrm{TBS}_{\mathrm{Sub}})$. Chứng thư này có cờ \texttt{CA:TRUE} và \texttt{pathlen:0}, nghĩa là chỉ được cấp chứng thư thực thể cuối, không được cấp thêm CA cấp dưới. Khóa $SK_{\mathrm{Sub}}$ dùng để ký chứng thư máy chủ.
    \item \textbf{Chứng thư máy chủ} (\texttt{server-a.test}): Được Intermediate CA ký phát với chữ ký $\sigma_{\mathrm{Server}} = \mathrm{Sign}_{SK_{\mathrm{Sub}}}(\mathrm{TBS}_{\mathrm{Server}})$, thiết lập \texttt{CA:FALSE}, SAN \texttt{DNS:server-a.test} và EKU \texttt{serverAuth}.
\end{itemize}

\clearpage

\begin{figure}[!htbp]
\centering
\resizebox{0.88\textwidth}{!}{%
\begin{tikzpicture}[
    % Định dạng các khối theo sắc xám học thuật nhất quán với Hình 3.1
    rootbox/.style={draw=black, very thick, double, fill=gray!18, rounded corners=6pt, text width=7.2cm, align=left, inner sep=5pt, font=\small, text=black},
    cabox/.style={draw=black, thick, fill=gray!10, rounded corners=6pt, text width=7.2cm, align=left, inner sep=5pt, font=\small, text=black},
    leafbox/.style={draw=black, thick, fill=white, rounded corners=6pt, text width=7.2cm, align=left, inner sep=5pt, font=\small, text=black},
    storebox/.style={draw=black, thick, dashed, fill=gray!6, rounded corners=6pt, text width=6.6cm, align=center, inner sep=5pt, font=\small, text=black},
    % Kiểu mũi tên nét liền (tiến trình ký) và nét đứt (tiến trình thẩm định)
    issuancearrow/.style={-{Stealth[scale=1.05]}, thick, black},
    valarrow/.style={-{Stealth[scale=1.1]}, very thick, dash pattern=on 5pt off 2.5pt, black},
    trustarrow/.style={-{Stealth[scale=1.1]}, very thick, dash pattern=on 6pt off 3pt, black}
]

    % --- KHỐI GỐC TIN CẬY (TRUST ANCHOR TRÊN CÙNG) ---
    \node[storebox] (store) at (0, 15.2) {
        \textbf{Kho chứng thư tin cậy (Root Trust Store)}\\[1pt]
        {\scriptsize Quản lý bởi Hệ điều hành / Trình duyệt (Relying Party)\\[1pt]Chứa bản sao khóa công khai $PK_{\mathrm{Root}}$ được tin cậy mặc định}
    };

    % --- KHỐI 1: Root CA ---
    \node[rootbox] (root) at (0, 10.4) {
        \textbf{Mắt xích 1: Root CA (Trust Anchor)}\\[1pt]
        {\scriptsize \texttt{Subject: CN = Lab Root CA}}\\[1pt]
        {\scriptsize \texttt{Issuer:  CN = Lab Root CA} (Tự ký - Self-signed)}\\[1pt]
        {\scriptsize Khóa: Cặp khóa gốc $(SK_{\mathrm{Root}}, PK_{\mathrm{Root}})$}\\[1pt]
        {\scriptsize \texttt{Basic Constraints: critical, CA:TRUE}}\\[1pt]
        {\scriptsize \texttt{SKI:} Mã băm nhận diện khóa $PK_{\mathrm{Root}}$}
    };

    % --- KHỐI 2: Intermediate CA ---
    \node[cabox] (subca) at (0, 5.2) {
        \textbf{Mắt xích 2: Intermediate CA (Sub-CA)}\\[1pt]
        {\scriptsize \texttt{Subject: CN = Lab Intermediate CA}}\\[1pt]
        {\scriptsize \texttt{Issuer:  CN = Lab Root CA}}\\[1pt]
        {\scriptsize Khóa: Cặp khóa trung gian $(SK_{\mathrm{Sub}}, PK_{\mathrm{Sub}})$}\\[1pt]
        {\scriptsize \texttt{Basic Constraints: critical, CA:TRUE, pathlen:0}}\\[1pt]
        {\scriptsize \texttt{AKI:} Trỏ tới $\mathrm{SKI}_{\mathrm{Root}}$ \textbar{} \texttt{SKI:} Băm khóa $PK_{\mathrm{Sub}}$}
    };

    % --- KHỐI 3: Server Certificate ---
    \node[leafbox] (server) at (0, 0.0) {
        \textbf{Mắt xích 3: Server Certificate (Thực thể cuối)}\\[1pt]
        {\scriptsize \texttt{Subject: CN = server-a.test}}\\[1pt]
        {\scriptsize \texttt{Issuer:  CN = Lab Intermediate CA}}\\[1pt]
        {\scriptsize Khóa: Khóa công khai $PK_{\mathrm{Server}}$ ($SK_{\mathrm{Server}}$ lưu an toàn trên Server)}\\[1pt]
        {\scriptsize \texttt{Basic Constraints: CA:FALSE} \textbar{} \texttt{EKU: serverAuth}}\\[1pt]
        {\scriptsize \texttt{SAN: DNS:server-a.test} \textbar{} \texttt{AKI:} Trỏ tới $\mathrm{SKI}_{\mathrm{Sub}}$}
    };

    % Neo Root CA vào Trust Store
    \draw[trustarrow] (root.north) -- (store.south)
        node[midway, right=8pt, align=left, font=\footnotesize\bfseries] {Neo niềm tin\\[1pt](Trust Anchor)};

    % ==============================================================
    % TIÊU ĐỀ HAI TIẾN TRÌNH Ở ĐỈNH (KHÔNG XUNG ĐỘT)
    % ==============================================================
    \node[align=center, font=\footnotesize\bfseries, text=black] at (-6.3, 15.2) {
        [TIẾN TRÌNH KÝ PHÁT]\\[2pt]
        {\scriptsize (Dùng khóa bí mật $SK$)}
    };

    \node[align=center, font=\footnotesize\bfseries, text=black] at (6.5, 15.2) {
        [TIẾN TRÌNH THẨM ĐỊNH]\\[2pt]
        {\scriptsize (Dùng khóa công khai $PK$)}
    };

    % ==============================================================
    % TIẾN TRÌNH KÝ PHÁT (BÊN TRÁI - TỪ TRÊN XUỐNG)
    % ==============================================================
    % Root CA ký Sub-CA
    \draw[issuancearrow] (root.west) -- (-4.5, 10.4) -- (-4.5, 5.2) -- (subca.west);
    \node[anchor=east, align=right, font=\footnotesize] at (-4.7, 7.8) {
        \textbf{1. Ký phát Sub-CA:}\\[1pt]
        $\sigma_{\mathrm{Sub}} = \mathrm{Sign}_{SK_{\mathrm{Root}}}(\mathrm{TBS}_{\mathrm{Sub}})$\\[1pt]
        {\scriptsize $SK_{\mathrm{Root}}$ ký lên chứng thư Sub-CA}
    };

    % Sub-CA ký Server Certificate
    \draw[issuancearrow] (subca.west) -- (-4.5, 5.2) -- (-4.5, 0.0) -- (server.west);
    \node[anchor=east, align=right, font=\footnotesize] at (-4.7, 2.6) {
        \textbf{2. Ký phát Server:}\\[1pt]
        $\sigma_{\mathrm{Server}} = \mathrm{Sign}_{SK_{\mathrm{Sub}}}(\mathrm{TBS}_{\mathrm{Server}})$\\[1pt]
        {\scriptsize $SK_{\mathrm{Sub}}$ ký lên chứng thư Server}
    };

    % ==============================================================
    % TIẾN TRÌNH THẨM ĐỊNH RFC 5280 (BÊN PHẢI - TỪ DƯỚI LÊN)
    % ==============================================================
    % Thẩm định Server Certificate
    \draw[valarrow] (server.east) -- (4.7, 0.0) -- (4.7, 5.2) -- (subca.east);
    \node[anchor=west, align=left, font=\footnotesize] at (4.9, 2.6) {
        \textbf{3. Thẩm định Server Certificate:}\\[1pt]
        $\mathrm{Verify}_{PK_{\mathrm{Sub}}}(\sigma_{\mathrm{Server}}) = \mathrm{True}$\\[1pt]
        {\scriptsize $\bullet$ Đối chiếu: $\mathrm{Issuer} \equiv \mathrm{Subject}_{\mathrm{Sub}}$}\\[1pt]
        {\scriptsize $\bullet$ Khớp khóa: $\mathrm{AKI}_{\mathrm{Server}} \equiv \mathrm{SKI}_{\mathrm{Sub}}$}
    };

    % Thẩm định Intermediate CA
    \draw[valarrow] (subca.east) -- (4.7, 5.2) -- (4.7, 10.4) -- (root.east);
    \node[anchor=west, align=left, font=\footnotesize] at (4.9, 7.8) {
        \textbf{4. Thẩm định Intermediate CA:}\\[1pt]
        $\mathrm{Verify}_{PK_{\mathrm{Root}}}(\sigma_{\mathrm{Sub}}) = \mathrm{True}$\\[1pt]
        {\scriptsize $\bullet$ Đối chiếu: $\mathrm{Issuer} \equiv \mathrm{Subject}_{\mathrm{Root}}$}\\[1pt]
        {\scriptsize $\bullet$ Khớp khóa: $\mathrm{AKI}_{\mathrm{Sub}} \equiv \mathrm{SKI}_{\mathrm{Root}}$}
    };

\end{tikzpicture}%
}

\caption{Chuỗi chứng thư Root CA $\rightarrow$ Intermediate CA $\rightarrow$ server trong lab: quan hệ ký phát và quan hệ thẩm định giữa các mắt xích.}
\label{fig:trust-chain-validation}
\end{figure}

\subsection{Quy trình xây dựng và xác minh}
\label{subsec:path-building-validation}

\noindent\textbf{Xây dựng đường dẫn.} Xuất phát từ chứng thư máy chủ, relying party tìm chứng thư của bên phát hành bằng cách khớp trường \texttt{Issuer} của chứng thư dưới với \texttt{Subject} của chứng thư trên, đồng thời khớp trường AKI với SKI (chỉ khớp tên là chưa đủ, vì một CA có thể có nhiều khóa cùng tên). Các chứng thư trung gian thường lấy từ chuỗi do máy chủ gửi kèm và chưa được tin cậy ở thời điểm này. Quá trình dừng khi gặp một chứng thư nằm trong kho trust anchor của client.

\noindent\textbf{Trust anchor.} Một Root CA tự ký chỉ trở thành trust anchor khi client được cấu hình để tin cậy nó (trong lab là tùy chọn \texttt{-CAfile}). Độ tin cậy của Root đến từ cấu hình đó, không đến từ chữ ký tự ký, vì bất kỳ ai cũng tạo được chứng thư tự ký cho chính mình.

\noindent\textbf{Xác minh đường dẫn.} Thuật toán của RFC~5280 quy định quy trình xử lý từng chứng thư trên đường dẫn từ trust anchor xuống chứng thư thực thể cuối qua 6 bước kiểm tra tuần tự:
\begin{enumerate}\setlength{\itemsep}{1pt}\setlength{\parskip}{0pt}
    \item \textbf{Xây dựng đường dẫn (Path Construction):} Truy vết các mắt xích từ chứng thư máy chủ ngược lên Trust Anchor dựa trên cặp định danh Issuer/Subject và AKI/SKI.
    \item \textbf{Xác minh chữ ký số (Signature Verification):} Dùng khóa công khai của bên cấp phát để kiểm tra tính toàn vẹn toán học của chữ ký trên từng chứng thư trong chuỗi.
    \item \textbf{Kiểm tra khung thời hạn hiệu lực (Validity Check):} Đối chiếu thời điểm xác minh với khoảng \texttt{[notBefore, notAfter]} của từng mắt xích.
    \item \textbf{Kiểm tra ràng buộc CA và extension liên quan (Constraints Check):} Thẩm tra \texttt{basicConstraints} (cờ \texttt{cA:TRUE} cho CA, giới hạn \texttt{pathLen}) và các ràng buộc chính sách.
    \item \textbf{Kiểm tra danh tính máy chủ và mục đích sử dụng (Identity \& Purpose Matching):} So khớp tên miền truy cập với SAN và đối chiếu cờ EKU (\texttt{serverAuth}) theo yêu cầu ứng dụng.
    \item \textbf{Kiểm tra trạng thái thu hồi (Revocation Status):} Tra cứu số sê-ri trong CRL hoặc truy vấn OCSP responder nếu chính sách kiểm tra thu hồi được kích hoạt.
\end{enumerate}

Cần lưu ý quan trọng rằng công cụ \texttt{openssl verify} không mặc định thực thi toàn bộ các bước trên. Hành vi xác minh phụ thuộc trực tiếp vào các tùy chọn dòng lệnh: kiểm tra hostname bắt buộc phải khai báo tường minh qua \texttt{-verify\_hostname}; kiểm tra mục đích TLS server cần chỉ định \texttt{-purpose sslserver}; và kiểm tra thu hồi đòi hỏi cờ \texttt{-crl\_check}. Nếu thiếu các tùy chọn này, lệnh chỉ thực hiện xác minh chuỗi cơ bản. Mỗi mắt xích phải đạt các điều kiện ở Bảng~\ref{tab:path-validation-conditions}. Chỉ cần một điều kiện không đạt là chứng thư bị từ chối theo chính sách xác minh được bật, dù các chữ ký vẫn đúng về mặt toán học. Đây là cơ sở của giả thuyết $H_1$ và được kiểm chứng bằng các ca TC02--TC08, TC10 ở Chương~\ref{chap:experiment}.

\begin{table}[!htbp]
\centering
\small
\renewcommand{\arraystretch}{1.25}
\begin{tabularx}{\textwidth}{>{\raggedright\arraybackslash}X>{\raggedright\arraybackslash}p{4.2cm}>{\raggedright\arraybackslash}p{2.8cm}}
\toprule
\textbf{Điều kiện kiểm tra} & \textbf{Áp dụng cho} & \textbf{Ca kiểm thử} \\
\midrule
Chữ ký của chứng thư được tạo bằng khóa của bên phát hành & Mọi chứng thư, trừ trust anchor & TC05 \\
\midrule
Chuỗi nối được tới trust anchor đã cấu hình & Toàn đường dẫn & TC02, TC06, TC10 \\
\midrule
Thời điểm xác minh nằm trong [\texttt{notBefore}, \texttt{notAfter}] & Mọi chứng thư & TC03 \\
\midrule
Ràng buộc CA: \texttt{CA:TRUE} với mắt xích trung gian, \texttt{pathLenConstraint} không bị vượt & Root, Intermediate & Không có ca riêng \\
\midrule
Mục đích sử dụng: Key Usage / EKU phù hợp (\texttt{serverAuth}) & Chứng thư máy chủ & TC08 \\
\midrule
Tên máy chủ khớp với SAN & Chứng thư máy chủ & TC04 \\
\midrule
Trạng thái thu hồi (chỉ khi được bật) & Mọi chứng thư & TC07 (xem Mục~\ref{sec:revocation-mechanisms}) \\
\bottomrule
\end{tabularx}
\caption{Các điều kiện xác minh đường dẫn và ca kiểm thử tương ứng}
\label{tab:path-validation-conditions}
\end{table}


Khi mọi điều kiện đều đạt, kết quả là PASS (TC01, TC09).

\newpage

\subsection{Quan sát từ lab}
\label{subsec:lab-observations}

Hình~\ref{fig:openssl-verify-terminal} trình bày hai lần chạy \texttt{openssl verify} với cùng chứng thư \texttt{server-a.crt}.

\begin{figure}[!htbp]
\centering
\includegraphics[width=0.90\textwidth]{openssl-verify-terminal.png}
\caption{Kết quả \texttt{openssl verify} với chuỗi đầy đủ (OK) và khi thiếu chứng thư Intermediate CA (lỗi 20, TC06).}
\label{fig:openssl-verify-terminal}
\end{figure}

\noindent\textbf{Lần 1 (đủ chuỗi):} Khi cung cấp Root qua \texttt{-CAfile} và Intermediate qua \texttt{-untrusted}, kết quả là \texttt{OK}, kèm \texttt{-purpose sslserver} và \texttt{-verify\_hostname server-a.test} nên EKU và SAN cũng được kiểm. Chuỗi in ra gồm ba mức: depth~0 là server, depth~1 là Lab Intermediate CA, depth~2 là Lab Root CA. Nhãn \texttt{(untrusted)} ở depth~0 và 1 nghĩa là hai chứng thư này không nằm trong kho tin cậy; chỉ Root là trust anchor.

\noindent\textbf{Lần 2 (thiếu Intermediate):} Khi bỏ \texttt{-untrusted}, OpenSSL báo \texttt{error 20 at 0 depth lookup: unable to get local issuer certificate}. Chữ ký trên chứng thư máy chủ vẫn đúng, nhưng không tìm được chứng thư của bên phát hành nên không dựng được đường dẫn tới trust anchor. Lỗi gắn với depth~0 vì đó là chứng thư không tìm được issuer. Đây là ca TC06.

Cần lưu ý hai điểm. Thứ nhất, cả hai lệnh không bật kiểm tra thu hồi, nên \texttt{OK} chỉ cho biết các điều kiện trong Bảng~\ref{tab:path-validation-conditions} (trừ thu hồi) đã đạt; trạng thái thu hồi được xét riêng ở Mục~\ref{sec:revocation-mechanisms}. Thứ hai, kết quả của TC06 phụ thuộc vào client: OpenSSL không tự bổ sung chứng thư trung gian còn thiếu, trong khi một số client khác có cơ chế bù (qua AIA hoặc bộ nhớ đệm). Điều này phù hợp với giới hạn của lab đã nêu ở Mục~6.4.

\section{Thu hồi và trạng thái chứng thư}
\label{sec:revocation-mechanisms}

Một chứng thư có chữ ký đúng và còn trong khoảng \texttt{[notBefore, notAfter]} vẫn có thể không còn đáng tin, ví dụ khi khóa bí mật bị lộ hoặc CA cấp sai (Mục~\ref{sec:certificate-lifecycle}, Bước~9). Relying party vì vậy cần một cách hỏi chứng thư đã bị thu hồi chưa. RFC~5280~\cite{rfc5280} định nghĩa danh sách thu hồi chứng thư (CRL), còn RFC~6960~\cite{rfc6960} định nghĩa giao thức trạng thái chứng thư trực tuyến (OCSP). Hai cơ chế trả lời cùng một câu hỏi nhưng khác nhau ở cách thông tin được phát hành và lấy về.

\subsection{Danh sách thu hồi chứng thư (CRL)}
\label{subsec:crl}

CRL là danh sách có chữ ký do CA (hoặc CRL issuer được ủy quyền) phát hành định kỳ (RFC~5280~\cite{rfc5280}, mục~5). Mỗi CRL chứa số sê-ri các chứng thư bị thu hồi, thời điểm thu hồi và mã lý do tùy chọn (ví dụ \texttt{keyCompromise}), cùng hai trường \texttt{thisUpdate} và \texttt{nextUpdate} cho biết CRL được tạo khi nào và dự kiến được thay khi nào. Chứng thư chỉ nơi công bố CRL qua extension \texttt{CRL Distribution Points}. Khi xác minh, relying party tải CRL, kiểm tra chữ ký bằng khóa của bên phát hành, kiểm tra CRL còn hiệu lực rồi tra số sê-ri của chứng thư (RFC~5280~\cite{rfc5280}, mục~6.3).

Điểm mạnh là client có thể dùng bản CRL đã lưu đệm đến hết \texttt{nextUpdate}, nên không cần liên lạc với CA ở mỗi kết nối. Điểm yếu là độ trễ: thu hồi chỉ có hiệu lực với client khi CRL mới được phát hành và tải về. Kích thước CRL cũng tăng theo số chứng thư bị thu hồi (RFC~5280 có cơ chế delta CRL để giảm bớt).

\subsection{Giao thức trạng thái chứng thư trực tuyến (OCSP)}
\label{subsec:ocsp}

Với OCSP, client gửi tới OCSP responder một truy vấn chứa định danh của chứng thư (băm tên và khóa của bên phát hành cùng số sê-ri). Địa chỉ responder nằm trong extension \texttt{Authority Information Access} (RFC~6960~\cite{rfc6960}). Responder trả về phản hồi có chữ ký với một trong ba trạng thái \texttt{good}, \texttt{revoked} hoặc \texttt{unknown}, kèm khoảng thời gian phản hồi có hiệu lực. Phản hồi do CA hoặc một responder được CA ủy quyền ký, và truy vấn có thể kèm nonce để chống phát lại. Máy chủ cũng có thể đính kèm sẵn phản hồi OCSP trong bắt tay TLS (OCSP stapling) để client không phải tự truy vấn (xem Mục~\ref{sec:trust-models-comparison}).

Điểm mạnh là thông tin theo từng chứng thư, nhỏ gọn và mới hơn CRL. Điểm yếu là cần responder trực tuyến (khả dụng), responder biết client đang kiểm tra chứng thư nào (riêng tư), và khi không liên lạc được responder thì client phải tự quyết định chấp nhận hay từ chối.

Trong Web PKI, BR từ phiên bản 2.0.1 (hiệu lực 15/03/2024) đã chuyển OCSP thành tùy chọn và đặt yêu cầu công bố CRL~\cite{cabforum_br}. Phân tích sâu hơn ở Mục~\ref{sec:trust-models-comparison}.

\subsection{So sánh CRL và OCSP}
\label{subsec:crl-vs-ocsp}

Bảng~\ref{tab:crl-vs-ocsp} tổng hợp so sánh giữa hai cơ chế CRL và OCSP.

\begin{table}[!htbp]
\centering
\small
\renewcommand{\arraystretch}{1.2}
\begin{tabularx}{\textwidth}{>{\raggedright\arraybackslash}p{3.2cm}>{\raggedright\arraybackslash}X>{\raggedright\arraybackslash}X}
\toprule
\textbf{Tiêu chí} & \textbf{CRL} & \textbf{OCSP} \\
\midrule
\textbf{Đơn vị thông tin} & Danh sách số sê-ri bị thu hồi của một CA & Trạng thái của một chứng thư \\
\midrule
\textbf{Độ mới} & Giới hạn bởi chu kỳ phát hành (\texttt{nextUpdate}) & Mới hơn, tùy cách responder cập nhật \\
\midrule
\textbf{Băng thông} & Lớn khi danh sách dài & Nhỏ mỗi truy vấn nhưng nhiều truy vấn \\
\midrule
\textbf{Riêng tư} & Không lộ chứng thư cụ thể đang kiểm tra & Responder biết chứng thư đang được kiểm tra \\
\midrule
\textbf{Khả dụng} & Dùng được bản đã lưu đệm đến \texttt{nextUpdate} & Phụ thuộc responder trực tuyến (trừ khi có stapling) \\
\midrule
\textbf{Trong lab} & Triển khai và kiểm thử (TC07) & Chỉ phân tích lý thuyết \\
\bottomrule
\end{tabularx}
\caption{So sánh CRL và OCSP}
\label{tab:crl-vs-ocsp}
\end{table}


\subsection{Trạng thái \texttt{good} không đồng nghĩa với chứng thư hợp lệ}
\label{subsec:good-vs-valid}

RFC~6960 nêu rõ: \texttt{good} chỉ cho biết tối thiểu là chứng thư chưa bị thu hồi, không nhất thiết có nghĩa là chứng thư từng được cấp hay thời điểm phản hồi nằm trong thời hạn của chứng thư (RFC~6960~\cite{rfc6960}, mục~2.2). OCSP chỉ trả lời một câu hỏi trong Bảng~\ref{tab:path-validation-conditions} (trạng thái thu hồi) và không thay thế các câu còn lại: chuỗi tới trust anchor, \texttt{[notBefore, notAfter]}, Basic Constraints, EKU và SAN. Một chứng thư hết hạn, sai tên máy chủ hoặc sai EKU vẫn có thể nhận \texttt{good}. Nguyên lý này áp dụng cả cho CRL: không thấy số sê-ri trong CRL chỉ có nghĩa là chưa bị liệt kê.

Sự cố DigiNotar minh họa rủi ro này. Theo Fox-IT~\cite{prins2011diginotar}, hành vi thông thường theo RFC~2560 (tiền thân của RFC~6960) là trả \texttt{good} cho số sê-ri không có bản ghi. Vì vậy các chứng thư giả không có trong hồ sơ vẫn có thể qua bước kiểm tra thu hồi. Ngày 1/9/2011, responder được chuyển sang danh sách trắng, trả \texttt{revoked} cho mọi số sê-ri không xác định.

\noindent\textbf{Liên hệ lab.} Lab chỉ kiểm tra CRL: Intermediate CA thu hồi chứng thư bằng \texttt{openssl ca -revoke}, phát hành CRL bằng \texttt{openssl ca -gencrl}, và client xác minh với CRL được chỉ định tường minh. TC07 kiểm tra rằng chứng thư đã bị liệt kê bị từ chối khi kiểm tra CRL được bật; kết quả cụ thể ở Mục~\ref{sec:exp-results}. Lệnh \texttt{openssl verify} mặc định không kiểm tra thu hồi, nên với client không bật CRL, chứng thư bị thu hồi vẫn có thể qua các bước còn lại. Điều này khớp với giả thuyết $H_2$.

\section{Liên hệ TLS}
\label{sec:tls-relation}

Các mục trước xét việc xác minh chứng thư một cách độc lập. Trong thực tế, chứng thư được trình bày trong bắt tay TLS 1.3~\cite{rfc8446}, và việc một chứng thư ``được chấp nhận'' là kết quả của nhiều kiểm tra do các bên khác nhau thực hiện. Mục này làm rõ vai trò của chứng thư và thông điệp \texttt{CertificateVerify} trong bắt tay, rồi phân định kiểm tra nào thuộc giao thức, thư viện, ứng dụng hay chính sách riêng của trình duyệt.

\subsection{Chứng thư và \texttt{CertificateVerify} trong bắt tay TLS 1.3}
\label{subsec:tls-cert-verify}

Sau khi hai bên thỏa thuận khóa phiên (ClientHello và ServerHello mang \texttt{key\_share}), máy chủ gửi ba thông điệp xác thực (RFC~8446~\cite{rfc8446}, mục~4.4):
\begin{itemize}
    \item \textbf{\texttt{Certificate}:} Chuỗi chứng thư, chứng thư thực thể cuối đứng đầu, tiếp theo là các chứng thư trung gian. Chứng thư của trust anchor có thể được lược bỏ khi client đã có sẵn (RFC~8446~\cite{rfc8446}, mục~4.4.2). Đây là lý do thiếu Intermediate (TC06) khiến client không dựng được đường dẫn.
    \item \textbf{\texttt{CertificateVerify}:} Chữ ký của máy chủ, tạo bằng khóa bí mật tương ứng với khóa công khai trong chứng thư thực thể cuối, tính trên giá trị băm của toàn bộ các thông điệp bắt tay trước đó (RFC~8446~\cite{rfc8446}, mục~4.4.3).
    \item \textbf{\texttt{Finished}:} Mã xác thực (MAC) trên toàn bộ bắt tay, xác nhận hai bên cùng suy ra một bộ khóa và một bản ghi bắt tay (RFC~8446~\cite{rfc8446}, mục~4.4.4).
\end{itemize}

Hai chữ ký trong bắt tay có vai trò khác nhau và bổ sung cho nhau:
\begin{itemize}
    \item Chữ ký của CA trên chứng thư trả lời câu hỏi ``khóa này thuộc về tên này theo CA mà client tin cậy''.
    \item \texttt{CertificateVerify} trả lời câu hỏi ``bên đang kết nối thật sự nắm khóa bí mật tương ứng, trong chính phiên này'', vì nội dung được ký bao gồm cả \texttt{key\_share} của phiên.
\end{itemize}

Thiếu \texttt{CertificateVerify}, kẻ tấn công có thể trình lại chứng thư công khai của máy chủ mà không có khóa bí mật. Thiếu xác minh đường dẫn, kẻ tấn công trình chứng thư tự tạo kèm chữ ký đúng, đúng như kịch bản thay thế khóa của Eve ở Mục~\ref{subsec:eve-mitm-attack}. Cần cả hai thì kết nối mới chống được kịch bản đó.

Hai chữ ký này cũng khác nhau về thuật toán. Chữ ký \texttt{CertificateVerify} phải dùng các thuật toán client chấp nhận trong \texttt{signature\_algorithms}, và với RSA thì dùng RSA-PSS (RFC~8446~\cite{rfc8446}, mục~4.2.3). Chữ ký của CA trên chứng thư (ví dụ \texttt{sha256WithRSAEncryption} ở Hình~\ref{fig:openssl-x509-terminal-real}) có thể dùng PKCS\#1 v1.5.

Giới hạn quan trọng là \texttt{CertificateVerify} chỉ chứng minh sở hữu khóa, không chứng minh chứng thư được cấp đúng. Với DigiNotar, kẻ tấn công vừa có chứng thư giả \texttt{*.google.com} do CA bị chiếm ký, vừa nắm khóa bí mật tương ứng, nên mọi kiểm tra mật mã trong bắt tay đều đạt. Theo Fox-IT~\cite{prins2011diginotar}, sự cố lộ ra vì Chrome có kiểm tra bổ sung (certificate pinning) cho chứng thư của Google. Đây là cơ sở cho Mục~\ref{sec:valid-signature-limits}.

\subsection{Phân định các kiểm tra}
\label{subsec:tls-checks}

RFC~8446 chỉ quy định phần xác thực nằm trong giao thức. Việc xác minh chứng thư chi tiết được dẫn tới RFC~5280 (RFC~8446~\cite{rfc8446}, mục~4.4.2.4; \cite{rfc5280}). Việc đối chiếu tên máy chủ do RFC~9525~\cite{rfc9525} quy định, và RFC này chỉ xét tên trong chứng thư máy chủ cuối, không xét việc xác minh chuỗi. Bảng~\ref{tab:tls-checks} phân định ai chịu trách nhiệm kiểm tra gì.

\begin{table}[!htbp]
\centering
\small
\renewcommand{\arraystretch}{1.12}
\begin{tabularx}{\textwidth}{>{\raggedright\arraybackslash}p{3.5cm}>{\raggedright\arraybackslash}X>{\raggedright\arraybackslash}p{2.5cm}>{\raggedright\arraybackslash}p{2.3cm}}
\toprule
\textbf{Kiểm tra} & \textbf{Bên thực hiện} & \textbf{Nguồn quy định} & \textbf{Trong lab} \\
\midrule
Cấu trúc thông điệp \texttt{Certificate}, kiểu chứng thư, khóa tương thích & Giao thức TLS & RFC~8446~\cite{rfc8446} & Không kiểm tra \\
\midrule
\texttt{CertificateVerify} (sở hữu khóa) và \texttt{Finished} & Giao thức TLS & RFC~8446~\cite{rfc8446} & Không kiểm tra \\
\midrule
Dựng và xác minh đường dẫn, thời hạn, Basic Constraints, Key Usage/EKU & Thư viện TLS, với trust anchor, thời điểm và mục đích sử dụng do ứng dụng cấu hình & RFC~5280~\cite{rfc5280} & TC01--TC03, TC05, TC06, TC08, TC10 \\
\midrule
Trạng thái thu hồi & Ứng dụng quyết định có kiểm tra không, dùng CRL hay OCSP, và hành vi khi lỗi; TLS chỉ cho phép đính kèm phản hồi OCSP (stapling) & RFC~5280~\cite{rfc5280}, RFC~6960~\cite{rfc6960} & TC07 (chỉ CRL) \\
\midrule
Khớp tên máy chủ với SAN & Ứng dụng, hoặc thư viện khi ứng dụng yêu cầu & RFC~9525~\cite{rfc9525} & TC04 \\
\midrule
Chính sách riêng của trình duyệt và Root Store: danh mục trust anchor, yêu cầu CT, pinning, giới hạn thời hạn chứng thư & Trình duyệt hoặc nhà cung cấp Root Store & Ngoài RFC~5280 và 8446; một phần theo BR~\cite{cabforum_br} & Không có \\
\bottomrule
\end{tabularx}
\caption{Phân định các kiểm tra trong kết nối TLS 1.3}
\label{tab:tls-checks}
\end{table}


Hai hệ quả cần nhấn mạnh:
\begin{itemize}
    \item Nhiều thư viện không tự bật kiểm tra tên máy chủ hay kiểm tra thu hồi; ứng dụng phải yêu cầu (trong lab là \texttt{-verify\_hostname} và kiểm tra CRL được bật tường minh). Vì vậy một kết nối TLS ``thành công'' chưa chắc đã qua đủ các kiểm tra ở Bảng~\ref{tab:path-validation-conditions}.
    \item Cùng một chứng thư có thể được chấp nhận ở client này và bị từ chối ở client khác tùy cấu hình và chính sách. Các ca kiểm thử trong lab phản ánh tầng ``thư viện OpenSSL kèm cấu hình'' chứ không phải chính sách của trình duyệt (xem Mục~\ref{sec:exp-limitations}).
\end{itemize}

\noindent\textbf{Phạm vi của công cụ trong lab.} Lệnh \texttt{openssl verify} đọc các chứng thư được cung cấp và chạy thuật toán xác minh đường dẫn theo hồ sơ PKIX (RFC~5280~\cite{rfc5280}) tại một thời điểm xác định. Lệnh này không thiết lập bắt tay TLS, không trao đổi \texttt{key\_share}, không tính giá trị băm của bản ghi bắt tay và không kiểm tra việc nắm giữ khóa bí mật thông qua \texttt{CertificateVerify} (RFC~8446~\cite{rfc8446}). Vì vậy kết quả TC01--TC10 phản ánh tầng xác minh chứng thư, không phản ánh toàn bộ một kết nối TLS. Việc kiểm tra bằng \texttt{openssl s\_server} và \texttt{openssl s\_client} là hướng mở rộng.

\chapter{Phân tích bảo mật, mô hình tấn công và so sánh}
\label{chap:threat-security}
\section{Tài sản, đối thủ và rủi ro}
\label{sec:threat-model}
Cần lập bảng tài sản cần bảo vệ (khóa CA, khóa máy chủ, trust store, quy trình xác minh) và các mối đe dọa: lộ khóa, CA cấp sai / bị xâm nhập (rogue certificates, phân tích sâu trường hợp DigiNotar~\cite{prins2011diginotar}), MITM, chứng thư hết hạn, kiểm tra thu hồi không đầy đủ.

\section{So sánh các cơ chế tin cậy và thu hồi}
Cần so sánh CA PKI, Web of Trust, TOFU theo giả định tin cậy; sau đó so sánh CRL, OCSP và chứng thư thời hạn ngắn theo độ mới, chi phí, khả dụng. Không kết luận một mô hình luôn vượt trội trong mọi môi trường.

\section{Giới hạn của chữ ký hợp lệ}
Cần liên hệ với kết quả ở Chương~\ref{chap:experiment}: một chữ ký đúng không tự chứng minh trust anchor, hostname, hiệu lực, EKU hoặc trạng thái thu hồi đều thỏa chính sách client.

\chapter{Thiết kế và đánh giá thực nghiệm}
\label{chap:experiment}

\section{Giả thuyết và môi trường}
\label{sec:exp-hypotheses-env}

\subsection{Các giả thuyết nghiên cứu}
\label{subsec:exp-hypotheses}
Để giải quyết các câu hỏi nghiên cứu (RQ) đặt ra tại Mục~\ref{sec:objectives-rqs}, thực nghiệm được thiết kế nhằm kiểm chứng ba giả thuyết kỹ thuật cốt lõi:
\begin{itemize}
    \item \textbf{Giả thuyết $H_1$ (Điều kiện cần của chữ ký số và vai trò chính sách client):} Chữ ký số hợp lệ chỉ là điều kiện cần chứ chưa đủ để một chứng thư số được client chấp nhận. Chứng thư chỉ được tin cậy khi thỏa mãn đồng thời các ràng buộc ngữ cảnh mà chính sách xác minh của client thực thi (như chuỗi tin cậy dẫn tới Trust Anchor được cấu hình, thời điểm kiểm tra nằm trong hạn hiệu lực, định danh máy chủ khớp SAN, mục đích sử dụng khớp EKU, và trạng thái thu hồi nếu client kích hoạt kiểm tra).
    \item \textbf{Giả thuyết $H_2$ (Tính phụ thuộc ngữ cảnh và chính sách của quyết định xác minh):} Quyết định thẩm định của client mang tính động và phụ thuộc trực tiếp vào các tham số ngữ cảnh (thời điểm kiểm tra, tên máy chủ, Trust Anchor chỉ định, mục đích kiểm tra) cùng chính sách kiểm tra thu hồi; việc thay đổi các tham số hoặc chính sách này trên cùng một chứng thư số sẽ trực tiếp làm đảo ngược quyết định thẩm định từ chấp thuận (\texttt{PASS}) sang từ chối (\texttt{FAIL}) hoặc ngược lại.
    \item \textbf{Giả thuyết $H_3$ (Đánh đổi về chi phí vận hành và tính khả thi của cơ chế thu hồi):} Cơ chế kiểm tra thu hồi qua danh sách thu hồi (CRL) bộc lộ sự đánh đổi rõ rệt về chi phí vận hành; kích thước tệp danh sách và thời gian xử lý thẩm định của client có xu hướng gia tăng theo số lượng chứng thư bị thu hồi.
\end{itemize}

\subsection{Môi trường thực nghiệm}
\label{subsec:exp-environment}
Toàn bộ quy trình sinh khóa, cấp phát chứng thư, thu hồi và thực thi bộ kiểm thử tự động được triển khai trên môi trường máy trạm với các thông số trích xuất trực tiếp từ hệ thống:
\begin{itemize}
    \item \textbf{Hệ điều hành \& Nhân Kernel:} Fedora Linux 43 (Container Image), chạy trên nền tảng WSL2 với nhân \texttt{Linux 6.6.87.2-microsoft-standard-WSL2} kiến trúc \texttt{x86\_64}.
    \item \textbf{Công cụ mật mã:} OpenSSL phiên bản \texttt{3.5.9} (phát hành ngày 29/09/2026, biên dịch ngày 30/09/2026 UTC), sử dụng thư viện OpenSSL 3.5.9 và cấu hình crypto-policy hệ thống tiêu chuẩn (\texttt{OPENSSLDIR: "/etc/pki/tls"}).
    \item \textbf{Thiết bị phần cứng:} Bộ vi xử lý Intel(R) Core(TM) i7-11800H @ 2.30GHz (thế hệ thứ 11, 16 vCPU); dung lượng bộ nhớ RAM ảo hóa cấp phát cho WSL2 là 3.7~GiB (khả dụng 3.1~GiB).
    \item \textbf{Thời điểm thực thi kiểm thử:} 01:58:15 (UTC+7), Thứ Tư, ngày 28 tháng 10 năm 2026.
\end{itemize}
Thông số môi trường thực tế được tự động kiểm tra và ghi lại qua kịch bản kiểm tra môi trường, thể hiện tổng hợp tại Hình~\ref{fig:env-summary}.

\begin{figure}[!htbp]
\centering
\includegraphics[width=0.92\textwidth]{env-04-summary-environment.png}
\caption{Thông số môi trường hệ điều hành, phần cứng và phiên bản OpenSSL thực tế của bài lab.}
\label{fig:env-summary}
\end{figure}

\section{Mô hình và quy trình}
\label{sec:exp-model-process}

\subsection{Mô hình phân cấp chứng thư (PKI Hierarchy)}
\label{subsec:pki-hierarchy}
Hệ thống phòng lab thiết lập mô hình PKI phân cấp 2 tầng theo chuẩn RFC~5280~\cite{rfc5280} nhằm mô phỏng cấu trúc hai tầng phổ biến trong các hệ thống PKI, được minh họa tại Hình~\ref{fig:lab-pki-hierarchy}.

\begin{figure}[!htbp]
\centering
\resizebox{0.95\textwidth}{!}{%
\begin{tikzpicture}[
    rootbox/.style={draw=black, very thick, double, fill=gray!18, rounded corners=6pt, text width=7.4cm, align=center, inner sep=6pt, font=\small, text=black},
    cabox/.style={draw=black, thick, fill=gray!10, rounded corners=6pt, text width=7.4cm, align=center, inner sep=6pt, font=\small, text=black},
    serverbox/.style={draw=black, thick, fill=white, rounded corners=6pt, text width=6.8cm, align=center, inner sep=6pt, font=\small, text=black},
    anomalybox/.style={draw=black, thick, dashed, fill=gray!6, rounded corners=6pt, text width=6.8cm, align=center, inner sep=6pt, font=\small, text=black},
    arrow/.style={-{Stealth[scale=1.1]}, thick, black},
    dashedarrow/.style={-{Stealth[scale=1.05]}, thick, dashed, black}
]

    % Tầng 1: Root CA
    \node[rootbox] (root) at (7.5, 10.2) {
        \textbf{Tổ chức chứng thực gốc (Root CA)}\\[2pt]
        {\scriptsize Tệp: \texttt{demo/out/root/root.crt} (RSA 4096-bit, Tự ký)}\\[1pt]
        {\scriptsize \texttt{basicConstraints = critical, CA:TRUE}}\\[1pt]
        {\scriptsize \texttt{keyUsage = critical, keyCertSign, cRLSign}}\\
        {\scriptsize \textit{(Nạp vào Trust Store của client qua cờ \texttt{-CAfile})}}
    };

    % Tầng 2: Intermediate CA
    \node[cabox] (subca) at (7.5, 5.8) {
        \textbf{Tổ chức chứng thực trung gian (Intermediate CA)}\\[2pt]
        {\scriptsize Tệp: \texttt{demo/out/intermediate/ca.crt} (RSA 4096-bit)}\\[1pt]
        {\scriptsize \texttt{basicConstraints = critical, CA:TRUE, pathlen:0}}\\[1pt]
        {\scriptsize \texttt{keyUsage = critical, keyCertSign, cRLSign}}\\
        {\scriptsize \textit{(Nạp vào client qua cờ \texttt{-untrusted} để dựng chuỗi tin cậy)}}
    };

    % Tầng 3 (Trái): End-Entity Servers hợp thức
    \node[serverbox] (servers) at (3.5, 0.3) {
        \textbf{Chứng thư máy chủ chuẩn (End-Entity)}\\[2pt]
        {\scriptsize Tệp: \texttt{server-a.crt}, \texttt{server-b.crt} (RSA 2048-bit)}\\[1pt]
        {\scriptsize \texttt{basicConstraints = CA:FALSE}}\\[1pt]
        {\scriptsize \texttt{keyUsage = digitalSignature, keyEncipherment}}\\[1pt]
        {\scriptsize \texttt{extendedKeyUsage = serverAuth}}\\[1pt]
        {\scriptsize \texttt{subjectAltName = DNS:server-a.test (server-b.test)}}
    };

    % Tầng 3 (Phải): Các chứng thư dị biệt phục vụ kiểm thử
    \node[anomalybox] (anomalies) at (11.5, 0.3) {
        \textbf{Chứng thư dị biệt / Kịch bản vi phạm}\\[2pt]
        {\scriptsize \texttt{client-only.crt}: EKU chỉ có \texttt{clientAuth} (TC08)}\\[1pt]
        {\scriptsize \texttt{tampered.crt}: Bị sửa đổi bit chữ ký số (TC05)}\\[1pt]
        {\scriptsize \texttt{rogue.crt}: Tự ký mạo danh, ngoài Root (TC02)}\\[1pt]
        {\scriptsize \texttt{other-ca/server.crt}: Thuộc CA độc lập khác (TC10)}
    };

    % Luồng từ Root CA xuống Intermediate CA
    \draw[arrow] (root.south) -- (subca.north)
        node[midway, right=6pt, font=\footnotesize\bfseries] {Ký phát hành CA con};

    % Tách nhánh từ Intermediate CA xuống tầng 3
    \coordinate (split) at (7.5, 3.2);
    \draw[thick] (subca.south) -- (split);

    \draw[arrow] (split) -| (servers.north)
        node[pos=0.35, above=3pt, font=\footnotesize\bfseries] {Ký cấp phát hợp thức};

    \draw[dashedarrow] (split) -| (anomalies.north)
        node[pos=0.35, above=3pt, font=\footnotesize\bfseries] {Cấp sai EKU / Can thiệp bit / Tự ký};

\end{tikzpicture}%
}

\caption{Mô hình phân cấp chứng thư X.509 hai tầng triển khai trong phòng lab thực nghiệm.}
\label{fig:lab-pki-hierarchy}
\end{figure}

\subsection{Cấu hình X.509 Extensions và Chính sách xác minh}
\label{subsec:x509-extensions-policy}
Để cô lập từng điều kiện kiểm thử, các chứng thư trong bài lab được cấu hình phần mở rộng (\textit{extensions}) theo mục đích cụ thể:
{\sloppy
\begin{itemize}
    \item \textbf{Root CA (\texttt{root.crt}):} Cấu hình \texttt{basicConstraints} = \texttt{critical, CA:TRUE}, \texttt{keyUsage} = \texttt{critical, keyCertSign, cRLSign}. Đóng vai trò là Trust Anchor độc lập, được nạp vào client qua cờ \texttt{-CAfile}.
    \item \textbf{Intermediate CA (\texttt{ca.crt}):} Cấu hình \texttt{basicConstraints} = \texttt{critical, CA:TRUE, pathlen:0}, \texttt{keyUsage} = \texttt{critical, keyCertSign, cRLSign}. Được nạp qua tham số \texttt{-untrusted} để chứng minh chuỗi dẫn xuất hợp lệ.
    \item \textbf{Chứng thư máy chủ chuẩn (\texttt{server-a.crt}, \texttt{server-b.crt}):} Cấu hình \texttt{basicConstraints} = \texttt{CA:FALSE}, \texttt{keyUsage} = \texttt{digitalSignature}, \texttt{keyEncipherment}, \texttt{extendedKeyUsage} = \texttt{serverAuth}, và \texttt{subjectAltName} = \texttt{DNS:server-a.test} (hoặc \texttt{DNS:server-b.test}).
    \item \textbf{Chứng thư dị biệt phục vụ kiểm thử:}
    \begin{itemize}
        \item \texttt{client-only.crt}: Cấu hình \texttt{extendedKeyUsage} = \texttt{clientAuth} (thiếu \texttt{serverAuth}) để phục vụ kiểm thử ràng buộc EKU (TC08).
        \item \texttt{rogue.crt}: Chứng thư tự ký mạo danh, không bắt nguồn từ Root CA trong trust store (TC02).
        \item \texttt{other-ca/server.crt}: Chứng thư được ký hợp thức nhưng bởi một CA hoàn toàn xa lạ không nằm trong danh mục \texttt{-CAfile} của client (TC10).
        \item \texttt{tampered.crt}: Chứng thư bị can thiệp đảo bit ở phần payload chữ ký số (TC05).
    \end{itemize}
\end{itemize}
}

\subsection{Quy trình thực nghiệm}
\label{subsec:exp-workflow}
Quy trình được tự động hóa hoàn toàn thông qua tệp điều khiển \texttt{Makefile} và kịch bản Bash theo các bước tuần tự:
\begin{enumerate}
    \item \texttt{make clean}: Dọn dẹp toàn bộ dữ liệu tạm và thư mục xuất bản cũ.
    \item \texttt{make ca}: Tạo cặp khóa RSA 4096-bit, khởi tạo cơ sở dữ liệu CA (\texttt{index.txt}, \texttt{serial}), và sinh chứng thư tự ký cho Root CA cùng Intermediate CA.
    \item \texttt{make servers}: Tạo khóa RSA 2048-bit, tạo CSR và ký cấp phát chứng thư cho các máy chủ mục tiêu với các extension tương ứng.
    \item \texttt{make revoke}: Thu hồi chứng thư \texttt{server-a.crt} bằng lệnh \texttt{openssl ca -revoke} với mã lý do \texttt{keyCompromise}, sau đó kết xuất danh sách thu hồi \texttt{ca.crl} bằng lệnh \texttt{openssl ca -gencrl}.
    \item \texttt{make demo-test}: Kích hoạt script \texttt{demo/scripts/run\_tests.sh} để thực thi 10 ca kiểm thử (TC01--TC10), đối soát exit code và lưu vết kết quả ra tệp \url{demo/out/results.tsv}. Trong đó, ca TC03 tính toán mốc tương lai qua tham số \texttt{-attime} với công thức \texttt{\$((\$(date +\%s) + 40 * 86400))} (+40 ngày kể từ lúc sinh, vượt hạn 30 ngày); ca TC04 truy vấn kiểm tra với tên miền \texttt{wrong.test} (không khớp với SAN \texttt{server-a.test}).
\end{enumerate}

\noindent\textbf{Lưu ý về thời hạn hiệu lực của chứng thư lab:} Toàn bộ chứng thư máy chủ trong bài lab được cấu hình thời hạn hiệu lực là 30 ngày (từ 24/10/2026 đến 23/11/2026). Sau mốc 23/11/2026, nếu tái lập kiểm thử mà không tạo mới chứng thư, các ca kiểm thử đạt mốc chuẩn (TC01 và TC09) sẽ chuyển sang trạng thái \texttt{FAIL} do lỗi hết hạn tự nhiên (\texttt{error 10: certificate has expired}). Để tái lập chính xác kết quả sau mốc thời gian này, người đánh giá có thể thực hiện dọn dẹp và khởi tạo lại toàn bộ lab bằng lệnh \texttt{make clean \&\& make demo-test}, hoặc truyền bổ sung tham số mốc thời gian cố định \texttt{-attime <epoch>} nằm trong khoảng hiệu lực ban đầu (ví dụ: \texttt{-attime 1792800000}).

\section{Test cases và kết quả}
\label{sec:exp-results}

\subsection{Bảng kết quả thực nghiệm quan sát}
\label{subsec:exp-results-table}
Quá trình kiểm thử thực hiện 10 ca kiểm thử độc lập (TC01--TC10) bằng công cụ \texttt{openssl verify}. Bảng~\ref{tab:exp-results} đối chiếu chi tiết giữa điều kiện kiểm thử, kết quả kỳ vọng và kết quả quan sát thực tế thu được từ hệ thống.

\begin{table}[!htbp]
\centering
\small
\renewcommand{\arraystretch}{1.12}
\setlength{\tabcolsep}{3.5pt}
\begin{tabularx}{\textwidth}{l >{\raggedright\arraybackslash}p{3.8cm} c c c >{\raggedright\arraybackslash\ttfamily\scriptsize}X c}
\toprule
\textbf{ID} & \textbf{Điều kiện kiểm thử cụ thể} & \textbf{Kỳ vọng} & \textbf{Thực tế} & \textbf{Exit} & \textrm{\textbf{Mã lỗi / Output quan sát được}} & \textbf{Đánh giá} \\
\midrule
\textbf{TC01} & Chuỗi tin cậy đầy đủ, hostname khớp SAN, EKU \texttt{serverAuth} hợp lệ & PASS & PASS & 0 & server-a.crt: OK & Khớp \\
\midrule
\textbf{TC02} & Chứng thư tự ký (không có chuỗi dẫn tới Trust Anchor cấu hình) & FAIL & FAIL & 2 & error 18: self-signed certificate & Khớp \\
\midrule
\textbf{TC03} & Thời điểm kiểm tra đặt tại mốc tương lai sau hạn (\texttt{-attime} $+40$ ngày) & FAIL & FAIL & 2 & error 10: certificate has expired & Khớp \\
\midrule
\textbf{TC04} & Hostname truy vấn (\texttt{wrong.test}) không trùng khớp SAN (\texttt{server-a.test}) & FAIL & FAIL & 2 & error 62: hostname mismatch & Khớp \\
\midrule
\textbf{TC05} & Chữ ký số trên chứng thư bị can thiệp trực tiếp vào payload nhị phân & FAIL & FAIL & 2 & error 7: certificate signature failure (invalid padding) & Khớp \\
\midrule
\textbf{TC06} & Bỏ qua chứng thư Intermediate CA trong tham số (\texttt{-untrusted}) & FAIL & FAIL & 2 & error 20: unable to get local issuer certificate & Khớp \\
\midrule
\textbf{TC07} & Chứng thư đã bị thu hồi và client bật kiểm tra CRL (\texttt{-crl\_check}) & FAIL & FAIL & 2 & error 23: certificate revoked & Khớp \\
\midrule
\textbf{TC08} & Chứng thư chỉ cấp EKU \texttt{clientAuth}, client đòi hỏi \texttt{sslserver} & FAIL & FAIL & 2 & error 26: unsuitable certificate purpose & Khớp \\
\midrule
\textbf{TC09} & Chứng thư còn hạn, chưa bị thu hồi, có nạp CRL hợp lệ & PASS & PASS & 0 & server-b.crt: OK & Khớp \\
\midrule
\textbf{TC10} & Chuỗi chứng thư hợp lệ nhưng không dẫn tới Root CA được cấu hình & FAIL & FAIL & 2 & error 20: unable to get local issuer certificate & Khớp \\
\bottomrule
\end{tabularx}
\caption{Kết quả quan sát thực nghiệm bộ 10 ca kiểm thử PKI (TC01--TC10).}
\label{tab:exp-results}
\end{table}


Hình~\ref{fig:suite-summary} trích xuất toàn văn kết quả thực thi tự động 10 ca kiểm thử từ giao diện dòng lệnh, thể hiện tính nhất quán giữa kết quả quan sát thực nghiệm và kỳ vọng thiết kế của bộ kiểm thử.

\begin{figure}[!htbp]
\centering
\includegraphics[width=\textwidth]{testcase-suite-summary.png}
\caption{Tổng hợp kết quả thực thi tự động bộ 10 ca kiểm thử PKI qua script \texttt{run\_tests.sh}.}
\label{fig:suite-summary}
\end{figure}

\subsection{Phân tích chi tiết các trường hợp quan sát}
\label{subsec:exp-analysis-cases}
Dựa trên kết quả quan sát tại Bảng~\ref{tab:exp-results}, các ca kiểm thử được phân nhóm và phân tích chuyên sâu theo nguyên tắc cô lập biến kiểm thử đối chiếu với mốc chuẩn TC01:
\begin{itemize}
    \item \textbf{Nhóm xác minh toàn vẹn mật mã (TC05):} So với mốc chuẩn TC01, ca kiểm thử này giữ nguyên mọi tham số và chuỗi xác minh, chỉ thay đổi duy nhất một bit trong payload chữ ký DER. Quá trình giải mã chữ ký bằng khóa công khai của bên phát hành gặp lỗi đệm PKCS\#1 v1.5 (\textit{invalid padding}), OpenSSL lập tức dừng thuật toán và báo \texttt{error 7: certificate signature failure}. Đây là bằng chứng cho thấy kiểm tra chữ ký là điều kiện tiên quyết bắt buộc.
    \item \textbf{Nhóm xác minh chuỗi tin cậy (TC02, TC06, TC10):} So với mốc chuẩn TC01, client giữ nguyên cấu hình xác minh nhưng nguồn gốc chứng thư mục tiêu hoặc chuỗi dẫn truyền bị thay đổi: bỏ qua cờ nạp Intermediate CA (\texttt{-untrusted}) ở TC06 dẫn đến đứt gãy chuỗi; chứng thư phát hành từ một CA độc lập không nằm trong Trust Store của client ở TC10; hoặc chứng thư tự ký mạo danh ở TC02. OpenSSL trả về \texttt{error 20} (\texttt{unable to get local issuer certificate}) cho TC06, TC10 và \texttt{error 18} (\texttt{self-signed certificate}) cho TC02.
    \item \textbf{Nhóm xác minh ngữ cảnh ứng dụng (TC03, TC04, TC08):} So với mốc chuẩn TC01 (sở hữu chuỗi tin cậy và chữ ký chuẩn), các ca này cô lập và chỉ thay đổi duy nhất một biến tham số phía client: mốc thời gian thẩm định vượt hạn (\texttt{-attime}) ở TC03 gây ra \texttt{error 10}; tên miền truy vấn (\texttt{wrong.test}) ở TC04 gây ra \texttt{error 62}; hoặc sử dụng chứng thư mang EKU chuyên biệt cho client (\texttt{client-only.crt}) ở TC08 gây ra \texttt{error 26} khi client đòi hỏi mục đích máy chủ (\texttt{-purpose sslserver}).
    \item \textbf{Nhóm kiểm tra trạng thái thu hồi (TC07, TC09):} So với mốc chuẩn TC01 (sử dụng cùng chứng thư \texttt{server-a.crt} vốn đã bị thu hồi trên CA), TC07 chỉ thay đổi duy nhất biến chính sách của client: bổ sung cờ kiểm tra thu hồi \texttt{-crl\_check} kèm tệp \texttt{ca.crl}, khiến OpenSSL lập tức từ chối với \texttt{error 23 (certificate revoked)}. Ngược lại, chứng thư \texttt{server-b.crt} chưa bị thu hồi vẫn vượt qua kiểm tra an toàn với cùng chính sách CRL nghiêm ngặt (TC09).
\end{itemize}

\section{Thảo luận và hạn chế}
\label{sec:exp-limitations}

\subsection{Đánh giá và trả lời các giả thuyết nghiên cứu}
\label{subsec:eval-hypotheses}

\paragraph{Đánh giá Giả thuyết $H_1$:}
Dữ liệu thực nghiệm từ Bảng~\ref{tab:exp-results} ủng hộ giả thuyết $H_1$ trong phạm vi 10 ca kiểm thử với OpenSSL. Trong số các ca bị từ chối, ngoại trừ TC05 bị hỏng chữ ký mật mã và TC02 là chứng thư tự ký (không có chuỗi dẫn tới Trust Anchor được cấu hình), toàn bộ 6 ca còn lại (TC03, TC04, TC06, TC07, TC08, TC10) đều sở hữu chữ ký số hợp lệ về mặt toán học do CA ký phát hành. Dù vậy, cả 6 ca này đều bị OpenSSL trả về mã lỗi (\texttt{exit code 2}) và từ chối xác minh do vi phạm các ràng buộc chính sách. Đáng chú ý, tại mốc chuẩn TC01, chứng thư \texttt{server-a.crt} thực chất đã nằm trong danh sách thu hồi của CA nhưng vẫn được OpenSSL chấp nhận (\texttt{OK}) do client không kích hoạt kiểm tra CRL; trong khi ở TC09, chứng thư \texttt{server-b.crt} thỏa mãn đồng thời cả chuỗi tin cậy, ngữ cảnh và trạng thái chưa bị thu hồi trên CRL nên đạt \texttt{OK}. Kết quả này xác nhận chữ ký số hợp lệ chỉ là điều kiện cần; quyết định chấp nhận chứng thư phụ thuộc hoàn toàn vào các ràng buộc do chính sách của client thực thi.

\paragraph{Đánh giá Giả thuyết $H_2$:}
Dữ liệu thực nghiệm xác nhận tính tương đối và phụ thuộc ngữ cảnh của quyết định xác minh: Thay đổi mốc thời gian kiểm tra (\texttt{-attime}) làm đảo ngược kết quả từ \texttt{OK} sang \texttt{error 10}; thay đổi tên máy chủ truy vấn làm phát sinh \texttt{error 62}; và việc kích hoạt cờ kiểm tra thu hồi (\texttt{-crl\_check}) ở TC07 trực tiếp chặn đứng chứng thư đã thu hồi mà trước đó đã vượt qua mốc chuẩn TC01. 

Để kiểm chứng sâu hơn bản chất nhân quả của $H_2$ trên cùng một chứng thư số mà không thay đổi mã định danh của 10 ca kiểm thử chính, hai lần chạy đối chứng bổ sung được thực hiện trực tiếp trên hệ thống:
\begin{itemize}
    \item \textit{Đối chứng thay đổi Trust Anchor trên cùng chứng thư \texttt{server-a.crt}:} Giữ nguyên chứng thư và Intermediate CA hợp lệ, nhưng thay đổi cờ \texttt{-CAfile} từ Root CA của hệ thống sang Root CA độc lập (\texttt{other-ca/root.crt}). Kết quả thẩm định bị từ chối tại độ sâu 1 với thông báo: \\
    \texttt{error 20: unable to get local issuer certificate} (\texttt{depth 1, exit 2}).
    \item \textit{Đối chứng thay đổi chính sách mục đích trên cùng chứng thư \texttt{client-only.crt}:} Tại TC08, chứng thư này bị từ chối khi client yêu cầu mục đích \texttt{sslserver} (\texttt{error 26}). Khi giữ nguyên chứng thư và chuỗi tin cậy nhưng thay đổi chính sách phía client sang \texttt{-purpose sslclient}, OpenSSL chấp thuận thẩm định thành công: \\
    \texttt{out/servers/client-only.crt: OK} (\texttt{exit code 0}).
\end{itemize}
Hai kết quả đối chứng này khẳng định rằng ngay cả khi chứng thư số và chữ ký không thay đổi, việc điều chỉnh chính sách kiểm tra của client (Trust Anchor được tin cậy hoặc mục đích sử dụng) sẽ trực tiếp đảo ngược quyết định thẩm định.

\paragraph{Đánh giá Giả thuyết $H_3$:}
Giả thuyết $H_3$ được kiểm chứng một phần qua thực nghiệm đo lường định lượng cục bộ với OpenSSL phiên bản 3.5.9, kết hợp phân tích kiến trúc lý thuyết RFC~5280. Để đánh giá chi phí vận hành của cơ chế CRL, đề tài tiến hành thực nghiệm đo lường dung lượng tệp danh sách thu hồi và thời gian thẩm định đường dẫn bằng lệnh \texttt{openssl verify} khi mở rộng số lượng mục thu hồi ($N$). Trên cùng cấu hình phần cứng ghi nhận tại Mục~\ref{subsec:exp-environment}, thời gian xác minh được lấy theo giá trị trung vị (\textit{median}) qua 10 lần thực thi độc lập:
\begin{itemize}
    \item Với $N = 0$ (CRL rỗng): Kích thước tệp CRL là 564 bytes ($\approx 0{,}55$~KB), thời gian thẩm định trung vị đạt 37{,}46~ms.
    \item Với $N = 100$: Kích thước tệp CRL đạt 5.309 bytes ($\approx 5{,}18$~KB), thời gian thẩm định trung vị là 35{,}71~ms.
    \item Với $N = 1.000$: Kích thước tệp CRL đạt 47.966 bytes ($\approx 46{,}84$~KB), thời gian thẩm định trung vị là 36{,}23~ms.
    \item Với $N = 10.000$: Kích thước tệp CRL tăng vọt lên 474.532 bytes ($\approx 463{,}41$~KB, gấp khoảng 841 lần so với $N=0$), thời gian thẩm định trung vị tăng lên 86{,}46~ms (tăng gấp hơn 2{,}3 lần).
\end{itemize}
Kết quả thực nghiệm định lượng cho thấy dung lượng tệp CRL tăng tuyến tính theo số lượng bản ghi thu hồi ($\approx 47$~bytes/bản ghi), và chi phí nạp/xử lý thẩm định của client bắt đầu tăng rõ rệt khi danh sách mở rộng lên quy mô lớn ($N = 10.000$). Điều này ủng hộ nhận định của $H_3$ về sự đánh đổi tài nguyên và giới hạn mở rộng của cơ chế CRL truyền thống, đòi hỏi các giải pháp tối ưu như OCSP hoặc danh sách thu hồi rút gọn.

\subsection{Hạn chế của bài lab}
\label{subsec:lab-limitations}
Mặc dù khảo sát và tái hiện được các điều kiện biên phổ biến theo chuẩn RFC~5280 trong môi trường lab, nghiên cứu tồn tại một số hạn chế kỹ thuật cần ghi nhận:
\begin{itemize}
    \item \textbf{Tính chất cục bộ của môi trường OpenSSL CLI:} Thực nghiệm sử dụng tiện ích dòng lệnh \texttt{openssl verify} phiên bản 3.5.9 trên mô hình Root CA nội bộ. Đây là phép thẩm định đường dẫn tĩnh ở tầng thư viện mật mã, không chạy qua tiến trình bắt tay mạng TLS 1.3 hoàn chỉnh (chưa kiểm tra các thông điệp \texttt{CertificateVerify} và \texttt{Finished}). Đồng thời, đề tài chưa mở rộng khảo sát đối sánh trên các thư viện mật mã khác (như GnuTLS, NSS, BoringSSL).
    \item \textbf{Phạm vi kiểm tra thu hồi của cờ \texttt{-crl\_check}:} Trong OpenSSL, cờ \texttt{-crl\_check} chỉ thẩm định trạng thái thu hồi của chứng thư đích cuối chuỗi (leaf certificate). Để kiểm tra trạng thái thu hồi của toàn bộ các chứng thư trung gian trong chuỗi tin cậy, client bắt buộc phải cấu hình cờ \texttt{-crl\_check\_all} và cung cấp đầy đủ tệp CRL phát hành từ Root CA.
    \item \textbf{Tính tiền định của dữ liệu tổng hợp:} Toàn bộ khóa và chứng thư số trong bài lab được sinh tổng hợp có chủ đích trong môi trường lab cô lập. Do các ca kiểm thử mang tính tiền định (deterministic), các điều kiện đầu vào giống nhau luôn mang lại kết quả thẩm định đồng nhất, vì vậy không đòi hỏi phương pháp lặp đo ngẫu nhiên thống kê trên các ca chức năng chính.
    \item \textbf{Khoảng cách với hành vi Web PKI thực tế:} Hành vi của các trình duyệt web và client hiện đại trong thực tế có sự phân hóa đáng kể và linh hoạt hơn so với tiện ích OpenSSL thuần túy:
    \begin{itemize}
        \item Một số client và trình duyệt có thể hỗ trợ cơ chế AIA Fetching để chủ động truy vấn và tải chứng thư Intermediate CA còn thiếu qua mạng (trong khi OpenSSL báo lỗi ngay ở TC06 nếu thiếu cờ \texttt{-untrusted}). Tuy nhiên, tính năng này khác nhau tùy theo cấu hình cụ thể của từng client và chưa được khảo sát trong bài lab.
        \item Nhiều trình duyệt web hiện đại đã loại bỏ việc tải và kiểm tra CRL trực tiếp từ điểm phân phối (CDP) theo cách OpenSSL thực hiện nhằm tránh rủi ro nghẽn mạng và độ trễ bắt tay; thay vào đó, các client áp dụng những cơ chế tối ưu như danh sách thu hồi rút gọn được đồng bộ định kỳ từ nhà phát triển trình duyệt (ví dụ Google CRLSet, Mozilla OneCRL) hoặc yêu cầu máy chủ gửi kèm phản hồi OCSP trong phiên TLS (OCSP Stapling). Chi tiết hành vi giữa các client trên Internet có sự khác biệt lớn và không hoàn toàn đồng nhất với các tham số dòng lệnh cục bộ của OpenSSL.
    \end{itemize}
\end{itemize}

\chapter{Ứng dụng, xu hướng và hướng phát triển PKI}
\label{chap:applications}

Các chương trước đi sâu vào trường hợp xác thực máy chủ TLS. Chương này mở rộng ra các lĩnh vực khác dùng PKI, trình bày cách PKI được áp dụng cho chữ ký số tại Việt Nam, rồi điểm qua các xu hướng đang thay đổi cách vận hành PKI: Certificate Transparency, tự động hóa và chứng thư ngắn hạn, mật mã hậu lượng tử. Phần cuối đề xuất hướng mở rộng cho đề tài.

\section{Các lĩnh vực ứng dụng PKI}
Mọi ứng dụng trong Bảng~\ref{tab:applications} dùng chung cấu trúc X.509, chuỗi tin cậy và thu hồi đã phân tích ở Chương~\ref{chap:x509}. Khác biệt nằm ở loại danh tính được chứng thực, mục đích sử dụng khóa (EKU) và thời điểm cần xét hiệu lực.

\begingroup\small
\begin{longtable}{L{2.6cm}L{2.9cm}L{2.6cm}L{4.0cm}}
\caption{Ứng dụng của PKI và đặc điểm xác minh}\label{tab:applications}\\
\toprule
Ứng dụng & Danh tính chứng thực & Mục đích (EKU) & Đặc điểm cần lưu ý \\
\midrule
\endfirsthead
\toprule
Ứng dụng & Danh tính chứng thực & Mục đích (EKU) & Đặc điểm cần lưu ý \\
\midrule
\endhead
HTTPS/TLS & Tên miền (SAN dNSName) & \texttt{serverAuth} & Web PKI công cộng; kiểm tra hostname; trường hợp nghiên cứu của báo cáo. \\
Xác thực client, VPN, Wi-Fi doanh nghiệp & Người dùng hoặc thiết bị & \texttt{clientAuth} & Thường dùng private PKI; thay mật khẩu bằng chứng thư (mTLS, 802.1X). TC08 cho thấy chứng thư loại này không dùng được cho máy chủ. \\
Chữ ký số văn bản, giao dịch điện tử & Cá nhân, tổ chức & Chống chối bỏ (\texttt{contentCommitment}) & Hiệu lực xét tại thời điểm ký nên cần dấu thời gian tin cậy~\cite{rfc3161}; cần lưu trữ dài hạn chuỗi và trạng thái thu hồi; khóa trên USB token, HSM hoặc ký từ xa. \\
Ký mã phần mềm & Nhà phát triển, tổ chức & \texttt{codeSigning} & Hệ điều hành kiểm tra trước khi cài hoặc chạy; dấu thời gian giúp chữ ký còn kiểm tra được sau khi chứng thư hết hạn; khóa ký bị lộ cho phép phát tán mã độc có chữ ký hợp lệ. \\
Thư điện tử (S/MIME) & Địa chỉ email & \texttt{emailProtection} & Ký và mã hóa thư đầu cuối~\cite{rfc8551}; người nhận cần chứng thư của người gửi để mã hóa. \\
Hộ chiếu, giấy tờ định danh điện tử & Cơ quan cấp giấy tờ & Riêng theo chuẩn & Chip lưu dữ liệu do cơ quan cấp ký; bên kiểm tra xác minh chữ ký để phát hiện giả mạo dữ liệu~\cite{icao9303}. \\
Thiết bị IoT & Thiết bị, nhà sản xuất & Tùy hệ thống & Chứng thư cài từ lúc sản xuất; vòng đời thiết bị dài, khó cập nhật trust anchor hoặc thay khóa. \\
\bottomrule
\end{longtable}
\endgroup

Điểm khác biệt quan trọng nhất là \emph{thời điểm xác minh}. Với TLS, client chỉ quan tâm chứng thư có hợp lệ ngay lúc kết nối hay không, như TC03 đã minh họa. Với chữ ký văn bản hoặc ký mã, câu hỏi là chứng thư có hợp lệ \emph{tại lúc ký} không, trong khi việc kiểm tra có thể diễn ra nhiều năm sau, khi chứng thư đã hết hạn. Dấu thời gian do một bên tin cậy ký (RFC 3161) gắn chữ ký với một thời điểm đã được chứng thực, nhờ đó vẫn kiểm tra được sau khi chứng thư hết hạn~\cite{rfc3161}.

\section{PKI và chữ ký số tại Việt Nam}
Tại Việt Nam, chữ ký số được dùng rộng rãi trong kê khai thuế, hải quan, bảo hiểm xã hội, hóa đơn điện tử, ngân hàng và dịch vụ công trực tuyến. Khung pháp lý hiện hành gồm Luật Giao dịch điện tử số 20/2023/QH15, có hiệu lực từ 01/7/2024~\cite{luatgdtd2023}, và Nghị định 23/2025/NĐ-CP quy định về chữ ký điện tử và dịch vụ tin cậy, ban hành ngày 21/02/2025, có hiệu lực từ 10/4/2025 và thay thế Nghị định 130/2018/NĐ-CP~\cite{nd232025}.

Về kiến trúc, hệ thống này là một PKI phân cấp đúng như mô hình ở Chương~\ref{chap:architecture}. Tổ chức cung cấp dịch vụ chứng thực chữ ký số quốc gia đóng vai trò Root CA và cấp chứng thư cho các tổ chức cung cấp dịch vụ chứng thực chữ ký số công cộng; các tổ chức này lại cấp chứng thư cho cá nhân và doanh nghiệp. Khác với Web PKI, trust anchor ở đây là Root CA quốc gia, được phần mềm ký và kiểm tra chữ ký của các cơ quan cài đặt theo quy định, chứ không dựa vào trust store mặc định của trình duyệt. Khóa bí mật của người dùng thường nằm trong USB token hoặc được giữ tại nhà cung cấp theo hình thức ký số từ xa. Các khái niệm của báo cáo áp dụng trực tiếp: người kiểm tra cần dựng được chuỗi tới Root CA quốc gia, kiểm tra thời hạn và trạng thái thu hồi của chứng thư người ký tại thời điểm ký, và dựa vào dấu thời gian khi kiểm tra lại về sau.

\section{Certificate Transparency}
Sự cố DigiNotar cho thấy chủ tên miền không có cách nào biết một CA đã cấp chứng thư cho tên của mình. Certificate Transparency (CT) giải quyết vấn đề này bằng các nhật ký công khai chỉ cho phép ghi thêm, tổ chức dưới dạng cây Merkle để ai cũng kiểm tra được log không bị sửa hay xóa~\cite{rfc6962,rfc9162}. CA gửi chứng thư (hoặc tiền chứng thư) vào log và nhận lại bằng chứng ghi log (SCT) để gắn kèm; các trình duyệt lớn yêu cầu chứng thư công cộng phải có SCT. Chủ tên miền và bên giám sát theo dõi log để phát hiện chứng thư lạ.

CT là cơ chế \emph{phát hiện}, không phải \emph{ngăn chặn}: một CA vẫn có thể ký chứng thư sai, nhưng hành vi đó để lại dấu vết công khai và CA phải chịu trách nhiệm. CT do đó bổ sung cho kiểm tra phía client ở Chương~\ref{chap:x509}, chứ không thay thế. Lab của đề tài là private PKI nên không dùng CT.

\section{Tự động hóa và chứng thư ngắn hạn}
Giao thức ACME (RFC 8555) cho phép máy chủ tự động chứng minh quyền kiểm soát tên miền, ví dụ đặt một tệp tại đường dẫn định trước (HTTP-01) hoặc một bản ghi DNS (DNS-01), rồi nhận và gia hạn chứng thư mà không cần thao tác thủ công~\cite{rfc8555}. ACME đóng vai trò RA tự động cho chứng thư xác minh tên miền và là nền tảng của Let's Encrypt.

Tự động hóa cho phép rút ngắn thời hạn chứng thư, và đây là cách giảm phụ thuộc vào thu hồi như đã phân tích ở Chương~\ref{chap:security}. Năm 2025, CA/Browser Forum thông qua Ballot SC-081v3 giảm dần thời hạn tối đa của chứng thư TLS công cộng từ 398 ngày xuống 200 ngày (từ 15/3/2026), 100 ngày (từ 15/3/2027) và 47 ngày (từ 15/3/2029), đồng thời rút ngắn thời gian được tái sử dụng kết quả xác minh tên miền~\cite{cabfsc081}. Với thời hạn ngắn như vậy, gia hạn thủ công không còn khả thi; các tổ chức cần quản lý vòng đời chứng thư tự động.

\section{Mật mã hậu lượng tử và PKI}
Thuật toán Shor trên một máy tính lượng tử đủ lớn sẽ phá được RSA và các lược đồ dựa trên logarit rời rạc, trong đó có ECDSA và ECDHE. Năm 2024, NIST công bố ba chuẩn đầu tiên: ML-KEM (FIPS 203) cho thiết lập khóa, ML-DSA (FIPS 204) và SLH-DSA (FIPS 205) cho chữ ký số~\cite{fips203,fips204,fips205}.

Hai phần của PKI và TLS chịu ảnh hưởng khác nhau, đúng như sự phân tách ở Chương~\ref{chap:crypto}. Phần \emph{thiết lập khóa phiên} cần chuyển sớm vì kẻ tấn công có thể ghi lại lưu lượng hôm nay để giải mã khi có máy tính lượng tử; các trình duyệt và máy chủ lớn đã triển khai trao đổi khóa lai kết hợp ECDHE với ML-KEM. Phần \emph{chữ ký trên chứng thư} chỉ cần an toàn tại thời điểm xác minh, nên ít gấp hơn, nhưng khó chuyển hơn: chữ ký và khóa công khai ML-DSA lớn hơn nhiều so với ECDSA, làm chuỗi chứng thư nặng hơn, và mọi CA, trust store, phần mềm kiểm tra đều phải cập nhật. Với chữ ký văn bản cần giá trị dài hạn, vấn đề còn quan trọng hơn vì chữ ký phải còn kiểm tra được sau nhiều năm.

\section{Hướng phát triển của đề tài}
Từ những giới hạn đã nêu, có thể mở rộng đề tài theo các hướng sau, sắp theo mức độ khả thi:
\begin{enumerate}
\item \textbf{OCSP trong lab:} chạy responder bằng \texttt{openssl ocsp} để so sánh với CRL về trạng thái \emph{good}, \emph{revoked}, \emph{unknown} và hành vi khi responder không phản hồi.
\item \textbf{Đo chi phí để kiểm chứng H3:} dùng \texttt{openssl speed} so sánh tốc độ ký và xác minh của RSA với ECDSA, kích thước CRL theo số chứng thư bị thu hồi.
\item \textbf{Xác thực hai chiều (mTLS):} cấp chứng thư \texttt{clientAuth} và yêu cầu client trình chứng thư khi kết nối.
\item \textbf{Ký văn bản có dấu thời gian:} dùng \texttt{openssl cms} và \texttt{openssl ts} để minh họa kiểm tra chữ ký sau khi chứng thư hết hạn.
\item \textbf{Chứng thư hậu lượng tử:} các bản OpenSSL từ 3.5 hỗ trợ ML-DSA, có thể dựng một chuỗi CA dùng ML-DSA và so sánh kích thước với RSA/ECDSA.
\end{enumerate}

\chapter*{Kết luận và hướng phát triển}
\label{chap:ket-luan}
\addcontentsline{toc}{chapter}{Kết luận và hướng phát triển}

\section*{1. Tổng kết và lời giải cho các câu hỏi nghiên cứu}

Báo cáo nghiên cứu đã giải quyết trọn vẹn bài toán nghiên cứu trọng tâm: \textit{Làm thế nào để phân phối và xác thực khóa công khai một cách đáng tin cậy, hạn chế tối đa nguy cơ giả mạo và tấn công thay thế khóa của kẻ đứng giữa (Eve), đồng thời đánh giá được các điều kiện chấp nhận và giới hạn an ninh của hệ thống trong thực tế?} Thông qua việc kết hợp chặt chẽ giữa nghiên cứu lý thuyết chuẩn mực quốc tế và thực nghiệm kỹ thuật trên môi trường lab OpenSSL, đề tài đưa ra lời giải đáp thỏa đáng cho 4 câu hỏi nghiên cứu (RQ1--RQ4):

\begin{itemize}
    \item \textbf{Trả lời RQ1 (Cơ chế liên kết danh tính và khóa công khai):} PKI giải quyết bài toán niềm tin phân tán trên mạng mở bằng cách chuyển dịch niềm tin sang các bên thứ ba có thẩm quyền (CA). Thông qua cấu trúc dữ liệu chuẩn hóa ASN.1 DER của khối \texttt{TBSCertificate} (RFC~5280~\cite{rfc5280}), CA đóng gói khóa công khai ($PK$) cùng thông tin định danh của chủ thể (trường SAN và DN) và ký xác thực bằng khóa bí mật của mình ($SK_{\mathrm{CA}}$). Chủ thể chứng minh quyền sở hữu khóa thông qua chữ ký tự thân trên yêu cầu ký chứng thư CSR (RFC~2986~\cite{rfc2986}), trong khi CA/RA thực thi quy trình thẩm tra hồ sơ độc lập theo các cấp độ xác thực (DV, OV, EV) quy định trong Baseline Requirements~\cite{cabforum_br}. Liên kết này bảo đảm trọn vẹn tính toàn vẹn và chống chối bỏ nhờ đặc tính mật mã của chữ ký số.

    \item \textbf{Trả lời RQ2 (Các điều kiện chấp nhận và từ chối chứng thư X.509):} Một chứng thư số chỉ được relying party chấp thuận khi và chỉ khi thỏa mãn đồng thời 6 điều kiện ngữ cảnh theo thuật toán thẩm định đường dẫn RFC~5280: (1) Xây dựng thành công chuỗi tin cậy liên tục dẫn về một Trust Anchor nằm trong Trust Store cục bộ; (2) Tính toàn vẹn mật mã của chữ ký số trên từng mắt xích; (3) Tuân thủ ràng buộc cơ bản Basic Constraints (\texttt{cA:TRUE} cho CA và \texttt{cA:FALSE} cho thực thể cuối); (4) Khớp mục đích sử dụng trong Key Usage và Extended Key Usage (như cờ \texttt{serverAuth} cho máy chủ TLS); (5) Định danh máy chủ truy cập trùng khớp với danh sách Subject Alternative Name (SAN theo RFC~9525~\cite{rfc9525}); và (6) Thời điểm xác minh nằm trong khung hiệu lực \texttt{[notBefore, notAfter]} đồng thời chứng thư chưa bị hủy bỏ trên danh sách thu hồi CRL hoặc OCSP. Việc vi phạm bất kỳ điều kiện nào đều dẫn tới việc từ chối ngay lập tức (được chứng minh qua các mã lỗi tương ứng tại TC02--TC08, TC10).

    \item \textbf{Trả lời RQ3 (Vai trò và sự đánh đổi của cơ chế thu hồi CRL/OCSP):} Cơ chế thu hồi là nhánh rẽ khẩn cấp thiết yếu để vô hiệu hóa chứng thư khi khóa bí mật bị lộ hoặc quy trình cấp phát có sai sót trước khi chứng thư hết hạn tự nhiên. Tuy nhiên, cơ chế này bộc lộ sự đánh đổi kỹ thuật rõ rệt: CRL (RFC~5280) gặp phải độ trễ truyền bá theo chu kỳ phát hành (\texttt{nextUpdate}) và kích thước tệp danh sách tăng tuyến tính theo số lượng mục thu hồi, gây lãng phí băng thông và tăng thời gian thẩm định của client; trong khi OCSP (RFC~6960~\cite{rfc6960}) đòi hỏi máy chủ responder trực tuyến liên tục, đối mặt với rủi ro lộ quyền riêng tư truy vấn và tính khả dụng của mạng. Đáng chú ý, trạng thái \texttt{good} của OCSP chỉ mang ý nghĩa số sê-ri chưa bị thu hồi, hoàn toàn không thay thế cho quy trình thẩm định chuỗi và ngữ cảnh toàn diện.

    \item \textbf{Trả lời RQ4 (Các điểm yếu và giới hạn an ninh cố hữu của PKI):} Nghiên cứu chỉ ra rằng an ninh của PKI phụ thuộc cốt lõi vào sự toàn vẹn của mô hình tin cậy và chính sách vận hành của các CA gốc. Minh chứng từ thảm họa DigiNotar (Fox-IT, 2012~\cite{prins2011diginotar}) cho thấy khi một CA bị thỏa hiệp, kẻ tấn công có thể phát hành các chứng thư giả mạo có chữ ký mật mã hoàn toàn hợp lệ, khiến mọi cơ chế xác minh cục bộ của client đều hoàn toàn bất lực. Bên cạnh đó, chữ ký số hợp lệ không đồng nghĩa với phần mềm an toàn (rủi ro mã độc ký số do lộ khóa riêng) hay chủ thể lương thiện (trang web lừa đảo sử dụng HTTPS hợp lệ). Cuối cùng, sự phát triển của máy tính lượng tử với thuật toán Shor (1994)~\cite{shor1994} đặt ra mối đe dọa sinh tử, có khả năng phá vỡ toàn bộ nền tảng toán học của RSA và ECDSA đang bảo vệ hệ thống PKI toàn cầu.
\end{itemize}

\section*{2. Đánh giá và kiểm chứng các giả thuyết nghiên cứu}

Dựa trên dữ liệu thực nghiệm tại Chương~\ref{chap:experiment}, đề tài đưa ra kết luận đánh giá cụ thể đối với ba giả thuyết nghiên cứu đã đặt ra:

\begin{itemize}
    \item \textbf{Giả thuyết $H_1$ (Điều kiện cần của chữ ký số và vai trò chính sách client):} \textit{Được ủng hộ hoàn toàn}. Dữ liệu thực nghiệm từ Bảng~\ref{tab:exp-results} chỉ ra rằng: ngoại trừ TC05 bị lỗi chữ ký và TC02 là chứng thư tự ký, toàn bộ 6 ca kiểm thử bị từ chối còn lại (TC03, TC04, TC06, TC07, TC08, TC10) đều sở hữu chữ ký số hoàn toàn nguyên vẹn về mặt toán học do CA ký phát. Tuy nhiên, chúng đều bị OpenSSL trả về mã lỗi và từ chối chấp nhận do vi phạm các ràng buộc chính sách ngữ cảnh. Đặc biệt, tại mốc chuẩn TC01, chứng thư thực chất đã bị thu hồi trên CA nhưng vẫn được OpenSSL chấp nhận (\texttt{OK}) khi client không bật kiểm tra CRL; trong khi ở TC09, chứng thư thỏa mãn đồng thời cả chuỗi tin cậy, ngữ cảnh và trạng thái chưa thu hồi trên CRL thì đạt \texttt{OK}. Kết quả khẳng định chắc chắn: chữ ký số hợp lệ chỉ là điều kiện cần chứ chưa đủ để một chứng thư được client tin cậy.

    \item \textbf{Giả thuyết $H_2$ (Tính phụ thuộc ngữ cảnh và sự đảo ngược quyết định):} \textit{Được ủng hộ hoàn toàn}. Quyết định thẩm định của client mang tính động và phụ thuộc trực tiếp vào các tham số ngữ cảnh. Khi giữ nguyên chứng thư số và chữ ký không đổi, việc điều chỉnh thời điểm kiểm tra (\texttt{-attime} ở TC03), tên miền truy vấn (TC04), kích hoạt kiểm tra thu hồi (\texttt{-crl\_check} ở TC07), thay đổi Trust Anchor chỉ định (đối chứng \texttt{other-ca/root.crt}) hoặc thay đổi mục đích sử dụng kiểm tra (\texttt{-purpose sslclient} trên \texttt{client-only.crt}) đều trực tiếp làm đảo ngược kết quả thẩm định từ chấp thuận (\texttt{OK}) sang từ chối (\texttt{FAIL}) và ngược lại.

    \item \textbf{Giả thuyết $H_3$ (Đánh đổi về chi phí vận hành của cơ chế thu hồi):} \textit{Được kiểm chứng một phần}. Thông qua phân tích chuẩn RFC~5280 kết hợp phép đo lường định lượng cục bộ với OpenSSL 3.5.9, kết quả thực nghiệm chứng minh rằng: khi số lượng mục thu hồi tăng từ $N=0$ lên $N=10.000$, kích thước tệp danh sách CRL tăng tuyến tính từ 564 bytes lên 474.532 bytes ($\approx 47$~bytes/mục thu hồi, tăng gấp 841 lần), đồng thời thời gian thẩm định trung vị của client tăng từ 37{,}46~ms lên 86{,}46~ms (tăng gấp hơn 2{,}3 lần). Dữ liệu ủng hộ nhận định về sự đánh đổi tài nguyên và giới hạn mở rộng của cơ chế CRL truyền thống.
\end{itemize}

\section*{3. Các đóng góp chính của đề tài}

Đề tài đã hoàn thành các mục tiêu nghiên cứu đề ra với ba đóng góp thiết thực:
\begin{enumerate}
    \item \textbf{Hệ thống hóa toàn diện cơ sở lý thuyết và tiêu chuẩn PKI:} Cung cấp bức tranh toàn cảnh, mạch lạc về cơ sở toán học mật mã, kiến trúc phân cấp, vòng đời chứng thư X.509 v3 và quy chuẩn quản trị CP/CPS theo các văn bản chuẩn mực quốc tế của IETF, ITU-T, NIST và CA/Browser Forum.
    \item \textbf{Xây dựng mô hình thực nghiệm và bộ kiểm thử tự động hóa:} Thiết lập thành công môi trường lab PKI phân cấp hai tầng độc lập bằng OpenSSL 3.5.9, tự động hóa quy trình qua Makefile và Bash script, kiểm chứng tường minh 10 ca kiểm thử điển hình đại diện cho các điều kiện biên của RFC~5280 và thực hiện phép đo đạc định lượng đối với cơ chế CRL.
    \item \textbf{Đề xuất mô hình phòng thủ chiều sâu và phân tích xu hướng công nghệ tương lai:} Tổng hợp mô hình an ninh 6 lớp cho PKI hiện đại, phân tích sâu sắc nguyên lý và giới hạn của cơ chế Certificate Transparency (RFC~9162), giải pháp tự động hóa ACME (RFC~8555) và vạch ra lộ trình di trú sang các chuẩn mật mã hậu lượng tử của NIST (FIPS 203, 204, 205).
\end{enumerate}

\section*{4. Hạn chế của nghiên cứu}

Bên cạnh các kết quả đạt được, nghiên cứu tồn tại một số hạn chế kỹ thuật cần được ghi nhận:
\begin{itemize}
    \item \textbf{Phạm vi kiểm thử cục bộ ở tầng OpenSSL CLI:} Thực nghiệm sử dụng lệnh \texttt{openssl verify} để kiểm tra thuật toán thẩm định đường dẫn tĩnh trên các tệp tin cục bộ, chưa thiết lập tiến trình bắt tay mạng TLS 1.3 hoàn chỉnh để kiểm tra sự phối hợp giữa chữ ký CA trên chứng thư và chữ ký bắt tay thời gian thực (\texttt{CertificateVerify}).
    \item \textbf{Giới hạn cơ chế thu hồi ngoại tuyến:} Nghiên cứu mới dừng lại ở việc kiểm thử danh sách thu hồi CRL tĩnh, chưa dựng máy chủ dịch vụ trạng thái chứng thư trực tuyến (OCSP Responder) hay kiểm nghiệm cơ chế đính kèm phản hồi OCSP (OCSP Stapling) trong giao tiếp mạng thực tế.
    \item \textbf{Chưa khảo sát đa dạng môi trường và thư viện:} Thực nghiệm chỉ triển khai trên thư viện OpenSSL 3.5.9 trong môi trường Linux Container, chưa mở rộng đối sánh hành vi xác minh trên các thư viện mật mã khác (như GnuTLS, NSS, BoringSSL) cũng như các cơ chế tối ưu riêng của trình duyệt web (AIA fetching, OneCRL, CRLSet).
    \item \textbf{Giới hạn về triển khai các công nghệ tiên tiến:} Các nội dung chuyên sâu về Certificate Transparency (CT Log, Monitor), giao thức tự động hóa ACME và các thuật toán mật mã hậu lượng tử (PQC) mới dừng lại ở mức độ phân tích kiến trúc lý thuyết và đối chiếu tiêu chuẩn, chưa được tích hợp và kiểm thử thực nghiệm trong môi trường lab của đề tài.
\end{itemize}

\section*{5. Hướng phát triển tiếp theo}

Từ các kết quả và hạn chế nêu trên, nhóm nghiên cứu đề xuất bốn hướng phát triển khả thi cho các công trình tiếp theo:
\begin{enumerate}
    \item \textbf{Mở rộng thực nghiệm sang bắt tay mạng TLS 1.3 và mTLS:} Cài đặt chuỗi chứng thư lên các máy chủ web dịch vụ thực tế (như Nginx, Caddy hoặc Envoy Proxy), triển khai xác thực hai chiều (Mutual TLS -- mTLS) và sử dụng công cụ phân tích gói tin (Wireshark/tcpdump) để khảo sát trực tiếp các thông điệp \texttt{Certificate}, \texttt{CertificateVerify} và \texttt{Finished}.
    \item \textbf{Xây dựng hạ tầng OCSP Responder và thử nghiệm OCSP Stapling:} Thiết lập dịch vụ OCSP trực tuyến độc lập bằng OpenSSL OCSP server hoặc Boulder, cấu hình máy chủ web hỗ trợ OCSP Stapling, đo lường độ trễ bắt tay mạng và đánh giá tính sẵn sàng của dịch vụ thu hồi.
    \item \textbf{Tự động hóa vòng đời chứng thư với giao thức ACME:} Triển khai một máy chủ ACME nội bộ (như Smallstep CA hoặc Boulder trong mạng thử nghiệm), cấu hình ACME Client (Certbot/Caddy) tự động thực hiện các thử thách HTTP-01 và DNS-01 để quản trị tự động việc cấp phát, cài đặt và gia hạn chứng thư.
    \item \textbf{Nghiên cứu thực nghiệm mật mã hậu lượng tử (PQC Migration):} Tích hợp thư viện Open Quantum Safe (OQS Provider) vào OpenSSL 3.x, sinh cặp khóa và cấp phát chứng thư số X.509 hậu lượng tử sử dụng các thuật toán chuẩn hóa mới của NIST gồm **ML-DSA** (FIPS 204) và **SLH-DSA** (FIPS 205). Tiến hành đo lường định lượng kích thước gói tin, độ trễ bắt tay TLS và đánh giá hiện tượng phân mảnh gói tin TCP/IP trong môi trường mạng thực tế nhằm chuẩn bị sẵn sàng cho kỷ nguyên hậu lượng tử.
\end{enumerate}


\bibliographystyle{unsrt}
\bibliography{report/references}
\end{document}
